% Generated mechanically from frozen input inputs/p03/ja/tex/sdga.tex.
% Do not edit this generated file; edit the bound locale source instead.

% OK, start here.
%

\setcounter{chapter}{23}
\chapter{微分次数付き層}



\phantomsection
\label{sdga-section-phantom}


\section{序論}
\label{sdga-section-introduction}

\noindent
本章は Differential Graded Algebra, Section
\ref{dga-section-introduction} で始めた議論の続きである。
概説論文として \cite{Keller-survey} がある。




\section{規約}
\label{sdga-section-conventions}

\noindent
本章でも、{\it 環}とは単位元 $1$ をもつ可換環であるという規約を維持する。
$R$ を環とするとき、{\it $R$-代数 $A$} とは、$R$-加群 $A$ と、
結合的かつ単位元をもつ乗法である $R$-双線形写像
$A \times A \to A$ を備えたものをいう。言い換えると、これは構造写像
$R \to A$ の像が $A$ の中心に含まれる単位的結合 $R$-代数である。








\section{次数付き代数の層}
\label{sdga-section-ga}

\noindent
この節は読み飛ばしてよい。

\begin{definition}
\label{sdga-definition-ga}
$(\mathcal{C}, \mathcal{O})$ を環付きサイトとする。
$(\mathcal{C},\mathcal O)$ 上の{\it 次数付き $\mathcal{O}$-代数の層}
または{\it 次数付き代数の層}とは、$n \in \mathbf{Z}$ で添字付けられた
$\mathcal{O}$-加群の族 $\mathcal{A}^n$ と、$\mathcal{O}$-双線形写像
$$
\mathcal{A}^n \times \mathcal{A}^m \to \mathcal{A}^{n + m},\quad
(a, b) \longmapsto ab
$$
で与えられる、次の性質をもつものをいう。これらの写像を乗法写像と呼ぶ。
\begin{enumerate}
\item 乗法は結合的である。
\item $\mathcal{A}^0$ には、乗法の両側単位元となる大域切断 $1$ が存在する。
\end{enumerate}
このような構造をしばしば $\mathcal{A}$ と書く。
{\it 次数付き $\mathcal{O}$-代数の準同型}
$f : \mathcal{A} \to \mathcal{B}$ とは、乗法写像と両立する
$\mathcal{O}$-加群の写像の族
$f^n : \mathcal{A}^n \to \mathcal{B}^n$ をいう。
\end{definition}

\noindent
次数付き $\mathcal{O}$-代数 $\mathcal{A}$ と対象
$U \in \Ob(\mathcal{C})$ が与えられたとき、記号
$$
\mathcal{A}(U) =
\Gamma(U, \mathcal{A}) =
\bigoplus\nolimits_{n \in \mathbf{Z}} \mathcal{A}^n(U)
$$
を用いる。これは次数付き $\mathcal{O}(U)$-代数である。

\begin{remark}
\label{sdga-remark-functoriality-ga}
$(f, f^\sharp) : (\Sh(\mathcal{C}), \mathcal{O}_\mathcal{C})
\to (\Sh(\mathcal{D}), \mathcal{O}_\mathcal{D})$
を環付きトポスの射とする。次が成り立つ。
\begin{enumerate}
\item $\mathcal{A}$ を次数付き $\mathcal{O}_\mathcal{C}$-代数とする。
$\mathcal{A}$ の乗法写像は乗法写像
$f_*\mathcal{A}^n \times f_*\mathcal{A}^m \to f_*\mathcal{A}^{n + m}$
を誘導し、$f^\sharp$ によってこれらを
$\mathcal{O}_\mathcal{D}$-双線形写像とみなせる。このようにして得られる次数付き
$\mathcal{O}_\mathcal{D}$-代数を $f_*\mathcal{A}$ と書く。
\item $\mathcal{B}$ を次数付き $\mathcal{O}_\mathcal{D}$-代数とする。
$\mathcal{B}$ の乗法写像は乗法写像
$f^*\mathcal{B}^n \times f^*\mathcal{B}^m \to f^*\mathcal{B}^{n + m}$
を誘導し、$f^\sharp$ を用いてこれらを
$\mathcal{O}_\mathcal{C}$-双線形写像とみなせる。このようにして得られる次数付き
$\mathcal{O}_\mathcal{C}$-代数を $f^*\mathcal{B}$ と書く。
\item 次数付き $\mathcal{O}_\mathcal{C}$-代数の準同型
$f^*\mathcal{B} \to \mathcal{A}$ の集合と、次数付き
$\mathcal{O}_\mathcal{C}$-代数の準同型
$\mathcal{B} \to f_*\mathcal{A}$ の集合との間には一対一対応がある。
\end{enumerate}
(3) は加群の層上の $f^*$ と $f_*$ の通常の随伴から直ちに従う。
\end{remark}




\section{次数付き加群の層}
\label{sdga-section-graded-modules}

\noindent
本節は読み飛ばしてよい。

\begin{definition}
\label{sdga-definition-gm}
$(\mathcal{C}, \mathcal{O})$ を環付きサイトとし、$\mathcal{A}$ を
$(\mathcal{C}, \mathcal{O})$ 上の次数付き代数の層とする。
$\mathcal{A}$ 上の（右）{\it 次数付き $\mathcal{A}$-加群}、または
（右）{\it 次数付き加群}とは、$n \in \mathbf{Z}$ で添字付けられた
$\mathcal{O}$-加群の族 $\mathcal{M}^n$ と、$\mathcal{O}$-双線形写像
$$
\mathcal{M}^n \times \mathcal{A}^m \to \mathcal{M}^{n + m},\quad
(x, a) \longmapsto xa
$$
であって、次の性質を満たすものとして与えられる。これらの写像を
乗法写像という。
\begin{enumerate}
\item 乗法は $(xa)a' = x(aa')$ を満たす。
\item $\mathcal{A}^0$ の単位切断 $1$ は、すべての $n$ に対して
$\mathcal{M}^n$ 上に恒等写像として作用する。
\end{enumerate}
この状況を表すために、単に「$\mathcal{M}$ を次数付き
$\mathcal{A}$-加群とする」としばしばいう。
次数付き $\mathcal{A}$-加群の{\it 準同型}
$f : \mathcal{M} \to \mathcal{N}$ とは、乗法写像と両立する
$\mathcal{O}$-加群の写像の族
$f^n : \mathcal{M}^n \to \mathcal{N}^n$ のことである。
（右）次数付き $\mathcal{A}$-加群の圏を
$\textit{Mod}(\mathcal{A})$ と書く。
\end{definition}

\noindent
{\it 左次数付き加群}もまったく同様に定義できるが、本章では
特に断らない限り右加群を用いる。

\medskip\noindent
次数付き $\mathcal{A}$-加群 $\mathcal{M}$ と対象
$U \in \Ob(\mathcal{C})$ に対して、記法
$$
\mathcal{M}(U) =
\Gamma(U, \mathcal{M}) =
\bigoplus\nolimits_{n \in \mathbf{Z}} \mathcal{M}^n(U)
$$
を用いる。これは（右）次数付き $\mathcal{A}(U)$-加群である。

\begin{lemma}
\label{sdga-lemma-gm-abelian}
$(\mathcal{C}, \mathcal{O})$ を環付きサイトとし、$\mathcal{A}$ を
次数付き $\mathcal{O}$-代数とする。圏 $\textit{Mod}(\mathcal{A})$
は次の性質をもつアーベル圏である。
\begin{enumerate}
\item $\textit{Mod}(\mathcal{A})$ は任意の直和をもつ。
\item $\textit{Mod}(\mathcal{A})$ は任意の余極限をもつ。
\item $\textit{Mod}(\mathcal{A})$ におけるフィルター余極限は完全である。
\item $\textit{Mod}(\mathcal{A})$ は任意の直積をもつ。
\item $\textit{Mod}(\mathcal{A})$ は任意の極限をもつ。
\end{enumerate}
次数付き $\mathcal{A}$-加群をその第 $n$ 項へ送る函子
$$
\textit{Mod}(\mathcal{A}) \longrightarrow \textit{Mod}(\mathcal{O}),\quad
\mathcal{M} \longmapsto \mathcal{M}^n
$$
はすべての極限および余極限と交換する。
\end{lemma}

\noindent
この補題は、極限と余極限を項ごとに取れることを述べている。
また、忘却関手
$$
\textit{Mod}(\mathcal{A})
\longrightarrow
\text{次数付き }\mathcal{O}\text{-加群}
$$
がすべての極限および余極限と交換することも述べている
（あるいは、好みによってはそれを含意している）。

\begin{proof}
補題の主張に現れる函子を
$\text{gr}^n : \textit{Mod}(\mathcal{A}) \to \textit{Mod}(\mathcal{O})$
と書く。次数付き $\mathcal{A}$-加群の準同型
$f : \mathcal{M} \to \mathcal{N}$ を考える。$f$ を次数付き
$\mathcal{O}$-加群の写像とみたときの核と余核には、定義
\ref{sdga-definition-gm} の乗法写像も備わる。したがって、これらは
$\textit{Mod}(\mathcal{A})$ における核と余核でもある。
ゆえに $\textit{Mod}(\mathcal{A})$ はアーベル圏であり、核と余核を
取る操作は $\text{gr}^n$ と交換する。

\medskip\noindent
極限と余極限の存在を証明するには、直積と直和の存在を示せば十分である。
Categories, Lemmas \ref{categories-lemma-limits-products-equalizers} and
\ref{categories-lemma-colimits-coproducts-coequalizers} を参照せよ。
同じ補題により、$\text{gr}^n$ が直和および直積と交換することを
示せば、極限および余極限とも交換することが従う。

\medskip\noindent
$\mathcal{M}_t$（$t \in T$）を次数付き $\mathcal{A}$-加群の族とする。
次数 $n$ の項が $\bigoplus_{t \in T} \mathcal{M}_t^n$ である
次数付き $\mathcal{A}$-加群を考えられる（乗法写像は明らかである）。
これが $\textit{Mod}(\mathcal{A})$ における直和であることは
容易に確かめられる。直積についても同様である。

\medskip\noindent
各 $n$ に対して $\text{gr}^n$ は完全関手であり、
$\textit{Mod}(\mathcal{A})$ の複体
$\mathcal{M}_1 \to \mathcal{M}_2 \to \mathcal{M}_3$ が完全であることと、
すべての $n$ に対して
$\text{gr}^n\mathcal{M}_1 \to \text{gr}^n\mathcal{M}_2 \to
\text{gr}^n\mathcal{M}_3$ が $\textit{Mod}(\mathcal{O})$ で
完全であることとは同値である。したがって (3) が従う。実際、
$\textit{Mod}(\mathcal{O})$ はグロタンディーク・アーベル圏なので、
そこでのフィルター余極限は完全である。Cohomology on Sites,
Section \ref{sites-cohomology-section-unbounded} を参照せよ。
\end{proof}






\section{次数付き加群層の次数付き圏}
\label{sdga-section-gm-gr-cat}

\noindent
本節は読み飛ばしてよい。本節は Differential Graded Algebra,
Example \ref{dga-example-gm-gr-cat} の類似である。次数付き圏に関する
規約については Differential Graded Algebra,
Section \ref{dga-section-graded} を参照せよ。

\medskip\noindent
$(\mathcal{C}, \mathcal{O})$ を環付きサイトとし、$\mathcal{A}$ を
$(\mathcal{C}, \mathcal{O})$ 上の次数付き代数の層とする。
$R = \Gamma(\mathcal{C}, \mathcal{O})$ 上の次数付き圏
$\textit{Mod}^{gr}(\mathcal{A})$ であって、付随する圏
$(\textit{Mod}^{gr}(\mathcal{A}))^0$ が次数付き
$\mathcal{A}$-加群の圏となるものを構成する。
$\textit{Mod}^{gr}(\mathcal{A})$ の対象としては、右次数付き
$\mathcal{A}$-加群を取る（Section \ref{sdga-section-graded-modules} を参照）。
次数付き $\mathcal{A}$-加群 $\mathcal{L}$ と $\mathcal{M}$ に対して
$$
\Hom_{\textit{Mod}^{gr}(\mathcal{A})}(\mathcal{L}, \mathcal{M}) =
\bigoplus\nolimits_{n \in \mathbf{Z}} \Hom^n(\mathcal{L}, \mathcal{M})
$$
とおく。ここで $\Hom^n(\mathcal{L}, \mathcal{M})$ は、次数 $n$ の
斉次な右 $\mathcal{A}$-加群写像
$f : \mathcal{L} \to \mathcal{M}$ の集合である。より正確には、$f$ は
$i \in \mathbf{Z}$ に対する写像の族
$f : \mathcal{L}^i \to \mathcal{M}^{i + n}$ であって、乗法写像と
両立するものとして与えられる。成分で書けば、
$$
\Hom^n(\mathcal{L}, \mathcal{M})
\subset
\prod\nolimits_{p + q = n}
\Hom_\mathcal{O}(\mathcal{L}^{-q}, \mathcal{M}^p)
$$
（添字の順序が逆転していることに注意せよ）は、
$$
f_{p, q}(m a) = f_{p - i, q + i}(m)a
$$
を満たす $f = (f_{p, q})$ からなる部分集合である。ここで $a$ は
$\mathcal{A}^i$ の局所切断、$m$ は $\mathcal{L}^{-q-i}$ の局所切断である。
次数付き $\mathcal{A}$-加群 $\mathcal{K}$、$\mathcal{L}$、
$\mathcal{M}$ に対して、写像
$$
\Hom^m(\mathcal{L}, \mathcal{M}) \times
\Hom^n(\mathcal{K}, \mathcal{L}) \longrightarrow
\Hom^{n + m}(\mathcal{K}, \mathcal{M})
$$
により $\textit{Mod}^{gr}(\mathcal{A})$ における合成を定義する。
これは右 $\mathcal{A}$-加群写像の通常の合成、すなわち
$(g,f)\mapsto g\circ f$ である。





\section{次数付き加群層のテンソル積}
\label{sdga-section-tensor-product}

\noindent
本節は読み飛ばしてよい。本節は Differential Graded Algebra,
Section \ref{dga-section-tensor-product} の一部の類似である。

\medskip\noindent
$(\mathcal{C}, \mathcal{O})$ を環付きサイトとし、$\mathcal{A}$ を
$(\mathcal{C},\mathcal{O})$ 上の次数付き代数の層とする。
$\mathcal{M}$ を右次数付き $\mathcal{A}$-加群、$\mathcal{N}$ を左次数付き
$\mathcal{A}$-加群とする。このとき{\it テンソル積}
$\mathcal{M} \otimes_\mathcal{A} \mathcal{N}$
を、次数 $n$ の項が
$$
(\mathcal{M} \otimes_\mathcal{A} \mathcal{N})^n =
\Coker\left(
\bigoplus\nolimits_{r + s + t = n} \mathcal{M}^r \otimes_\mathcal{O}
\mathcal{A}^s \otimes_\mathcal{O} \mathcal{N}^t
\longrightarrow
\bigoplus\nolimits_{p + q = n} \mathcal{M}^p \otimes_\mathcal{O} \mathcal{N}^q
\right)
$$
である次数付き $\mathcal{O}$-加群として定義する。ここで、この写像は
$\mathcal{M}^r\otimes_\mathcal{O}\mathcal{A}^s
\otimes_\mathcal{O}\mathcal{N}^t$ の局所切断
$x\otimes a\otimes y$ を $xa\otimes y-x\otimes ay$ へ送る。
この定義により、
$(\mathcal{M} \otimes_\mathcal{A} \mathcal{N})^n$
は前層
$U \mapsto (\mathcal{M}(U) \otimes_{\mathcal{A}(U)} \mathcal{N}(U))^n$
の層化である。ここで、次数付き加群のテンソル積は
Differential Graded Algebra,
Section \ref{dga-section-tensor-product} で定義したものである。

\medskip\noindent
左次数付き $\mathcal{A}$-加群 $\mathcal{N}$ を固定すると、関手
$$
- \otimes_\mathcal{A} \mathcal{N} :
\textit{Mod}(\mathcal{A})
\longrightarrow
\text{Gr}(\textit{Mod}(\mathcal{O})) =
\text{次数付き }\mathcal{O}\text{-加群}
$$
を得る。記法 $\text{Gr}(-)$ については Homology,
Definition \ref{homology-definition-graded} を参照せよ。
次数付き $\mathcal{O}$-加群の次数付き圏を
$\text{Gr}^{gr}(\textit{Mod}(\mathcal{O}))$ と書く。Differential Graded
Algebra, Example \ref{dga-example-graded-category-graded-objects} を参照せよ。
上の関手は、次数付き圏の関手
$$
- \otimes_\mathcal{A} \mathcal{N} :
\textit{Mod}^{gr}(\mathcal{A})
\longrightarrow
\text{Gr}^{gr}(\textit{Mod}(\mathcal{O}))
$$
へ拡張できる。これは $\mathcal{M} \to \mathcal{M}'$ の次数 $n$ の
準同型を、$\mathcal{M}\otimes_\mathcal{A}\mathcal{N}$ から
$\mathcal{M}'\otimes_\mathcal{A}\mathcal{N}$ への誘導される次数 $n$ の
写像へ送る。






\section{次数付き加群層の内部 Hom}
\label{sdga-section-internal-hom-graded}

\noindent
本節は読み飛ばすことを強く勧める。

\medskip\noindent
Section \ref{sdga-section-gm-gr-cat} の構成を層化したものが必要となる。
$(\mathcal{C}, \mathcal{O})$、$\mathcal{A}$、$\mathcal{M}$、
$\mathcal{L}$ を Section \ref{sdga-section-gm-gr-cat} のとおりとする。このとき
$$
\SheafHom^{gr}_\mathcal{A}(\mathcal{M}, \mathcal{L})
$$
を、次数 $n$ の項
$$
\SheafHom^n_\mathcal{A}(\mathcal{M}, \mathcal{L})
\subset
\prod\nolimits_{p + q = n}
\SheafHom_\mathcal{O}(\mathcal{L}^{-q}, \mathcal{M}^p)
$$
が、
$$
f_{p, q}(m a) = f_{p - i, q + i}(m)a
$$
を満たす局所切断 $f=(f_{p,q})$ からなる部分層であるような、
次数付き $\mathcal{O}$-加群として定義する。ここで $a$ は
$\mathcal{A}^i$ の局所切断、$m$ は $\mathcal{L}^{-q-i}$ の
局所切断である。Section \ref{sdga-section-gm-gr-cat} と同様に、合成写像
$$
\SheafHom^{gr}_\mathcal{A}(\mathcal{L}, \mathcal{M}) \otimes_\mathcal{O}
\SheafHom^{gr}_\mathcal{A}(\mathcal{K}, \mathcal{L})
\longrightarrow
\SheafHom^{gr}_\mathcal{A}(\mathcal{K}, \mathcal{M})
$$
がある。左辺は Section \ref{sdga-section-tensor-product} で定義した
次数付き $\mathcal{O}$-加群のテンソル積である。この写像は合成写像
$$
\SheafHom^m_\mathcal{A}(\mathcal{L}, \mathcal{M}) \otimes_\mathcal{O}
\SheafHom^n_\mathcal{A}(\mathcal{K}, \mathcal{L}) \longrightarrow
\SheafHom^{n + m}_\mathcal{A}(\mathcal{K}, \mathcal{M})
$$
で与えられ、これは（局所的には）通常の合成として定義される。

\medskip\noindent
これらの定義により、合成と両立する次数付き $R$-加群として
$$
\Hom_{\textit{Mod}^{gr}(\mathcal{A})}(\mathcal{L}, \mathcal{M}) =
\Gamma(\mathcal{C}, \SheafHom^{gr}_\mathcal{A}(\mathcal{L}, \mathcal{M}))
$$
が成り立つ。







\section{次数付き双加群の層とテンソル--Hom 随伴}
\label{sdga-section-graded-bimodules}

\noindent
本節は読み飛ばしてよい。

\begin{definition}
\label{sdga-definition-bimodule}
$(\mathcal{C},\mathcal{O})$ を環付きサイトとし、$\mathcal{A}$ と
$\mathcal{B}$ を $(\mathcal{C},\mathcal{O})$ 上の次数付き代数の層とする。
{\it 次数付き $(\mathcal{A},\mathcal{B})$-双加群}とは、
$n\in\mathbf{Z}$ で添字付けられた $\mathcal{O}$-加群の族
$\mathcal{M}^n$ と、$\mathcal{O}$-双線形写像
$$
\mathcal{M}^n \times \mathcal{B}^m \to \mathcal{M}^{n + m},\quad
(x, b) \longmapsto xb
$$
および
$$
\mathcal{A}^n \times \mathcal{M}^m \to \mathcal{M}^{n + m},\quad
(a, x) \longmapsto ax
$$
であって、次の性質を満たすものとして与えられる。これらの写像を
乗法写像という。
\begin{enumerate}
\item 乗法は $a(a'x)=(aa')x$ および $(xb)b'=x(bb')$ を満たす。
\item $(ax)b = a(xb)$ である。
\item $\mathcal{A}^0$ の単位切断 $1$ は乗法により恒等写像として作用する。
\item $\mathcal{B}^0$ の単位切断 $1$ は乗法により恒等写像として作用する。
\end{enumerate}
このような構造をしばしば $\mathcal{M}$ と書く。
次数付き $(\mathcal{A},\mathcal{B})$-双加群の{\it 準同型}
$f:\mathcal{M}\to\mathcal{N}$ とは、乗法写像と両立する
$\mathcal{O}$-加群の写像の族
$f^n:\mathcal{M}^n\to\mathcal{N}^n$ のことである。
\end{definition}

\noindent
次数付き $(\mathcal{A},\mathcal{B})$-双加群 $\mathcal{M}$ と対象
$U\in\Ob(\mathcal{C})$ に対して、記法
$$
\mathcal{M}(U) =
\Gamma(U, \mathcal{M}) =
\bigoplus\nolimits_{n \in \mathbf{Z}} \mathcal{M}^n(U)
$$
を用いる。これは次数付き
$(\mathcal{A}(U),\mathcal{B}(U))$-双加群である。

\medskip\noindent
$(\mathcal{C},\mathcal{O})$ を環付きサイトとし、$\mathcal{A}$ と
$\mathcal{B}$ を $(\mathcal{C},\mathcal{O})$ 上の次数付き代数の層とする。
$\mathcal{M}$ を右次数付き $\mathcal{A}$-加群、$\mathcal{N}$ を次数付き
$(\mathcal{A},\mathcal{B})$-双加群とする。この場合、Section
\ref{sdga-section-tensor-product} で定義した次数付きテンソル積
$$
\mathcal{M} \otimes_\mathcal{A} \mathcal{N}
$$
は、明らかな乗法写像をもつ右次数付き $\mathcal{B}$-加群である。
この構成は関手および次数付き圏の関手
$$
\otimes_\mathcal{A} \mathcal{N} :
\textit{Mod}(\mathcal{A})
\longrightarrow
\textit{Mod}(\mathcal{B})
\quad\text{and}\quad
\otimes_\mathcal{A} \mathcal{N} :
\textit{Mod}^{gr}(\mathcal{A})
\longrightarrow
\textit{Mod}^{gr}(\mathcal{B})
$$
を定義する。これは $\mathcal{M}$ から $\mathcal{M}'$ への次数 $n$ の
準同型を、$\mathcal{M}\otimes_\mathcal{A}\mathcal{N}$ から
$\mathcal{M}'\otimes_\mathcal{A}\mathcal{N}$ への誘導される次数 $n$ の
写像へ送る。

\medskip\noindent
$(\mathcal{C},\mathcal{O})$ を環付きサイトとし、$\mathcal{A}$ と
$\mathcal{B}$ を $(\mathcal{C},\mathcal{O})$ 上の次数付き代数の層とする。
$\mathcal{N}$ を次数付き $(\mathcal{A},\mathcal{B})$-双加群、
$\mathcal{L}$ を右次数付き $\mathcal{B}$-加群とする。この場合、Section
\ref{sdga-section-internal-hom-graded} で定義した次数付き内部 Hom
$$
\SheafHom_\mathcal{B}^{gr}(\mathcal{N}, \mathcal{L})
$$
は右次数付き $\mathcal{A}$-加群であり、その乗法写像は
\footnote{ここでは符号を一切用いないという規約を採用している。}
$$
\SheafHom^n_\mathcal{B}(\mathcal{N}, \mathcal{L})
\times \mathcal{A}^m
\longrightarrow
\SheafHom^{n + m}_\mathcal{B}(\mathcal{N}, \mathcal{L})
$$
である。これは $U$ 上の
$\SheafHom^n_\mathcal{B}(\mathcal{N},\mathcal{L})$ の切断
$f=(f_{p,q})$ と $U$ 上の $\mathcal{A}^m$ の切断 $a$ を、$U$ 上の
$\SheafHom^{n+m}_\mathcal{B}(\mathcal{N},\mathcal{L})$ の切断 $fa$ へ
送る。後者は写像の族
$$
\mathcal{N}^{-q - m}|_U \xrightarrow{a \cdot -}
\mathcal{N}^{-q}|_U \xrightarrow{f_{p, q}}
\mathcal{M}^p|_U
$$
として定義される。この定義が良定義であることの確認は省略する。
この構成は関手および次数付き圏の関手
$$
\SheafHom_\mathcal{B}^{gr}(\mathcal{N}, -) :
\textit{Mod}(\mathcal{B})
\longrightarrow
\textit{Mod}(\mathcal{A})
\quad\text{and}\quad
\SheafHom_\mathcal{B}^{gr}(\mathcal{N}, -) :
\textit{Mod}^{gr}(\mathcal{B})
\longrightarrow
\textit{Mod}^{gr}(\mathcal{A})
$$
を定義する。これは $\mathcal{L}$ から $\mathcal{L}'$ への次数 $n$ の
準同型を、$\SheafHom_\mathcal{B}^{gr}(\mathcal{N},\mathcal{L})$ から
$\SheafHom_\mathcal{B}^{gr}(\mathcal{N},\mathcal{L}')$ への誘導される
次数 $n$ の写像へ送る。

\begin{lemma}
\label{sdga-lemma-tensor-hom-adjunction-gr}
$(\mathcal{C},\mathcal{O})$ を環付きサイトとし、$\mathcal{A}$ と
$\mathcal{B}$ を $(\mathcal{C},\mathcal{O})$ 上の次数付き代数の層とする。
$\mathcal{M}$ を右次数付き $\mathcal{A}$-加群、$\mathcal{N}$ を次数付き
$(\mathcal{A},\mathcal{B})$-双加群、$\mathcal{L}$ を右次数付き
$\mathcal{B}$-加群とする。上の規約のもとで
$$
\Hom_{\textit{Mod}^{gr}(\mathcal{B})}(
\mathcal{M} \otimes_\mathcal{A} \mathcal{N}, \mathcal{L}) =
\Hom_{\textit{Mod}^{gr}(\mathcal{A})}(
\mathcal{M}, \SheafHom_\mathcal{B}^{gr}(\mathcal{N}, \mathcal{L}))
$$
および
$$
\SheafHom_\mathcal{B}^{gr}(
\mathcal{M} \otimes_\mathcal{A} \mathcal{N}, \mathcal{L}) =
\SheafHom_\mathcal{A}^{gr}(
\mathcal{M}, \SheafHom_\mathcal{B}^{gr}(\mathcal{N}, \mathcal{L}))
$$
が成り立ち、これらは $\mathcal{M}$、$\mathcal{N}$、$\mathcal{L}$
について関手的である。
\end{lemma}

\begin{proof}
省略する。ヒント：両辺を、$\mathcal{A}$-双線形で右から
$\mathcal{B}$-線形な次数付き写像
$\psi : \mathcal{M} \times \mathcal{N} \to \mathcal{L}$
として解釈すればよい。
\end{proof}

\noindent
$(\mathcal{C},\mathcal{O})$ を環付きサイトとし、$\mathcal{A}$ と
$\mathcal{B}$ を $(\mathcal{C},\mathcal{O})$ 上の次数付き代数の層とする。
上の特別な場合として、次数付き $\mathcal{O}$-代数の準同型
$\varphi:\mathcal{A}\to\mathcal{B}$ が与えられているとする。
このとき関手および次数付き圏の関手
$$
\otimes_{\mathcal{A}, \varphi} \mathcal{B} :
\textit{Mod}(\mathcal{A})
\longrightarrow
\textit{Mod}(\mathcal{B})
\quad\text{and}\quad
\otimes_{\mathcal{A}, \varphi} \mathcal{B} :
\textit{Mod}^{gr}(\mathcal{A})
\longrightarrow
\textit{Mod}^{gr}(\mathcal{B})
$$
を得る。一方、スカラー制限関手
$$
res_\varphi :
\textit{Mod}(\mathcal{B})
\longrightarrow
\textit{Mod}(\mathcal{A})
\quad\text{and}\quad
res_\varphi :
\textit{Mod}^{gr}(\mathcal{B})
\longrightarrow
\textit{Mod}^{gr}(\mathcal{A})
$$
もある。上の補題を用いると、これらの関手が互いに随伴であることが
示せる（スカラー制限と底変換についての通常の場合と同様である）。
実際、$\mathcal{B}$ を次数付き $(\mathcal{A},\mathcal{B})$-双加群と
みなしたものを ${}_\mathcal{A}\mathcal{B}_\mathcal{B}$ と書く。
任意の右次数付き $\mathcal{B}$-加群 $\mathcal{L}$ に対して
$$
\SheafHom_\mathcal{B}^{gr}({}_\mathcal{A}\mathcal{B}_\mathcal{B}, \mathcal{L})
= res_\varphi(\mathcal{L})
$$
が右次数付き $\mathcal{A}$-加群として成り立つ。したがって Lemma
\ref{sdga-lemma-tensor-hom-adjunction-gr} により、関手的同型
$$
\Hom_{\textit{Mod}^{gr}(\mathcal{B})}(
\mathcal{M} \otimes_{\mathcal{A}, \varphi} \mathcal{B}, \mathcal{L}) =
\Hom_{\textit{Mod}^{gr}(\mathcal{A})}(
\mathcal{M}, res_\varphi(\mathcal{L}))
$$
を得る。$\varphi$ が文脈から明らかな場合、この式では通常
$\varphi$ への依存を省略する。同様にして、次数付き
$\mathcal{O}$-加群の等式
$$
\SheafHom^{gr}_\mathcal{B}(
\mathcal{M} \otimes_\mathcal{A} \mathcal{B}, \mathcal{L}) =
\SheafHom_\mathcal{A}^{gr}(\mathcal{M}, \mathcal{L})
$$
を得る。




\section{次数付き加群層の引き戻しと順像}
\label{sdga-section-functoriality-graded}

\noindent
本節は読み飛ばすことを勧める。

\medskip\noindent
$(f, f^\sharp) : (\Sh(\mathcal{C}), \mathcal{O}_\mathcal{C})
\to (\Sh(\mathcal{D}), \mathcal{O}_\mathcal{D})$
を環付きトポスの射とする。$\mathcal{A}$ を次数付き
$\mathcal{O}_\mathcal{C}$-代数、$\mathcal{B}$ を次数付き
$\mathcal{O}_\mathcal{D}$-代数とする。次数付き
$f^{-1}\mathcal{O}_\mathcal{D}$-代数の準同型
$$
\varphi : f^{-1}\mathcal{B} \to \mathcal{A}
$$
が与えられているとする。スカラー制限とスカラー拡大の随伴により、
これは次数付き $\mathcal{O}_\mathcal{C}$-代数の準同型
$\varphi:f^*\mathcal{B}\to\mathcal{A}$ と同じデータである。同値な言い方を
すれば、$\varphi$ は次数付き $\mathcal{O}_\mathcal{D}$-代数の準同型
$$
\varphi : \mathcal{B} \to f_*\mathcal{A}
$$
とみなせる。Remark \ref{sdga-remark-functoriality-ga} を参照せよ。

\medskip\noindent
関手
$$
f_* :
\textit{Mod}(\mathcal{A})
\longrightarrow
\textit{Mod}(\mathcal{B})
$$
を定義しよう。次数付き $\mathcal{A}$-加群 $\mathcal{M}$ に対して、
$f_*\mathcal{M}$ を、次数 $n$ の項が $f_*\mathcal{M}^n$ である
次数付き $\mathcal{B}$-加群と定義する。乗法として
$$
f_*\mathcal{M}^n \times \mathcal{B}^m
\xrightarrow{(\text{id}, \varphi^m)}
f_*\mathcal{M}^n \times f_*\mathcal{A}^m
\xrightarrow{f_*\mu_{n, m}}
f_*\mathcal{M}^{n + m}
$$
を用いる。ここで $\mu_{n,m}:\mathcal{M}^n\times\mathcal{A}^m
\to\mathcal{M}^{n+m}$ は、$\mathcal{M}$ の $\mathcal{A}$-加群としての
乗法写像である。ここでは $f_*$ が直積と交換することを用いた。
この構成は $\mathcal{M}$ について明らかに関手的であり、所望の関手を得る。

\medskip\noindent
次に関手
$$
f^* :
\textit{Mod}(\mathcal{B})
\longrightarrow
\textit{Mod}(\mathcal{A})
$$
を、関手の合成
$$
\textit{Mod}(\mathcal{B})
\xrightarrow{f^{-1}}
\textit{Mod}(f^{-1}\mathcal{B})
\xrightarrow{ - \otimes_{f^{-1}\mathcal{B}} \mathcal{A}}
\textit{Mod}(\mathcal{A})
$$
として定義する。まず、次数付き $\mathcal{B}$-加群 $\mathcal{N}$ に対して、
$f^{-1}\mathcal{N}$ を、次数 $n$ の項が $f^{-1}\mathcal{N}^n$ である
次数付き $f^{-1}\mathcal{B}$-加群と定義する。乗法として
$$
f^{-1}\nu_{n, m} :
f^{-1}\mathcal{N}^n \times f^{-1}\mathcal{B}^m
\longrightarrow
f^{-1}\mathcal{N}^{n + m}
$$
を用いる。ここで $\nu_{n,m}:\mathcal{N}^n\times\mathcal{B}^m
\to\mathcal{N}^{n+m}$ は、$\mathcal{N}$ の $\mathcal{B}$-加群としての
乗法写像である。ここでは $f^{-1}$ が直積と交換することを用いた。
この構成は $\mathcal{N}$ について明らかに関手的であり、関手
$f^{-1}$ を得る。これを踏まえ、Section
\ref{sdga-section-graded-bimodules} のテンソル積に関する議論を用いて関手
$$
- \otimes_{f^{-1}\mathcal{B}} \mathcal{A} :
\textit{Mod}(f^{-1}\mathcal{B})
\longrightarrow
\textit{Mod}(\mathcal{A})
$$
を定義できる。最後に、上で予告したとおり
$$
f^*\mathcal{N} =
f^{-1}\mathcal{N} \otimes_{f^{-1}\mathcal{B}, \varphi} \mathcal{A}
$$
とおく。

\medskip\noindent
関手 $f_*$ と $f^*$ は、容易に次数付き圏の関手
$$
f_* :
\textit{Mod}^{gr}(\mathcal{A})
\longrightarrow
\textit{Mod}^{gr}(\mathcal{B})
\quad\text{and}\quad
f^* :
\textit{Mod}^{gr}(\mathcal{B})
\longrightarrow
\textit{Mod}^{gr}(\mathcal{A})
$$
へ拡張できる。これらは基礎となる対象には同じ操作を施し、次数 $n$ の
斉次射には上の構成の関手性により定義される操作を施す。

\begin{lemma}
\label{sdga-lemma-adjunction-push-pull-gr}
上の状況では
$$
\Hom_{\textit{Mod}^{gr}(\mathcal{B})}(
\mathcal{N}, f_*\mathcal{M}) =
\Hom_{\textit{Mod}^{gr}(\mathcal{A})}(
f^*\mathcal{N}, \mathcal{M})
$$
が成り立つ。
\end{lemma}

\begin{proof}
省略する。ヒント：まず、アーベル群の層に対する $f^{-1}$ と $f_*$ の
随伴を用いて、$\textit{Mod}(\mathcal{B})$ と
$\textit{Mod}(f^{-1}\mathcal{B})$ の間の関手として $f^{-1}$ と $f_*$ が
随伴であることを示す。次に、Section \ref{sdga-section-graded-bimodules}
で与えた底変換とスカラー制限の随伴を用いる。
\end{proof}





\section{局所化と次数付き加群の層}
\label{sdga-section-localize-graded}

\noindent
本節は読み飛ばすことを勧める。

\medskip\noindent
$(\mathcal{C},\mathcal{O})$ を環付きサイトとし、
$U\in\Ob(\mathcal{C})$ とする。対応する局所化射を
$$
j :
(\Sh(\mathcal{C}/U), \mathcal{O}_U)
\longrightarrow
(\Sh(\mathcal{C}), \mathcal{O})
$$
と書く（Modules on Sites,
Section \ref{sites-modules-section-localize}）。以下では次の事実を用いる。
$\mathcal{O}_U$-加群 $\mathcal{M}_i$（$i=1,2$）と
$\mathcal{O}$-加群 $\mathcal{A}$ に対して、標準写像
$$
j_! :
\Hom_{\mathcal{O}_U}(
\mathcal{M}_1 \otimes_{\mathcal{O}_U} \mathcal{A}|_U, \mathcal{M}_2)
\longrightarrow
\Hom_\mathcal{O}(
j_!\mathcal{M}_1 \otimes_\mathcal{O} \mathcal{A}, j_!\mathcal{M}_2)
$$
がある。実際、Modules on Sites,
Lemma \ref{sites-modules-lemma-j-shriek-and-tensor} により
$j_!(\mathcal{M}_1 \otimes_{\mathcal{O}_U} \mathcal{A}|_U) =
j_!\mathcal{M}_1 \otimes_\mathcal{O} \mathcal{A}$ である。

\medskip\noindent
$\mathcal{A}$ を次数付き $\mathcal{O}$-代数とする。
$\mathcal{A}$ の $\mathcal{C}/U$ への制限を $\mathcal{A}_U$ と書く。
すなわち $\mathcal{A}_U=j^*\mathcal{A}=j^{-1}\mathcal{A}$ である。
Section \ref{sdga-section-functoriality-graded} では随伴関手
$$
j_* :
\textit{Mod}^{gr}(\mathcal{A}_U)
\longrightarrow
\textit{Mod}^{gr}(\mathcal{A})
\quad\text{and}\quad
j^* :
\textit{Mod}^{gr}(\mathcal{A})
\longrightarrow
\textit{Mod}^{gr}(\mathcal{A}_U)
$$
を構成し、$j^*$ は $j_*$ の左随伴であった。さらに、$j^*$ の左随伴である
完全関手
$$
j_! :
\textit{Mod}^{gr}(\mathcal{A}_U)
\longrightarrow
\textit{Mod}^{gr}(\mathcal{A})
$$
も存在する。実際、次数付き $\mathcal{A}_U$-加群 $\mathcal{M}$ に対して、
$j_!\mathcal{M}$ を、次数 $n$ の項が $j_!\mathcal{M}^n$ である
次数付き $\mathcal{A}$-加群と定義する。乗法写像として
$$
j_!\mu_{n, m} :
j_!\mathcal{M}^n \times \mathcal{A}^m \to
j_!\mathcal{M}^{n + m}
$$
を用いる。ここで $\mu_{m,n}:\mathcal{M}^n\times\mathcal{A}^m
\to\mathcal{M}^{n+m}$ は与えられた乗法写像である。次数付き
$\mathcal{A}_U$-加群の次数 $n$ の斉次写像
$f:\mathcal{M}\to\mathcal{M}'$ に対して、次数 $n$ の斉次写像
$j_!f:j_!\mathcal{M}\to j_!\mathcal{M}'$ を得る。これで所望の関手を得る。

\begin{lemma}
\label{sdga-lemma-extension-by-zero-graded}
上の状況では
$$
\Hom_{\textit{Mod}^{gr}(\mathcal{A})}(
j_!\mathcal{M}, \mathcal{N}) =
\Hom_{\textit{Mod}^{gr}(\mathcal{A}_U)}(
\mathcal{M}, j^*\mathcal{N})
$$
が成り立つ。
\end{lemma}

\begin{proof}
Modules on Sites, Section \ref{sites-modules-section-localize} の議論により、
$\mathcal{O}$-加群上の関手 $j_!$ と $j^*$ は随伴である。したがって、
$\mathcal{O}$-加群構造だけをみれば
$$
\Hom_{\text{Gr}^{gr}(\textit{Mod}(\mathcal{O}))}(
j_!\mathcal{M}, \mathcal{N}) =
\Hom_{\text{Gr}^{gr}(\textit{Mod}(\mathcal{O}_U))}(
\mathcal{M}, j^*\mathcal{N})
$$
が成り立つ（$\text{Gr}^{gr}(\textit{Mod}(\mathcal{O}))$ は次数付き
$\mathcal{O}$-加群の次数付き圏を表すことを思い出そう）。さらに、
これらの同一視が左辺の $\mathcal{A}$-加群写像を右辺の
$\mathcal{A}_U$-加群写像へ対応させることを確かめる必要がある。
これを確かめるため、$\mathcal{O}_U$-線形写像
$f^n : \mathcal{M}^n \to j^*\mathcal{N}^{n + d}$
が $\mathcal{O}$-線形写像
$g^n : j_!\mathcal{M}^n \to \mathcal{N}^{n + d}$
に対応するとする。次の図式が可換であることと
$$
\xymatrix{
\mathcal{M}^n \otimes_{\mathcal{O}_U} \mathcal{A}_U^m
\ar[r]_{f^n \otimes 1} \ar[d] &
j^*\mathcal{N}^{n + d} \otimes_{\mathcal{O}_U} \mathcal{A}_U^m \ar[d] \\
\mathcal{M}^{n + m} \ar[r]^{f^{n + m}} &
j^*\mathcal{N}^{n + m + d}
}
$$
次の図式が可換であることとが同値であることを示せば十分である。
$$
\xymatrix{
j_!\mathcal{M}^n \otimes_\mathcal{O} \mathcal{A}^m
\ar[r]_{g^n \otimes 1} \ar[d] &
\mathcal{N}^{n + d} \otimes_\mathcal{O} \mathcal{A}_U^m \ar[d] \\
j_!\mathcal{M}^{n + m} \ar[r]^{g^{n + m}} &
\mathcal{N}^{n + m + d}
}
$$
ところが
\begin{align*}
\Hom_{\mathcal{O}_U}(\mathcal{M}^n \otimes_{\mathcal{O}_U} \mathcal{A}_U^m,
j^*\mathcal{N}^{n + d + m})
& =
\Hom_\mathcal{O}(j_!(\mathcal{M}^n \otimes_{\mathcal{O}_U} \mathcal{A}_U^m),
\mathcal{N}^{n + d + m}) \\
& =
\Hom_\mathcal{O}(j_!\mathcal{M}^n \otimes_\mathcal{O} \mathcal{A}^m,
\mathcal{N}^{n + d + m})
\end{align*}
である。これは上ですでに用いた Modules on Sites,
Lemma \ref{sites-modules-lemma-j-shriek-and-tensor} による。
最初の群における第一の図式の可換性への障害が、最後の群における
第二の図式の可換性への障害へ写ることの確認は省略する。
\end{proof}

\begin{lemma}
\label{sdga-lemma-tensor-with-extension-by-zero}
上の状況で、$\mathcal{M}$ を右次数付き $\mathcal{A}_U$-加群、
$\mathcal{N}$ を左次数付き $\mathcal{A}$-加群とする。このとき
$$
j_!\mathcal{M} \otimes_\mathcal{A} \mathcal{N} =
j_!(\mathcal{M} \otimes_{\mathcal{A}_U} \mathcal{N}|_U)
$$
が次数付き $\mathcal{O}$-加群として成り立ち、$\mathcal{M}$ と
$\mathcal{N}$ について関手的である。
\end{lemma}

\begin{proof}
$j_!\mathcal{M}\otimes_\mathcal{A}\mathcal{N}$ の次数 $n$ の成分は、
標準写像
$$
\bigoplus\nolimits_{r + s + t = n}
j_!\mathcal{M}^r \otimes_\mathcal{O}
\mathcal{A}^s \otimes_\mathcal{O}
\mathcal{N}^t
\longrightarrow
\bigoplus\nolimits_{p + q = n}
j_!\mathcal{M}^p \otimes_\mathcal{O} \mathcal{N}^q
$$
の余核であった。Section \ref{sdga-section-tensor-product} を参照せよ。
Modules on Sites, Lemma \ref{sites-modules-lemma-j-shriek-and-tensor}
により、これは写像
$$
\bigoplus\nolimits_{r + s + t = n}
j_!(\mathcal{M}^r \otimes_{\mathcal{O}_U}
\mathcal{A}^s|_U \otimes_{\mathcal{O}_U}
\mathcal{N}^t|_U)
\longrightarrow
\bigoplus\nolimits_{p + q = n}
j_!(\mathcal{M}^p \otimes_{\mathcal{O}_U} \mathcal{N}^q|_U)
$$
の余核と同じであり、結論を得る。別の証明として、上で引用した補題の
証明における米田の議論を繰り返すこともできる。
\end{proof}








\section{次数付き加群の層上のシフト関手}
\label{sdga-section-shift}

\noindent
本節は読み飛ばすことを強く勧める。次数付き代数上の次数付き加群の層は、
Differential Graded Algebra,
Remark \ref{dga-remark-graded-shift-functors} で論じた現象の一例である。

\medskip\noindent
$(\mathcal{C},\mathcal{O})$ を環付きサイトとし、$\mathcal{A}$ を
$(\mathcal{C},\mathcal{O})$ 上の次数付き代数の層、$\mathcal{M}$ を
次数付き $\mathcal{A}$-加群、$k\in\mathbf{Z}$ とする。
$\mathcal{M}$ の{\it 第 $k$ シフト} $\mathcal{M}[k]$ を、次数 $n$ の部分が
$$
(\mathcal{M}[k])^n = \mathcal{M}^{n + k}
$$
すなわち $\mathcal{M}$ の次数 $n+k$ の部分である次数付き
$\mathcal{A}$-加群として定義する。その乗法写像
$$
(\mathcal{M}[k])^n \times \mathcal{A}^m
\longrightarrow
(\mathcal{M}[k])^{n + m}
$$
には、単に $\mathcal{M}$ の乗法写像
$$
\mathcal{M}^{n + k} \times \mathcal{A}^m
\longrightarrow
\mathcal{M}^{n + m + k}
$$
を用いる。これが関手 $[k]$ を定義すること、
$[k+l]=[k]\circ[l]$ であること、および
$$
\Hom_{\textit{Mod}^{gr}(\mathcal{A})}(\mathcal{L}, \mathcal{M}[k]) =
\Hom_{\textit{Mod}^{gr}(\mathcal{A})}(\mathcal{L}, \mathcal{M})[k]
$$
が（符号を介さずに）$\mathcal{M}$ と $\mathcal{L}$ について関手的に
成り立つことは明らかである。したがって Differential Graded Algebra,
Remark \ref{dga-remark-graded-shift-functors} で論じたように、次数付き
$\mathcal{A}$-加群の次数付き圏は、次数付き $\mathcal{A}$-加群の通常の圏と
シフト関手から実際に復元できる。

\begin{lemma}
\label{sdga-lemma-gm-grothendieck-abelian}
$(\mathcal{C},\mathcal{O})$ を環付きサイトとし、$\mathcal{A}$ を
次数付き $\mathcal{O}$-代数とする。圏 $\textit{Mod}(\mathcal{A})$
はグロタンディーク・アーベル圏である。
\end{lemma}

\begin{proof}
Lemma \ref{sdga-lemma-gm-abelian} とグロタンディーク・アーベル圏の定義
（Injectives, Definition \ref{injectives-definition-grothendieck-conditions}）
により、$\textit{Mod}(\mathcal{A})$ が生成対象をもつことを示せば十分である。
$$
\mathcal{G} = \bigoplus\nolimits_{k, U} j_{U!}\mathcal{A}_U[k]
$$
は生成対象であると主張する。ここで直和は $\mathcal{C}$ のすべての対象
$U$ と $k\in\mathbf{Z}$ にわたる。実際、次数付き
$\mathcal{A}$-加群 $\mathcal{M}$ に対して、$\mathcal{G}$ から
$\mathcal{M}$ への非零写像が存在しないならば、すべての $k$ と $U$ について
$$
\Hom_{\textit{Mod}(\mathcal{A})}(j_{U!}\mathcal{A}_U[k], \mathcal{M}) =
\Hom_{\textit{Mod}(\mathcal{A}_U)}(\mathcal{A}_U[k], \mathcal{M}|_U) =
\Gamma(U, \mathcal{M}^{-k})
$$
は零である。したがって $\mathcal{M}$ は零である。
\end{proof}









\section{微分次数付き代数の層}
\label{sdga-section-dga}

\noindent
本節は Differential Graded Algebra,
Section \ref{dga-section-dga} の類似である。

\begin{definition}
\label{sdga-definition-dga}
$(\mathcal{C},\mathcal{O})$ を環付きサイトとする。
$(\mathcal{C},\mathcal{O})$ 上の{\it 微分次数付き
$\mathcal{O}$-代数の層}、または{\it 微分次数付き代数の層}とは、
$\mathcal{O}$-加群の余鎖複体 $\mathcal{A}^\bullet$ と
$\mathcal{O}$-双線形写像
$$
\mathcal{A}^n \times \mathcal{A}^m \to \mathcal{A}^{n + m},\quad
(a, b) \longmapsto ab
$$
であって、次の性質を満たすものとして与えられる。これらの写像を
乗法写像という。
\begin{enumerate}
\item 乗法は結合的である。
\item $\mathcal{A}^0$ には、乗法の両側単位元となる大域切断 $1$ がある。
\item $U\in\Ob(\mathcal{C})$、$a\in\mathcal{A}^n(U)$、
$b\in\mathcal{A}^m(U)$ に対して
$$
\text{d}^{n + m}(ab) = \text{d}^n(a)b + (-1)^n a\text{d}^m(b)
$$
が成り立つ。
\end{enumerate}
このような構造をしばしば $(\mathcal{A},\text{d})$ と書く。
$(\mathcal{A},\text{d})$ から $(\mathcal{B},\text{d})$ への
{\it 微分次数付き $\mathcal{O}$-代数の準同型}とは、乗法写像と両立する
$\mathcal{O}$-加群の複体の写像
$f:\mathcal{A}^\bullet\to\mathcal{B}^\bullet$ のことである。
\end{definition}

\noindent
微分次数付き $\mathcal{O}$-代数 $(\mathcal{A},\text{d})$ と対象
$U\in\Ob(\mathcal{C})$ に対して、記法
$$
\mathcal{A}(U) =
\Gamma(U, \mathcal{A}) =
\bigoplus\nolimits_{n \in \mathbf{Z}} \mathcal{A}^n(U)
$$
を用いる。これは微分次数付き $\mathcal{O}(U)$-代数である。

\medskip\noindent
可能な限り、微分次数付き $\mathcal{O}$-代数
$(\mathcal{A},\text{d})$ を、定義の Leibniz 則を満たす次数 $1$ の作用素
$\text{d}:\mathcal{A}\to\mathcal{A}$ を備えた次数付き
$\mathcal{O}$-代数 $\mathcal{A}$ とみなす（ここで $\mathcal{A}$ は
次数付き $\mathcal{O}$-加群とみなしている）。

\begin{remark}
\label{sdga-remark-functoriality-dga}
$(f,f^\sharp):(\Sh(\mathcal{C}),\mathcal{O}_\mathcal{C})
\to (\Sh(\mathcal{D}), \mathcal{O}_\mathcal{D})$
を環付きトポスの射とする。
\begin{enumerate}
\item $(\mathcal{A},\text{d})$ を微分次数付き
$\mathcal{O}_\mathcal{C}$-代数とする。その順像は微分次数付き
$\mathcal{O}_\mathcal{D}$-代数 $(f_*\mathcal{A},\text{d})$ である。
ここで $f_*\mathcal{A}$ は Remark \ref{sdga-remark-functoriality-ga} のとおりで、
写像 $f_*\mathcal{A}^n\to f_*\mathcal{A}^{n+1}$ として
$\text{d}=f_*\text{d}$ である。Leibniz 則の確認は省略する。
\item $\mathcal{B}$ を微分次数付き $\mathcal{O}_\mathcal{D}$-代数とする。
その引き戻しは微分次数付き $\mathcal{O}_\mathcal{C}$-代数
$(f^*\mathcal{B},\text{d})$ である。ここで $f^*\mathcal{B}$ は Remark
\ref{sdga-remark-functoriality-ga} のとおりで、写像
$f^*\mathcal{B}^n\to f^*\mathcal{B}^{n+1}$ として
$\text{d}=f^*\text{d}$ である。Leibniz 則の確認は省略する。
\item 微分次数付き $\mathcal{O}_\mathcal{C}$-代数の準同型
$f^*\mathcal{B}\to\mathcal{A}$ の集合と、微分次数付き
$\mathcal{O}_\mathcal{D}$-代数の準同型
$\mathcal{B}\to f_*\mathcal{A}$ の集合との間には一対一対応がある。
\end{enumerate}
(3) は加群の層上の $f^*$ と $f_*$ の通常の随伴から直ちに従う。
\end{remark}
















\section{微分次数付き加群の層}
\label{sdga-section-modules}

\noindent
本節は Differential Graded Algebra,
Section \ref{dga-section-modules} の類似である。

\begin{definition}
\label{sdga-definition-dgm}
$(\mathcal{C},\mathcal{O})$ を環付きサイトとし、
$(\mathcal{A},\text{d})$ を $(\mathcal{C},\mathcal{O})$ 上の
微分次数付き代数の層とする。$\mathcal{A}$ 上の（右）
{\it 微分次数付き $\mathcal{A}$-加群}、または（右）
{\it 微分次数付き加群}とは、余鎖複体 $\mathcal{M}^\bullet$ と
$\mathcal{O}$-双線形写像
$$
\mathcal{M}^n \times \mathcal{A}^m \to \mathcal{M}^{n + m},\quad
(x, a) \longmapsto xa
$$
であって、次の性質を満たすものとして与えられる。これらの写像を
乗法写像という。
\begin{enumerate}
\item 乗法は $(xa)a'=x(aa')$ を満たす。
\item $\mathcal{A}^0$ の単位切断 $1$ は、すべての $n$ に対して
$\mathcal{M}^n$ 上に恒等写像として作用する。
\item $U\in\Ob(\mathcal{C})$、$x\in\mathcal{M}^n(U)$、
$a\in\mathcal{A}^m(U)$ に対して
$$
\text{d}^{n + m}(xa) = \text{d}^n(x)a + (-1)^n x\text{d}^m(a)
$$
が成り立つ。
\end{enumerate}
この状況を表すために、単に「$\mathcal{M}$ を微分次数付き
$\mathcal{A}$-加群とする」としばしばいう。
$\mathcal{M}$ から $\mathcal{N}$ への
{\it 微分次数付き $\mathcal{A}$-加群の準同型}とは、乗法写像と両立する
$\mathcal{O}$-加群の複体の写像
$f:\mathcal{M}^\bullet\to\mathcal{N}^\bullet$ のことである。
（右）微分次数付き $\mathcal{A}$-加群の圏を
$\textit{Mod}(\mathcal{A},\text{d})$ と書く。
\end{definition}

\noindent
{\it 左微分次数付き加群}もまったく同様に定義できるが、本章では
特に断らない限り右加群を用いる。

\medskip\noindent
微分次数付き $\mathcal{A}$-加群 $\mathcal{M}$ と対象
$U\in\Ob(\mathcal{C})$ に対して、記法
$$
\mathcal{M}(U) =
\Gamma(U, \mathcal{M}) =
\bigoplus\nolimits_{n \in \mathbf{Z}} \mathcal{M}^n(U)
$$
を用いる。これは（右）微分次数付き $\mathcal{A}(U)$-加群である。

\begin{lemma}
\label{sdga-lemma-dgm-abelian}
$(\mathcal{C},\mathcal{O})$ を環付きサイトとし、
$(\mathcal{A},\text{d})$ を微分次数付き $\mathcal{O}$-代数とする。
圏 $\textit{Mod}(\mathcal{A},\text{d})$ は次の性質をもつアーベル圏である。
\begin{enumerate}
\item $\textit{Mod}(\mathcal{A},\text{d})$ は任意の直和をもつ。
\item $\textit{Mod}(\mathcal{A},\text{d})$ は任意の余極限をもつ。
\item $\textit{Mod}(\mathcal{A},\text{d})$ におけるフィルター余極限は完全である。
\item $\textit{Mod}(\mathcal{A},\text{d})$ は任意の直積をもつ。
\item $\textit{Mod}(\mathcal{A},\text{d})$ は任意の極限をもつ。
\end{enumerate}
微分次数付き $\mathcal{A}$-加群をその基礎となる次数付き加群へ送る忘却関手
$$
\textit{Mod}(\mathcal{A}, \text{d})
\longrightarrow
\textit{Mod}(\mathcal{A})
$$
はすべての極限および余極限と交換する。
\end{lemma}

\begin{proof}
補題の主張に現れる函子を
$F : \textit{Mod}(\mathcal{A}, \text{d}) \to \textit{Mod}(\mathcal{A})$
と書く。圏 $\textit{Mod}(\mathcal{A})$ は性質 (1)--(5) をもつことに
注意せよ。Lemma \ref{sdga-lemma-gm-abelian} を参照せよ。

\medskip\noindent
次数付き $\mathcal{A}$-加群の準同型
$f:\mathcal{M}\to\mathcal{N}$ を考える。$f$ を次数付き
$\mathcal{A}$-加群の写像とみたときの核と余核には、定義
\ref{sdga-definition-dgm} の微分も備わる。したがって、これらは
$\textit{Mod}(\mathcal{A},\text{d})$ における核と余核でもある。
ゆえに $\textit{Mod}(\mathcal{A},\text{d})$ はアーベル圏であり、
核と余核を取る操作は $F$ と交換する。

\medskip\noindent
極限と余極限の存在を証明するには、直積と直和の存在を示せば十分である。
Categories, Lemmas \ref{categories-lemma-limits-products-equalizers} and
\ref{categories-lemma-colimits-coproducts-coequalizers} を参照せよ。
同じ補題により、$F$ が直和および直積と交換することを示せば、
極限および余極限とも交換することが従う。

\medskip\noindent
$\mathcal{M}_t$（$t\in T$）を微分次数付き $\mathcal{A}$-加群の族とする。
直和 $\bigoplus\mathcal{M}_t$ を次数付き $\mathcal{A}$-加群として考える。
次数付き加群の直和は項ごとに取られるので、
$\bigoplus\mathcal{M}_t$ に微分が備わることは明らかである。
これが $\textit{Mod}(\mathcal{A},\text{d})$ における直和であることは
容易に確かめられる。直積についても同様である。

\medskip\noindent
$F$ は完全関手であり、$\textit{Mod}(\mathcal{A},\text{d})$ の複体
$\mathcal{M}_1\to\mathcal{M}_2\to\mathcal{M}_3$ が完全であることと、
$F(\mathcal{M}_1) \to F(\mathcal{M}_2) \to F(\mathcal{M}_3)$ が
$\textit{Mod}(\mathcal{A})$ で完全であることとは同値である。
$\textit{Mod}(\mathcal{A})$ におけるフィルター余極限は完全なので、
(3) が従う。
\end{proof}

\noindent
Lemmas \ref{sdga-lemma-dgm-abelian} and \ref{sdga-lemma-gm-abelian} を合わせると、
アーベル圏の間の完全忠実関手
$$
\textit{Mod}(\mathcal{A}, \text{d}) \longrightarrow \text{Comp}(\mathcal{O})
$$
を得る。微分次数付き $\mathcal{A}$-加群 $\mathcal{M}$ に対して、
コホモロジー $\mathcal{O}$-加群 $H^i(\mathcal{M})$ を、$\mathcal{M}$ に
対応する $\mathcal{O}$-加群の複体のコホモロジーとして定義する。
したがって、微分次数付き $\mathcal{A}$-加群の短完全列
$0 \to \mathcal{K} \to \mathcal{L} \to \mathcal{M} \to 0$ は、
コホモロジー加群の長完全列
\begin{equation}
\label{sdga-equation-les}
H^n(\mathcal{K}) \to H^n(\mathcal{L}) \to H^n(\mathcal{M}) \to
H^{n + 1}(\mathcal{K})
\end{equation}
を与える。Homology,
Lemma \ref{homology-lemma-long-exact-sequence-cochain} を参照せよ。

\medskip\noindent
さらに、以後は加群の複体に用いるすべての用語を流用する。例えば、
すべての $k\in\mathbf{Z}$ に対して $H^k(\mathcal{M})=0$ ならば、
微分次数付き $\mathcal{A}$-加群 $\mathcal{M}$ は{\it 非輪状}であるという。
微分次数付き $\mathcal{A}$-加群の準同型
$\mathcal{M}\to\mathcal{N}$ がすべての $k\in\mathbf{Z}$ に対して同型
$H^k(\mathcal{M})\to H^k(\mathcal{N})$ を誘導するならば、これを
{\it 擬同型}という。以下同様である。


















\section{加群の微分次数付き圏}
\label{sdga-section-dgm-dg-cat}

\noindent
本節は Differential Graded Algebra,
Example \ref{dga-example-dgm-dg-cat} の類似である。微分次数付き圏に関する
規約については Differential Graded Algebra,
Section \ref{dga-section-dga-categories} を参照せよ。

\medskip\noindent
$(\mathcal{C},\mathcal{O})$ を環付きサイトとし、
$(\mathcal{A},\text{d})$ を $(\mathcal{C},\mathcal{O})$ 上の
微分次数付き代数の層とする。$R=\Gamma(\mathcal{C},\mathcal{O})$ 上の
微分次数付き圏
$$
\textit{Mod}^{dg}(\mathcal{A}, \text{d})
$$
であって、その付随する複体の圏が微分次数付き
$\mathcal{A}$-加群の圏
$$
\textit{Mod}(\mathcal{A}, \text{d}) =
\text{Comp}(\textit{Mod}^{dg}(\mathcal{A}, \text{d}))
$$
となるものを構成する。$\textit{Mod}^{dg}(\mathcal{A},\text{d})$ の
対象として右微分次数付き $\mathcal{A}$-加群を取る。Section
\ref{sdga-section-modules} を参照せよ。微分次数付き $\mathcal{A}$-加群
$\mathcal{L}$ と $\mathcal{M}$ に対して、次数付き $R$-加群として
$$
\Hom_{\textit{Mod}^{dg}(\mathcal{A}, \text{d})}(\mathcal{L}, \mathcal{M}) =
\Hom_{\textit{Mod}^{gr}(\mathcal{A})}(\mathcal{L}, \mathcal{M}) =
\bigoplus\nolimits_{n \in \mathbf{Z}} \Hom^n(\mathcal{L}, \mathcal{M})
$$
とおく。Section \ref{sdga-section-gm-gr-cat} を参照せよ。言い換えれば、
次数 $n$ の部分 $\Hom^n(\mathcal{L},\mathcal{M})$ は、次数 $n$ の
斉次な右 $\mathcal{A}$-加群写像からなる $R$-加群である。
$f\in\Hom^n(\mathcal{L},\mathcal{M})$ に対して
$$
\text{d}(f) =
\text{d}_\mathcal{M} \circ f - (-1)^n f \circ \text{d}_\mathcal{L}
$$
とおく。この式を解釈するために、$\text{d}_\mathcal{M}$ と
$\text{d}_\mathcal{L}$ を次数付き $\mathcal{O}$-加群写像とみなし、
次数付き $\mathcal{O}$-加群写像の合成を用いる。$\text{d}(f)$ が次数付き
$\mathcal{O}$-加群写像として次数 $n+1$ の斉次写像であることは明らかである。
また $\mathcal{A}$-線形である。実際、$\mathcal{M}$ と $\mathcal{A}$ の
斉次な局所切断 $x$ と $a$ に対して
\begin{align*}
\text{d}(f)(xa)
& =
\text{d}_\mathcal{M}(f(x) a) - (-1)^n f (\text{d}_\mathcal{L}(xa)) \\
& =
\text{d}_\mathcal{M}(f(x)) a + (-1)^{\deg(x) + n} f(x) \text{d}(a) 
- (-1)^n f(\text{d}_\mathcal{L}(x)) a - (-1)^{n + \deg(x)} f(x) \text{d}(a) \\
& = \text{d}(f)(x) a
\end{align*}
となる（$\Hom$ 上の微分の定義にこの符号がなければ、この計算は
成り立たないことに注意せよ）。

\medskip\noindent
微分次数付き $\mathcal{A}$-加群 $\mathcal{K}$、$\mathcal{L}$、
$\mathcal{M}$ に対して、Section \ref{sdga-section-gm-gr-cat} では層の写像の
通常の合成により、すでに合成
$$
\Hom^m(\mathcal{L}, \mathcal{M}) \times
\Hom^n(\mathcal{K}, \mathcal{L}) \longrightarrow
\Hom^{n + m}(\mathcal{K}, \mathcal{M})
$$
を定義した。これは微分次数付き加群の写像
$$
\Hom_{\textit{Mod}^{dg}(\mathcal{A}, \text{d})}(\mathcal{L}, \mathcal{M})
\otimes_R
\Hom_{\textit{Mod}^{dg}(\mathcal{A}, \text{d})}(\mathcal{K}, \mathcal{L})
\longrightarrow
\Hom_{\textit{Mod}^{dg}(\mathcal{A}, \text{d})}(\mathcal{K}, \mathcal{M})
$$
を定義する。これは Differential Graded Algebra,
Definition \ref{dga-definition-dga-category} で要求されるとおりである。
実際、
\begin{align*}
\text{d}(g \circ f) & =
\text{d}_\mathcal{M} \circ g \circ f
- (-1)^{n + m} g \circ f \circ \text{d}_\mathcal{K} \\
& =
\left(\text{d}_\mathcal{M} \circ g - (-1)^m g \circ \text{d}_L\right) \circ f
+ (-1)^m g \circ \left(\text{d}_\mathcal{L} \circ f
- (-1)^n f \circ \text{d}_\mathcal{K}\right) \\
& =
\text{d}(g) \circ f + (-1)^m g \circ \text{d}(f)
\end{align*}
である。ここで $f$ の次数は $n$、$g$ の次数は $m$ であり、所望の式を得る。






\section{微分次数付き加群の層のテンソル積}
\label{sdga-section-tensor-product-dg}

\noindent
本節は Differential Graded Algebra,
Section \ref{dga-section-tensor-product} の一部の類似である。

\medskip\noindent
$(\mathcal{C},\mathcal{O})$ を環付きサイトとし、
$(\mathcal{A},\text{d})$ を $(\mathcal{C},\mathcal{O})$ 上の
微分次数付き代数の層とする。$\mathcal{M}$ を右微分次数付き
$\mathcal{A}$-加群、$\mathcal{N}$ を左微分次数付き
$\mathcal{A}$-加群とする。この状況で{\it テンソル積}
$\mathcal{M}\otimes_\mathcal{A}\mathcal{N}$ を次のように定義する。
次数付き $\mathcal{O}$-加群としては Section
\ref{sdga-section-tensor-product} の構成で与えられ、微分
$$
\text{d}_{\mathcal{M} \otimes_\mathcal{A} \mathcal{N}} :
(\mathcal{M} \otimes_\mathcal{A} \mathcal{N})^n
\longrightarrow
(\mathcal{M} \otimes_\mathcal{A} \mathcal{N})^{n + 1}
$$
を備える。この微分は、$\mathcal{M}$ と $\mathcal{N}$ の斉次な局所切断
$x$ と $y$ に対して
$$
\text{d}_{\mathcal{M} \otimes_\mathcal{A} \mathcal{N}}(x \otimes y) =
\text{d}_\mathcal{M}(x) \otimes y +
(-1)^{\deg(x)}x \otimes \text{d}_\mathcal{N}(y)
$$
という規則で定義する。これが良定義であることを示すには、
$\text{d}_{\mathcal{M}\otimes_\mathcal{A}\mathcal{N}}$ が
$xa\otimes y-x\otimes ay$ の形の元を零へ写すことを示さなければならない。
ここで $x$、$a$、$y$ はそれぞれ $\mathcal{M}$、$\mathcal{A}$、
$\mathcal{N}$ の斉次な局所切断である。計算すると
\begin{align*}
&
\text{d}_{\mathcal{M} \otimes_\mathcal{A} \mathcal{N}}(
xa \otimes y - x \otimes ay) \\
& =
\text{d}_\mathcal{M}(xa) \otimes y + (-1)^{\deg(x) + \deg(a)}
xa \otimes \text{d}_\mathcal{N}(y)
-\text{d}_\mathcal{M}(x) \otimes ay - (-1)^{\deg(x)}
x \otimes \text{d}_\mathcal{N}(ay) \\
& =
\text{d}_\mathcal{M}(x)a \otimes y + (-1)^{\deg(x)}x\text{d}(a) \otimes y
+ (-1)^{\deg(x) + \deg(a)} xa \otimes \text{d}_\mathcal{N}(y) \\
&
-\text{d}_\mathcal{M}(x) \otimes ay - (-1)^{\deg(x)}
x \otimes \text{d}(a)y - (-1)^{\deg(x) +
\deg(a)} x\otimes a\text{d}_\mathcal{N}(y)
\end{align*}
となる。ここで元
$$
\text{d}_\mathcal{M}(x)a \otimes y - \text{d}_\mathcal{M}(x) \otimes ay,\quad
x\text{d}(a) \otimes y - x \otimes \text{d}(a)y,\quad
\text{and}\quad
xa \otimes \text{d}_\mathcal{N}(y) - x\otimes a\text{d}_\mathcal{N}(y)
$$
はいずれも $\mathcal{M}\otimes_\mathcal{A}\mathcal{N}$ で零へ写るので、
結論を得る。等式
$\text{d}_{\mathcal{M} \otimes_\mathcal{A} \mathcal{N}} \circ
\text{d}_{\mathcal{M} \otimes_\mathcal{A} \mathcal{N}} = 0$
の確認は省略する。

\medskip\noindent
左微分次数付き $\mathcal{A}$-加群 $\mathcal{N}$ を固定すると、関手
$$
- \otimes_\mathcal{A} \mathcal{N} :
\textit{Mod}(\mathcal{A}, \text{d})
\longrightarrow
\text{Comp}(\mathcal{O})
$$
を得る。右辺は $\mathcal{O}$-加群の複体の圏である。これは微分次数付き圏の関手
$$
- \otimes_\mathcal{A} \mathcal{N} :
\textit{Mod}^{dg}(\mathcal{A}, \text{d})
\longrightarrow
\text{Comp}^{dg}(\mathcal{O})
$$
へ拡張できる。基礎となる次数付き対象上では、次数 $n$ の準同型
$f:\mathcal{M}\to\mathcal{M}'$ を次数 $n$ の写像
$f \otimes \text{id}_\mathcal{N} :
\mathcal{M} \otimes_\mathcal{A} \mathcal{N} \to
\mathcal{M}' \otimes_\mathcal{A} \mathcal{N}$ へ送る。これは Section
\ref{sdga-section-tensor-product} で行った構成そのものである。この構成が
機能することを示すには、写像
$$
\Hom_{\textit{Mod}^{dg}(\mathcal{A}, \text{d})}(\mathcal{M}, \mathcal{M}')
\longrightarrow
\Hom_{\text{Comp}^{dg}(\mathcal{O})}(
\mathcal{M} \otimes_\mathcal{A} \mathcal{N},
\mathcal{M}' \otimes_\mathcal{A} \mathcal{N})
$$
が微分と両立することを確かめなければならない。上の $f$ に対して、
$$
(\text{d}_{\mathcal{M}'} \circ f - (-1)^n f \circ \text{d}_\mathcal{M})
\otimes \text{id}_\mathcal{N}
$$
が
$$
\text{d}_{\mathcal{M}' \otimes_\mathcal{A} \mathcal{N}}
\circ (f \otimes \text{id}_\mathcal{N})
- (-1)^n (f \otimes \text{id}_\mathcal{N}) \circ
\text{d}_{\mathcal{M} \otimes_\mathcal{A} \mathcal{N}}
$$
に等しいことを示す必要がある。これらの作用素が
$x\otimes y$ の形の局所切断に及ぼす作用を計算する。ここで $x$ と $y$ は
それぞれ $\mathcal{M}$ と $\mathcal{N}$ の斉次な局所切断である。
第一の作用素からは
$$
(\text{d}_{\mathcal{M}'}(f(x)) - (-1)^n f(\text{d}_\mathcal{M}(x))) \otimes y
$$
を得て、第二の作用素からは
\begin{align*}
&\text{d}_{\mathcal{M}' \otimes_\mathcal{A} \mathcal{N}}(f(x) \otimes y)
- (-1)^n (f \otimes \text{id}_\mathcal{N})(
\text{d}_{\mathcal{M} \otimes_\mathcal{A} \mathcal{N}}(x \otimes y) \\
& =
\text{d}_{\mathcal{M}'}(f(x)) \otimes y +
(-1)^{\deg(x) + n}f(x) \otimes \text{d}_\mathcal{N}(y) \\
&
-(-1)^n f(\text{d}_\mathcal{M}(x)) \otimes y
-(-1)^n (-1)^{\deg(x)}f(x) \otimes \text{d}_\mathcal{N}(y)
\end{align*}
を得る。これは確かに同じ局所切断である。










\section{微分次数付き加群の層の内部 Hom}
\label{sdga-section-internal-hom-dg}

\noindent
Section \ref{sdga-section-dgm-dg-cat} の構成を層化したものが必要となる。
$(\mathcal{C},\mathcal{O})$、$\mathcal{A}$、$\mathcal{M}$、
$\mathcal{L}$ を Section \ref{sdga-section-dgm-dg-cat} のとおりとする。
このとき、次数付き $\mathcal{O}$-加群として
$$
\SheafHom^{dg}_\mathcal{A}(\mathcal{M}, \mathcal{L}) =
\SheafHom^{gr}_\mathcal{A}(\mathcal{M}, \mathcal{L}) =
\bigoplus\nolimits_{n \in \mathbf{Z}}
\SheafHom^n_\mathcal{A}(\mathcal{M}, \mathcal{L})
$$
と定義する。Section \ref{sdga-section-internal-hom-graded} を参照せよ。
言い換えれば、$U$ 上の次数 $n$ の部分
$\SheafHom^n_\mathcal{A}(\mathcal{L},\mathcal{M})$ の切断 $f$ は、
次数 $n$ の斉次な右 $\mathcal{A}_U$-加群写像
$\mathcal{L}|_U\to\mathcal{M}|_U$ である。このような $f$ に対して
$$
\text{d}(f) =
\text{d}_\mathcal{M}|_U \circ f - (-1)^n f \circ \text{d}_\mathcal{L}|_U
$$
とおく。この式を解釈するために、$\text{d}_\mathcal{M}|_U$ と
$\text{d}_\mathcal{L}|_U$ を次数付き $\mathcal{O}_U$-加群写像とみなし、
次数付き $\mathcal{O}_U$-加群写像の合成を用いる。
$\text{d}(f)$ が次数付き $\mathcal{O}_U$-加群写像として次数 $n+1$ の
斉次写像であることは明らかである。Section
\ref{sdga-section-dgm-dg-cat} とまったく同じ計算により、
$\text{d}(f)$ は $\mathcal{A}_U$-線形である。

\medskip\noindent
Section \ref{sdga-section-dgm-dg-cat} と同様に、合成写像
$$
\SheafHom^{dg}_\mathcal{A}(\mathcal{L}, \mathcal{M}) \otimes_\mathcal{O}
\SheafHom^{dg}_\mathcal{A}(\mathcal{K}, \mathcal{L})
\longrightarrow
\SheafHom^{dg}_\mathcal{A}(\mathcal{K}, \mathcal{M})
$$
がある。左辺は Section \ref{sdga-section-tensor-product-dg} で定義した
微分次数付き $\mathcal{O}$-加群のテンソル積である。この写像は合成写像
$$
\SheafHom^m(\mathcal{L}, \mathcal{M}) \otimes_\mathcal{O}
\SheafHom^n(\mathcal{K}, \mathcal{L}) \longrightarrow
\SheafHom^{n + m}(\mathcal{K}, \mathcal{M})
$$
で与えられ、これは（局所的には）通常の合成として定義される。
局所切断上で Section \ref{sdga-section-dgm-dg-cat} とまったく同じ計算を
行うと、合成写像が微分次数付き $\mathcal{O}$-加群の準同型であることが
わかる。

\medskip\noindent
これらの定義により、合成と両立する次数付き $R$-加群として
$$
\Hom_{\textit{Mod}^{dg}(\mathcal{A})}(\mathcal{L}, \mathcal{M}) =
\Gamma(\mathcal{C}, \SheafHom^{dg}_\mathcal{A}(\mathcal{L}, \mathcal{M}))
$$
が成り立つ。









\section{微分次数付き双加群の層とテンソル--Hom 随伴}
\label{sdga-section-dg-bimodules}

\noindent
本節は Differential Graded Algebra, Section
\ref{dga-section-tensor-product} の一部に対応する。

\begin{definition}
\label{sdga-definition-dg-bimodule}
$(\mathcal{C}, \mathcal{O})$ を環付きサイトとする。
$\mathcal{A}$ と $\mathcal{B}$ を
$(\mathcal{C}, \mathcal{O})$ 上の微分次数付き代数の層とする。
{\it 微分次数付き $(\mathcal{A}, \mathcal{B})$-双加群}とは、
$\mathcal{O}$-加群の複体 $\mathcal{M}^\bullet$ であって、
$\mathcal{O}$-双線型写像
$$
\mathcal{M}^n \times \mathcal{B}^m \to \mathcal{M}^{n + m},\quad
(x, b) \longmapsto xb
$$
および
$$
\mathcal{A}^n \times \mathcal{M}^m \to \mathcal{M}^{n + m},\quad
(a, x) \longmapsto ax
$$
を備えるものをいう。これらを乗法写像と呼び、次の性質を課す。
\begin{enumerate}
\item 乗法は $a(a'x) = (aa')x$ および
$(xb)b' = x(bb')$ を満たす。
\item $(ax)b = a(xb)$ である。
\item $\text{d}(ax) = \text{d}(a) x + (-1)^{\deg(a)}a \text{d}(x)$ および
$\text{d}(xb) = \text{d}(x) b + (-1)^{\deg(x)}x \text{d}(b)$ である。
\item $\mathcal{A}^0$ の単位切断 $1$ は乗法により恒等写像として作用する。
\item $\mathcal{B}^0$ の単位切断 $1$ は乗法により恒等写像として作用する。
\end{enumerate}
このような構造をしばしば $\mathcal{M}$ と記し、ときには
${}_\mathcal{A}\mathcal{M}_\mathcal{B}$ とも書く。
微分次数付き $(\mathcal{A}, \mathcal{B})$-双加群の
{\it 準同型} $f : \mathcal{M} \to \mathcal{N}$ とは、
乗法写像と両立する $\mathcal{O}$-加群の複体の写像
$f : \mathcal{M}^\bullet \to \mathcal{N}^\bullet$ のことである。
\end{definition}

\noindent
微分次数付き $(\mathcal{A}, \mathcal{B})$-双加群 $\mathcal{M}$ と
対象 $U \in \Ob(\mathcal{C})$ に対して、記法
$$
\mathcal{M}(U) =
\Gamma(U, \mathcal{M}) =
\bigoplus\nolimits_{n \in \mathbf{Z}} \mathcal{M}^n(U)
$$
を用いる。これは微分次数付き
$(\mathcal{A}(U), \mathcal{B}(U))$-双加群である。

\medskip\noindent
微分次数付き $(\mathcal{A}, \mathcal{B})$-双加群
$\mathcal{M}$ は、右微分次数付き $\mathcal{B}$-加群であると同時に
左微分次数付き $\mathcal{A}$-加群であり、両者の次数付けと微分が一致し、
かつ $\mathcal{A}$-加群構造と $\mathcal{B}$-加群構造が可換であるものに
ほかならない。正確には次のように述べられる。

\begin{lemma}
\label{sdga-lemma-what-makes-a-bimodule-dg}
$(\mathcal{C}, \mathcal{O})$ を環付きサイトとする。
$\mathcal{A}$ と $\mathcal{B}$ を
$(\mathcal{C}, \mathcal{O})$ 上の微分次数付き代数の層とし、
$\mathcal{N}$ を右微分次数付き $\mathcal{B}$-加群とする。
$\mathcal{N}$ 上の、与えられた微分次数付き
$\mathcal{B}$-加群構造と両立する
$(\mathcal{A}, \mathcal{B})$-双加群構造と、
微分次数付き $\mathcal{O}$-代数の準同型
$$
\mathcal{A}
\longrightarrow
\SheafHom^{dg}_\mathcal{B}(\mathcal{N}, \mathcal{N})
$$
との間には一対一の対応がある。
\end{lemma}

\begin{proof}
省略する。
\end{proof}

\noindent
$(\mathcal{C}, \mathcal{O})$ を環付きサイトとする。
$\mathcal{A}$ と $\mathcal{B}$ を
$(\mathcal{C}, \mathcal{O})$ 上の微分次数付き代数の層とする。
$\mathcal{M}$ を右微分次数付き $\mathcal{A}$-加群とし、
$\mathcal{N}$ を微分次数付き
$(\mathcal{A}, \mathcal{B})$-双加群とする。このとき、
Section \ref{sdga-section-tensor-product-dg} で定義した微分次数付きテンソル積
$$
\mathcal{M} \otimes_\mathcal{A} \mathcal{N}
$$
は右微分次数付き $\mathcal{B}$-加群であり、その乗法写像は
Section \ref{sdga-section-graded-bimodules} のものと同じである。
この構成は関手および次数付き圏の関手
$$
\otimes_\mathcal{A} \mathcal{N} :
\textit{Mod}(\mathcal{A}, \text{d})
\longrightarrow
\textit{Mod}(\mathcal{B}, \text{d})
\quad\text{and}\quad
\otimes_\mathcal{A} \mathcal{N} :
\textit{Mod}^{dg}(\mathcal{A}, \text{d})
\longrightarrow
\textit{Mod}^{dg}(\mathcal{B}, \text{d})
$$
を定める。すなわち、$\mathcal{M}$ から $\mathcal{M}'$ への次数 $n$ の
準同型を、$\mathcal{M} \otimes_\mathcal{A} \mathcal{N}$ から
$\mathcal{M}' \otimes_\mathcal{A} \mathcal{N}$ への誘導された
次数 $n$ の写像に送る。

\medskip\noindent
$(\mathcal{C}, \mathcal{O})$ を環付きサイトとする。
$\mathcal{A}$ と $\mathcal{B}$ を
$(\mathcal{C}, \mathcal{O})$ 上の微分次数付き代数の層とする。
$\mathcal{N}$ を微分次数付き
$(\mathcal{A}, \mathcal{B})$-双加群とし、$\mathcal{L}$ を
右微分次数付き $\mathcal{B}$-加群とする。このとき、
Section \ref{sdga-section-internal-hom-dg} で定義した微分次数付き内部 Hom
$$
\SheafHom_\mathcal{B}^{dg}(\mathcal{N}, \mathcal{L})
$$
は右微分次数付き $\mathcal{A}$-加群であり、その右次数付き
$\mathcal{A}$-加群構造は
Section \ref{sdga-section-graded-bimodules} で定義したものである。
別の定義として、次の合成を乗法に用いることもできる。
$$
\SheafHom_\mathcal{B}^{dg}(\mathcal{N}, \mathcal{L})
\otimes_\mathcal{O} \mathcal{A}
\to
\SheafHom_\mathcal{B}^{dg}(\mathcal{N}, \mathcal{L})
\otimes_\mathcal{O} \SheafHom_\mathcal{B}^{dg}(\mathcal{N}, \mathcal{N})
\to
\SheafHom_\mathcal{B}^{dg}(\mathcal{N}, \mathcal{L})
$$
ここで、最初の矢印は
Lemma \ref{sdga-lemma-what-makes-a-bimodule-dg} から得られ、
二つ目の矢印は Section \ref{sdga-section-internal-hom-dg} の合成である。
これらの矢印はいずれも微分と両立するので、確かに微分次数付き
$\mathcal{A}$-加群が得られる。この構成は関手および
微分次数付き圏の関手
$$
\SheafHom_\mathcal{B}^{dg}(\mathcal{N}, -) :
\textit{Mod}(\mathcal{B}, \text{d})
\longrightarrow
\textit{Mod}(\mathcal{A})
\quad\text{and}\quad
\SheafHom_\mathcal{B}^{dg}(\mathcal{N}, -) :
\textit{Mod}^{dg}(\mathcal{B}, \text{d})
\longrightarrow
\textit{Mod}^{dg}(\mathcal{A}, \text{d})
$$
を定める。すなわち、$\mathcal{L}$ から $\mathcal{L}'$ への次数 $n$ の
準同型を、$\SheafHom_\mathcal{B}^{dg}(\mathcal{N}, \mathcal{L})$ から
$\SheafHom_\mathcal{B}^{dg}(\mathcal{N}, \mathcal{L}')$ への
誘導された次数 $n$ の写像に送る。

\begin{lemma}
\label{sdga-lemma-tensor-hom-adjunction-dg}
$(\mathcal{C}, \mathcal{O})$ を環付きサイトとする。
$\mathcal{A}$ と $\mathcal{B}$ を
$(\mathcal{C}, \mathcal{O})$ 上の微分次数付き代数の層とする。
$\mathcal{M}$ を右微分次数付き $\mathcal{A}$-加群とし、
$\mathcal{N}$ を微分次数付き
$(\mathcal{A}, \mathcal{B})$-双加群とし、$\mathcal{L}$ を
右微分次数付き $\mathcal{B}$-加群とする。上の規約のもとで
$$
\Hom_{\textit{Mod}^{dg}(\mathcal{B}, \text{d})}(
\mathcal{M} \otimes_\mathcal{A} \mathcal{N}, \mathcal{L}) =
\Hom_{\textit{Mod}^{dg}(\mathcal{A}, \text{d})}(
\mathcal{M}, \SheafHom_\mathcal{B}^{dg}(\mathcal{N}, \mathcal{L}))
$$
および
$$
\SheafHom_\mathcal{B}^{dg}(
\mathcal{M} \otimes_\mathcal{A} \mathcal{N}, \mathcal{L}) =
\SheafHom_\mathcal{A}^{dg}(
\mathcal{M}, \SheafHom_\mathcal{B}^{dg}(\mathcal{N}, \mathcal{L}))
$$
が成り立ち、これらは $\mathcal{M}$、$\mathcal{N}$、$\mathcal{L}$ に
関して関手的である。
\end{lemma}

\begin{proof}
省略する。ヒント：次数付きの水準では、
Lemma \ref{sdga-lemma-tensor-hom-adjunction-gr} でこれを既に示した。
したがって、これらの同型が微分と両立することを確認すれば十分であり、
それは局所切断の水準での計算により確かめられる。
\end{proof}

\noindent
$(\mathcal{C}, \mathcal{O})$ を環付きサイトとする。
$\mathcal{A}$ と $\mathcal{B}$ を
$(\mathcal{C}, \mathcal{O})$ 上の微分次数付き代数の層とする。
上の特別な場合として、微分次数付き $\mathcal{O}$-代数の準同型
$\varphi : \mathcal{A} \to \mathcal{B}$ が与えられたとする。
このとき、関手および微分次数付き圏の関手
$$
\otimes_{\mathcal{A}, \varphi} \mathcal{B} :
\textit{Mod}(\mathcal{A}, \text{d})
\longrightarrow
\textit{Mod}(\mathcal{B}, \text{d})
\quad\text{and}\quad
\otimes_{\mathcal{A}, \varphi} \mathcal{B} :
\textit{Mod}^{dg}(\mathcal{A}, \text{d})
\longrightarrow
\textit{Mod}^{dg}(\mathcal{B}, \text{d})
$$
を得る。一方、スカラー制限関手
$$
res_\varphi :
\textit{Mod}(\mathcal{B}, \text{d})
\longrightarrow
\textit{Mod}(\mathcal{A}, \text{d})
\quad\text{and}\quad
res_\varphi :
\textit{Mod}^{dg}(\mathcal{B}, \text{d})
\longrightarrow
\textit{Mod}^{dg}(\mathcal{A}, \text{d})
$$
がある。上の補題を用いると、これらの関手が互いに随伴であることが
分かる（スカラー制限と底変換における通常の場合と同様である）。
実際、$\mathcal{B}$ を微分次数付き
$(\mathcal{A}, \mathcal{B})$-双加群とみなしたものを
${}_\mathcal{A}\mathcal{B}_\mathcal{B}$ と書く。
すると、任意の右微分次数付き $\mathcal{B}$-加群 $\mathcal{L}$ に対して
$$
\SheafHom_\mathcal{B}^{dg}({}_\mathcal{A}\mathcal{B}_\mathcal{B}, \mathcal{L})
= res_\varphi(\mathcal{L})
$$
が右微分次数付き $\mathcal{A}$-加群として成り立つ。したがって
Lemma \ref{sdga-lemma-tensor-hom-adjunction-gr} により、関手的同型
$$
\Hom_{\textit{Mod}^{dg}(\mathcal{B}, \text{d})}(
\mathcal{M} \otimes_{\mathcal{A}, \varphi} \mathcal{B}, \mathcal{L}) =
\Hom_{\textit{Mod}^{dg}(\mathcal{A}, \text{d})}(
\mathcal{M}, res_\varphi(\mathcal{L}))
$$
を得る。文脈から $\varphi$ が明らかな場合、この式では通常
$\varphi$ への依存を省略する。同様にして、次数付き
$\mathcal{O}$-加群の等式
$$
\SheafHom^{dg}_\mathcal{B}(
\mathcal{M} \otimes_\mathcal{A} \mathcal{B}, \mathcal{L}) =
\SheafHom_\mathcal{A}^{dg}(\mathcal{M}, \mathcal{L})
$$
を得る。









\section{微分次数付き加群の層に対する引き戻しと順像}
\label{sdga-section-functoriality-dg}

\noindent
$(f, f^\sharp) : (\Sh(\mathcal{C}), \mathcal{O}_\mathcal{C})
\to (\Sh(\mathcal{D}), \mathcal{O}_\mathcal{D})$
を環付きトポスの射とする。$\mathcal{A}$ を微分次数付き
$\mathcal{O}_\mathcal{C}$-代数とし、$\mathcal{B}$ を微分次数付き
$\mathcal{O}_\mathcal{D}$-代数とする。微分次数付き
$f^{-1}\mathcal{O}_\mathcal{D}$-代数の準同型
$$
\varphi : f^{-1}\mathcal{B} \to \mathcal{A}
$$
が与えられているとする。スカラー制限とスカラー拡大の随伴により、
これは微分次数付き $\mathcal{O}_\mathcal{C}$-代数の準同型
$\varphi : f^*\mathcal{B} \to \mathcal{A}$ と同じものである。
同値な言い方をすれば、$\varphi$ は微分次数付き
$\mathcal{O}_\mathcal{D}$-代数の準同型
$$
\varphi : \mathcal{B} \to f_*\mathcal{A}
$$
とみなせる。Remark \ref{sdga-remark-functoriality-dga} を参照せよ。

\medskip\noindent
関手
$$
f_* :
\textit{Mod}(\mathcal{A}, \text{d})
\longrightarrow
\textit{Mod}(\mathcal{B}, \text{d})
$$
を定義しよう。微分次数付き $\mathcal{A}$-加群 $\mathcal{M}$ に対し、
$f_*\mathcal{M}$ を Section \ref{sdga-section-functoriality-graded} で
構成した次数付き $\mathcal{B}$-加群で、微分が写像
$f_*d : f_*\mathcal{M}^n \to f_*\mathcal{M}^{n + 1}$ により
与えられるものと定義する。この構成は明らかに $\mathcal{M}$ に関して
関手的であり、所望の関手を得る。

\medskip\noindent
関手
$$
f^* :
\textit{Mod}(\mathcal{B}, \text{d})
\longrightarrow
\textit{Mod}(\mathcal{A}, \text{d})
$$
を定義しよう。微分次数付き $\mathcal{B}$-加群 $\mathcal{N}$ に対し、
$f^*\mathcal{N}$ を Section \ref{sdga-section-functoriality-graded} で
構成した次数付き $\mathcal{A}$-加群と定義する。ここで
$$
f^*\mathcal{N} = f^{-1}\mathcal{N} \otimes_{f^{-1}\mathcal{B}} \mathcal{A}
$$
であった。$f^{-1}\mathcal{N}$ は微分
$f^{-1}\text{d} : f^{-1}\mathcal{N}^n \to f^{-1}\mathcal{N}^{n + 1}$
を備えるので、このテンソル積は
Section \ref{sdga-section-dg-bimodules} で論じたテンソル積の一例とみなせ、
そこから微分が得られる。この構成は明らかに $\mathcal{N}$ に関して
関手的であり、関手 $f^*$ を得る。

\medskip\noindent
関手 $f_*$ と $f^*$ は直ちに微分次数付き圏の関手
$$
f_* :
\textit{Mod}^{dg}(\mathcal{A}, \text{d})
\longrightarrow
\textit{Mod}^{dg}(\mathcal{B}, \text{d})
\quad\text{and}\quad
f^* :
\textit{Mod}^{dg}(\mathcal{B}, \text{d})
\longrightarrow
\textit{Mod}^{dg}(\mathcal{A}, \text{d})
$$
に拡張される。これらは台となる対象には同じ作用をし、次数 $n$ の
斉次射に対しては上の構成の関手性によって定義される。

\begin{lemma}
\label{sdga-lemma-adjunction-push-pull-dg}
上の状況では
$$
\Hom_{\textit{Mod}^{dg}(\mathcal{B}, \text{d})}(
\mathcal{N}, f_*\mathcal{M}) =
\Hom_{\textit{Mod}^{dg}(\mathcal{A}, \text{d})}(
f^*\mathcal{N}, \mathcal{M})
$$
が成り立つ。
\end{lemma}

\begin{proof}
省略する。ヒント：Lemma \ref{sdga-lemma-adjunction-push-pull-gr} により、
台となる次数付き圏についてはこれは成立する。計算により、
これらの同型が微分と両立することが分かる。
\end{proof}
















\section{局所化と微分次数付き加群の層}
\label{sdga-section-localize-dg}

\noindent
$(\mathcal{C}, \mathcal{O})$ を環付きサイトとする。
$U \in \Ob(\mathcal{C})$ とし、対応する局所化射
$$
j :
(\Sh(\mathcal{C}/U), \mathcal{O}_U)
\longrightarrow
(\Sh(\mathcal{C}), \mathcal{O})
$$
を $j$ と記す
(Modules on Sites, Section \ref{sites-modules-section-localize})。
以下では次の事実を用いる。$\mathcal{O}_U$-加群
$\mathcal{M}_i$、$i = 1, 2$ と $\mathcal{O}$-加群
$\mathcal{A}$ に対して、標準写像
$$
j_! :
\Hom_{\mathcal{O}_U}(
\mathcal{M}_1 \otimes_{\mathcal{O}_U} \mathcal{A}|_U, \mathcal{M}_2)
\longrightarrow
\Hom_\mathcal{O}(
j_!\mathcal{M}_1 \otimes_\mathcal{O} \mathcal{A}, j_!\mathcal{M}_2)
$$
がある。実際、Modules on Sites,
Lemma \ref{sites-modules-lemma-j-shriek-and-tensor} により
$j_!(\mathcal{M}_1 \otimes_{\mathcal{O}_U} \mathcal{A}|_U) =
j_!\mathcal{M}_1 \otimes_\mathcal{O} \mathcal{A}$ である。

\medskip\noindent
$\mathcal{A}$ を微分次数付き $\mathcal{O}$-代数とする。
$\mathcal{A}$ の $\mathcal{C}/U$ への制限を
$\mathcal{A}_U$ と記す。すなわち
$\mathcal{A}_U = j^*\mathcal{A} = j^{-1}\mathcal{A}$ である。
Section \ref{sdga-section-functoriality-dg} では随伴関手
$$
j_* :
\textit{Mod}^{dg}(\mathcal{A}_U, \text{d})
\longrightarrow
\textit{Mod}^{dg}(\mathcal{A}, \text{d})
\quad\text{and}\quad
j^* :
\textit{Mod}^{dg}(\mathcal{A}, \text{d})
\longrightarrow
\textit{Mod}^{dg}(\mathcal{A}_U, \text{d})
$$
を構成した。ここで $j^*$ は $j_*$ の左随伴である。さらに、
$j_*$ の右随伴となる完全関手
$$
j_! :
\textit{Mod}^{dg}(\mathcal{A}_U, \text{d})
\longrightarrow
\textit{Mod}^{dg}(\mathcal{A}, \text{d})
$$
があると主張する。実際、微分次数付き
$\mathcal{A}_U$-加群 $\mathcal{M}$ に対し、
$j_!\mathcal{M}$ を Section \ref{sdga-section-localize-graded} で構成した
次数付き $\mathcal{A}$-加群で、微分が
$j_!\text{d} : j_!\mathcal{M}^n \to j_!\mathcal{M}^{n + 1}$
であるものと定義する。微分次数付き $\mathcal{A}_U$-加群の
次数 $n$ の斉次写像
$f : \mathcal{M} \to \mathcal{M}'$ が与えられると、
微分次数付き $\mathcal{A}$-加群の次数 $n$ の斉次写像
$j_!f : j_!\mathcal{M} \to j_!\mathcal{M}'$ を得る。
この構成が微分と両立することの直接的な確認は省略する。
こうして所望の関手を得る。

\begin{lemma}
\label{sdga-lemma-extension-by-zero-dg}
上の状況では
$$
\Hom_{\textit{Mod}^{dg}(\mathcal{A}, \text{d})}(
j_!\mathcal{M}, \mathcal{N}) =
\Hom_{\textit{Mod}^{dg}(\mathcal{A}_U, \text{d})}(
\mathcal{M}, j^*\mathcal{N})
$$
が成り立つ。
\end{lemma}

\begin{proof}
省略する。ヒント：Lemma \ref{sdga-lemma-extension-by-zero-graded} により、
この補題は次数付きの水準で成立することを既に見た。したがって、
得られる同型が微分と両立することだけを確認すればよい。
\end{proof}

\begin{lemma}
\label{sdga-lemma-tensor-with-extension-by-zero-dg}
上の状況で、$\mathcal{M}$ を右微分次数付き
$\mathcal{A}_U$-加群とし、$\mathcal{N}$ を左微分次数付き
$\mathcal{A}$-加群とする。このとき
$$
j_!\mathcal{M} \otimes_\mathcal{A} \mathcal{N} =
j_!(\mathcal{M} \otimes_{\mathcal{A}_U} \mathcal{N}|_U)
$$
が $\mathcal{O}$-加群の複体として成り立ち、
$\mathcal{M}$ と $\mathcal{N}$ に関して関手的である。
\end{lemma}

\begin{proof}
次数付き加群としては、
Lemma \ref{sdga-lemma-tensor-with-extension-by-zero} から従う。
この同型が微分と両立することの確認は省略する。
\end{proof}








\section{微分次数付き加群の層上のシフト関手}
\label{sdga-section-shift-dg}

\noindent
$(\mathcal{C}, \mathcal{O})$ を環付きサイトとする。
$\mathcal{A}$ を $(\mathcal{C}, \mathcal{O})$ 上の
微分次数付き代数の層とし、$\mathcal{M}$ を微分次数付き
$\mathcal{A}$-加群とする。$k \in \mathbf{Z}$ とする。
$\mathcal{M}$ の{\it 第 $k$ シフト}を $\mathcal{M}[k]$ と記し、
次のように定義する。
\begin{enumerate}
\item 次数付き $\mathcal{A}$-加群としての $\mathcal{M}[k]$ は
Section \ref{sdga-section-shift} で定義したものとする。
\item 微分
$d_{\mathcal{M}[k]} : (\mathcal{M}[k])^n \to (\mathcal{M}[k])^{n + 1}$
を
$(-1)^k\text{d}_\mathcal{M} : \mathcal{M}^{n + k} \to \mathcal{M}^{n + k + 1}$
と定義する。
\end{enumerate}
次数 $n$ の斉次な $\mathcal{A}$-加群準同型
$f : \mathcal{L} \to \mathcal{M}$ に対し、
$f[k] : \mathcal{L}[k] \to \mathcal{M}[k]$ を $f$ と同じ成分写像で
与える。このとき $f[k]$ は次数 $n$ の斉次な
$\mathcal{A}$-加群写像である。これにより写像
$$
\Hom_{\textit{Mod}^{dg}(\mathcal{A}, \text{d})}(\mathcal{L}, \mathcal{M})
\longrightarrow
\Hom_{\textit{Mod}^{dg}(\mathcal{A}, \text{d})}
(\mathcal{L}[k], \mathcal{M}[k])
$$
を得る。この写像は微分と両立する（$\mathcal{L}$ と
$\mathcal{M}$ の微分の符号が同じだけ変化することから従う）。
これらの選択は Differential Graded Algebra,
Definition \ref{dga-definition-shift} における選択と両立する。
微分次数付き圏の関手
$$
[k] :
\textit{Mod}^{dg}(\mathcal{A}, \text{d})
\longrightarrow
\textit{Mod}^{dg}(\mathcal{A}, \text{d})
$$
を定義したこと、および $[k + l] = [k] \circ [l]$ は明らかである。

\medskip\noindent
台となる次数付き加群上で Section \ref{sdga-section-shift} に定義した同型
$$
\Hom_{\textit{Mod}^{dg}(\mathcal{A}, \text{d})}(\mathcal{L}, \mathcal{M}[k])
=
\Hom_{\textit{Mod}^{dg}(\mathcal{A}, \text{d})}(\mathcal{L}, \mathcal{M})[k]
$$
は微分と両立すると主張する。これを見るため、
$f : \mathcal{L} \to \mathcal{M}[k]$ を次数 $n$ の斉次な右
$\mathcal{A}$-加群写像とする。これは左辺の次数 $n$ の元である。
$f$ と{\bf 同じ}成分写像をもつ、次数 $n + k$ の斉次な
$\mathcal{A}$-加群写像
$f' : \mathcal{L} \to \mathcal{M}$ を $f'$ と記す。
規約により、これは右辺の次数 $n$ の対応する元である。
左辺の微分の定義から
$$
\text{d}_{LHS}(f) =
\text{d}_{\mathcal{M}[k]} \circ f - (-1)^n f \circ \text{d}_\mathcal{L} =
(-1)^k\text{d}_\mathcal{M} \circ f - (-1)^n f \circ \text{d}_\mathcal{L}
$$
を得る。一方、右辺の微分については
$$
\text{d}_{RHS}(f') =
(-1)^k\left(
\text{d}_\mathcal{M} \circ f' - (-1)^{n + k} f' \circ \text{d}_\mathcal{L}
\right) =
(-1)^k\text{d}_\mathcal{M} \circ f' - (-1)^n f' \circ \text{d}_\mathcal{L}
$$
を得る。これらの写像は同じ成分写像をもつので、証明は完了する。










\section{ホモトピー圏}
\label{sdga-section-homotopy}

\noindent
本節は Differential Graded Algebra,
Section \ref{dga-section-homotopy} に対応する。

\begin{definition}
\label{sdga-definition-homotopy}
$(\mathcal{C}, \mathcal{O})$ を環付きサイトとする。
$\mathcal{A}$ を $(\mathcal{C}, \mathcal{O})$ 上の
微分次数付き代数の層とする。
$f, g : \mathcal{M} \to \mathcal{N}$ を微分次数付き
$\mathcal{A}$-加群の準同型とする。
{\it $f$ と $g$ の間のホモトピー}とは、次数 $-1$ の斉次な
次数付き $\mathcal{A}$-加群写像
$h : \mathcal{M} \to \mathcal{N}$ であって
$$
f - g = \text{d}_\mathcal{N} \circ h + h \circ \text{d}_\mathcal{M}
$$
を満たすものをいう。ホモトピーが存在するとき、$f$ と $g$ は
{\it ホモトピックである}という。
\end{definition}

\noindent
定義の状況で、写像
$a : \mathcal{K} \to \mathcal{M}$ と $c : \mathcal{N} \to \mathcal{L}$
があると
$$
\begin{matrix}
h\text{ はホモトピーである} \\
\text{すなわち }f\text{ と } g\text{ の間のものである}
\end{matrix}
\quad
\Rightarrow
\quad
\begin{matrix}
c \circ h \circ a\text{ はホモトピーである}\\
\text{すなわち }
c \circ f \circ a\text{ と } c\circ g \circ a\text{ の間のものである}
\end{matrix}
$$
であることが分かる。したがって
$\textit{Mod}(\mathcal{A}, \text{d})$ における射のホモトピー類の
合成を定義できる。

\begin{definition}
\label{sdga-definition-complexes-notation}
$(\mathcal{C}, \mathcal{O})$ を環付きサイトとする。
$\mathcal{A}$ を $(\mathcal{C}, \mathcal{O})$ 上の
微分次数付き代数の層とする。{\it ホモトピー圏}とは、
$K(\textit{Mod}(\mathcal{A}, \text{d}))$ と記される次の圏である。
その対象は $\textit{Mod}(\mathcal{A}, \text{d})$ の対象であり、
その射は微分次数付き $\mathcal{A}$-加群準同型のホモトピー類である。
\end{definition}

\noindent
記法 $K(\textit{Mod}(\mathcal{A}, \text{d}))$ は標準的ではないが、
少なくとも Stacks project の他の箇所における $K(-)$ の用法と整合する。

\medskip\noindent
Differential Graded Algebra, Definition
\ref{dga-definition-homotopy-category-of-dga-category} では、
微分次数付き圏 $\mathcal{S}$ の複体の圏
$\text{Comp}(\mathcal{S})$ とホモトピー圏
$K(\mathcal{S})$ の意味を定義した。これを微分次数付き圏
$\textit{Mod}^{dg}(\mathcal{A}, \text{d})$ に適用すると
$$
\textit{Mod}(\mathcal{A}, \text{d}) =
\text{Comp}(\textit{Mod}^{dg}(\mathcal{A}, \text{d}))
$$
を得る（Section \ref{sdga-section-dgm-dg-cat} の議論を参照せよ）。
さらに
$$
K(\textit{Mod}(\mathcal{A}, \text{d})) =
K(\textit{Mod}^{dg}(\mathcal{A}, \text{d}))
$$
を得る。この最後の等式を見るには、$h$ が
Definition \ref{sdga-definition-homotopy} のとおりであるとき
$$
\text{d}_{\Hom_{\textit{Mod}^{dg}(\mathcal{A}, \text{d})}
(\mathcal{M}, \mathcal{N})}(h) =
\text{d}_\mathcal{N} \circ h + h \circ \text{d}_\mathcal{M}
$$
という等式が成り立つことに注意すればよい。詳細は省略する。

\begin{lemma}
\label{sdga-lemma-homotopy-direct-sums}
$(\mathcal{C}, \mathcal{O})$ を環付きサイトとする。
$\mathcal{A}$ を $(\mathcal{C}, \mathcal{O})$ 上の
微分次数付き代数の層とする。
ホモトピー圏 $K(\textit{Mod}(\mathcal{A}, \text{d}))$ は
直和と直積をもつ。
\end{lemma}

\begin{proof}
省略する。ヒント：Lemma \ref{sdga-lemma-dgm-abelian} の直和と直積を
そのまま用いればよい。これでよいのは、これらの関手が次数付き
$\mathcal{A}$-加群の圏への忘却関手と可換することを既に見たこと、
およびアーベル群の族の圏上で $\prod$ と $\bigoplus$ が
完全関手であることによる。
\end{proof}












\section{錐と三角形}
\label{sdga-section-conclude-triangulated}

\noindent
本節では Differential Graded Algebra,
Section \ref{dga-section-review} の内容を用いて、微分次数付き
$\mathcal{A}$-加群の圏のホモトピー圏が三角圏であることを示す。

\begin{lemma}
\label{sdga-lemma-axioms-AB}
$(\mathcal{C}, \mathcal{O})$ を環付きサイトとする。
$\mathcal{A}$ を $(\mathcal{C}, \mathcal{O})$ 上の
微分次数付き代数の層とする。微分次数付き圏
$\textit{Mod}^{dg}(\mathcal{A}, \text{d})$ は
Differential Graded Algebra,
Section \ref{dga-section-review} の公理 (A) と (B) を満たす。
\end{lemma}

\begin{proof}
微分次数付き $\mathcal{A}$-加群 $\mathcal{M}$ と
$\mathcal{N}$ が与えられたとする。明らかな方法で定義される
微分次数付き $\mathcal{A}$-加群
$\mathcal{M} \oplus \mathcal{N}$ を考える。このとき余入射
$i : \mathcal{M} \to \mathcal{M} \oplus \mathcal{N}$ と
$j : \mathcal{N} \to \mathcal{M} \oplus \mathcal{N}$、および射影
$p : \mathcal{M} \oplus \mathcal{N} \to \mathcal{N}$ と
$q : \mathcal{M} \oplus \mathcal{N} \to \mathcal{M}$ は
微分次数付き $\mathcal{A}$-加群の射である。したがって
$i, j, p, q$ は次数 $0$ の斉次かつ閉な射、すなわち
$\text{d}(i) = 0$ などを満たす。ゆえに、この直和は
Differential Graded Algebra,
Definition \ref{dga-definition-dg-direct-sum} の意味での
微分次数付き和である。これで公理 (A) が証明された。

\medskip\noindent
公理 (B) は Section \ref{sdga-section-shift-dg} で示した。
\end{proof}

\noindent
$(\mathcal{C}, \mathcal{O})$ を環付きサイトとする。
$\mathcal{A}$ を $(\mathcal{C}, \mathcal{O})$ 上の
微分次数付き代数の層とする。
$\textit{Mod}(\mathcal{A}, \text{d})$ の列
$$
0 \to \mathcal{K} \to \mathcal{L} \to \mathcal{N} \to 0
$$
が{\it 許容短完全列}と呼ばれるのは
(Differential Graded Algebra,
Section \ref{dga-section-review})、それが
$\textit{Mod}(\mathcal{A})$ で分裂するときであった。
言い換えれば、次数付き $\mathcal{A}$-加群の列として分裂するときである。
次数付き $\mathcal{A}$-加群としての分裂を
$s : \mathcal{N} \to \mathcal{L}$ と
$\pi : \mathcal{L} \to \mathcal{K}$ と記す。
Lemma \ref{sdga-lemma-axioms-AB} と Differential Graded Algebra,
Lemma \ref{dga-lemma-get-triangle} を合わせると、三角形
$$
\mathcal{K} \to \mathcal{L} \to \mathcal{N} \to \mathcal{K}[1]
$$
を得る。ここで Differential Graded Algebra,
Lemma \ref{dga-lemma-get-triangle} の証明における矢印
$\mathcal{N} \to \mathcal{K}[1]$ は
$$
\delta = \pi \circ
\text{d}_{\Hom_{\textit{Mod}^{dg}(\mathcal{A}, \text{d})}
(\mathcal{L}, \mathcal{M})}(s) =
\pi \circ \text{d}_\mathcal{L} \circ s -
\pi \circ s \circ \text{d}_\mathcal{N} =
\pi \circ \text{d}_\mathcal{L} \circ s
$$
として構成される。このひどく込み入った記法は御容赦いただきたい。
いずれにせよ、この状況では Differential Graded Algebra,
Lemma \ref{dga-lemma-get-triangle} で構成された境界写像 $\delta$ は、
台となる $\mathcal{O}$-加群の複体上で、Stacks project 全体を通じて
項ごとに分裂する複体の短完全列に用いられる通常の境界写像と一致する。
Derived Categories, Definition \ref{derived-definition-split-ses} を参照せよ。

\begin{definition}
\label{sdga-definition-cone}
$(\mathcal{C}, \mathcal{O})$ を環付きサイトとする。
$\mathcal{A}$ を $(\mathcal{C}, \mathcal{O})$ 上の
微分次数付き代数の層とする。
$f : \mathcal{K} \to \mathcal{L}$ を微分次数付き
$\mathcal{A}$-加群の準同型とする。$f$ の{\it 錐}とは、
次のように定義される微分次数付き $\mathcal{A}$-加群 $C(f)$ である。
\begin{enumerate}
\item 台となる $\mathcal{O}$-加群の複体は、対応する
$\mathcal{A}$-加群の複体の写像
$f : \mathcal{K}^\bullet \to \mathcal{L}^\bullet$ の錐である。
すなわち $C(f)^n = \mathcal{L}^n \oplus \mathcal{K}^{n + 1}$ であり、
微分は
$$
d_{C(f)} =
\left(
\begin{matrix}
\text{d}_\mathcal{L} & f \\
0 & -\text{d}_\mathcal{K}
\end{matrix}
\right)
$$
である。
\item 乗法写像
$$
C(f)^n \times \mathcal{A}^m \to C(f)^{n + m}
$$
は、乗法写像
$\mathcal{L}^n \times \mathcal{A}^m \to \mathcal{L}^{n + m}$
と乗法写像
$\mathcal{K}^{n + 1} \times \mathcal{A}^m \to \mathcal{K}^{n + 1 + m}$
の直和である。
\end{enumerate}
これは明らかな写像から誘導される微分次数付き
$\mathcal{A}$-加群の標準準同型
$i : \mathcal{L} \to C(f)$ と $p : C(f) \to \mathcal{K}[1]$
を備える。
\end{definition}

\noindent
定義の状況では、列
$$
0 \to \mathcal{L} \to C(f) \to \mathcal{K}[1] \to 0
$$
は許容短完全列であることに注意せよ。

\begin{lemma}
\label{sdga-lemma-axiom-C}
$(\mathcal{C}, \mathcal{O})$ を環付きサイトとする。
$\mathcal{A}$ を $(\mathcal{C}, \mathcal{O})$ 上の
微分次数付き代数の層とする。微分次数付き圏
$\textit{Mod}^{dg}(\mathcal{A}, \text{d})$ は
Differential Graded Algebra,
Situation \ref{dga-situation-ABC} で定式化された公理 (C) を満たす。
\end{lemma}

\begin{proof}
$f : \mathcal{K} \to \mathcal{L}$ を微分次数付き
$\mathcal{A}$-加群の準同型とする。上で許容短完全列
$$
0 \to \mathcal{L} \to C(f) \to \mathcal{K}[1] \to 0
$$
を得た。証明を終えるには、これに付随する境界写像
$$
\delta : \mathcal{K}[1] \to \mathcal{L}[1]
$$
（上の議論を参照）が $f[1]$ に等しいことを示さなければならない。
切断 $s : \mathcal{K}[1] \to C(f)$ として、次数 $n$ では埋め込み
$\mathcal{K}^{n + 1} \to C(f)^n$ を用いる。写像 $\pi$ は次数 $n$ で
射影 $C(f)^n \to \mathcal{L}^n$ により与えられる。したがって最後に
$$
\delta = \pi \circ \text{d}_{C(f)} \circ s
$$
を計算しなければならない（上の議論を参照）。
行列記法ではこれは
$$
\left(
\begin{matrix}
1 & 0
\end{matrix}
\right)
\left(
\begin{matrix}
\text{d}_\mathcal{L} & f \\
0 & -\text{d}_\mathcal{K}
\end{matrix}
\right)
\left(
\begin{matrix}
0 \\
1
\end{matrix}
\right) = f
$$
に等しく、望む結果を得る。
\end{proof}

\noindent
以上で、Differential Graded Algebra,
Section \ref{dga-section-review} で証明されたすべての補題が
微分次数付き圏 $\textit{Mod}^{dg}(\mathcal{A}, \text{d})$
に対して有効であることが分かった。特に次を得る。

\begin{proposition}
\label{sdga-proposition-homotopy-category-triangulated}
$(\mathcal{C}, \mathcal{O})$ を環付きサイトとする。
$\mathcal{A}$ を $(\mathcal{C}, \mathcal{O})$ 上の
微分次数付き代数の層とする。
ホモトピー圏 $K(\textit{Mod}(\mathcal{A}, \text{d}))$ は三角圏であり、
\begin{enumerate}
\item シフト関手は Section \ref{sdga-section-shift-dg} で構成したものである。
\item 識別三角形は $K(\textit{Mod}(\mathcal{A}, \text{d}))$ の
三角形のうち、三角形として
$$
\mathcal{K} \to \mathcal{L} \to \mathcal{N}
\xrightarrow{\delta} \mathcal{K}[1],\quad\quad
\delta = \pi \circ \text{d}_\mathcal{L} \circ s
$$
に同型なものである。ここでこの三角形は、上で
$\textit{Mod}(\mathcal{A}, \text{d})$ の許容短完全列
$0 \to \mathcal{K} \to \mathcal{L} \to \mathcal{N} \to 0$
から構成したものである。
\end{enumerate}
\end{proposition}

\begin{proof}
Section \ref{sdga-section-homotopy} にあるように
$K(\textit{Mod}(\mathcal{A}, \text{d})) =
K(\textit{Mod}^{dg}(\mathcal{A}, \text{d}))$ であった。
これを踏まえると、この命題は Lemmas \ref{sdga-lemma-axioms-AB} と
\ref{sdga-lemma-axiom-C}、および Differential Graded Algebra,
Proposition \ref{dga-proposition-ABC-homotopy-category-triangulated}
から従う。
\end{proof}

\begin{remark}
\label{sdga-remark-cone-identity}
$(\mathcal{C}, \mathcal{O})$ を環付きサイトとする。
$\mathcal{A}$ を $(\mathcal{C}, \mathcal{O})$ 上の
微分次数付き代数の層とする。
$C = C(\text{id}_\mathcal{A})$ を、微分次数付き
$\mathcal{A}$-加群の写像とみなした恒等写像
$\mathcal{A} \to \mathcal{A}$ の錐とする。このとき
$$
\Hom_{\textit{Mod}(\mathcal{A}, \text{d})}(C, \mathcal{M}) =
\{(x, y) \in
\Gamma(\mathcal{C}, \mathcal{M}^0) \times
\Gamma(\mathcal{C}, \mathcal{M}^{-1}) \mid
x = \text{d}(y)\}
$$
である。左から右への写像は $f$ を対 $(x, y)$ に送る。
ここで $x$ は
$C^{-1} = \mathcal{A}^{-1} \oplus \mathcal{A}^0$
の大域切断 $(0, 1)$ の像であり、$y$ は
$C^0 = \mathcal{A}^0 \oplus \mathcal{A}^1$
の大域切断 $(1, 0)$ の像である。
\end{remark}

\begin{lemma}
\label{sdga-lemma-dgm-grothendieck-abelian}
$(\mathcal{C}, \mathcal{O})$ を環付きサイトとする。
$(\mathcal{A}, \text{d})$ を微分次数付き $\mathcal{O}$-代数とする。
圏 $\textit{Mod}(\mathcal{A}, \text{d})$ は
Grothendieck アーベル圏である。
\end{lemma}

\begin{proof}
Lemma \ref{sdga-lemma-dgm-abelian} と Grothendieck アーベル圏の定義
(Injectives, Definition \ref{injectives-definition-grothendieck-conditions})
により、$\textit{Mod}(\mathcal{A}, \text{d})$ が生成対象をもつことを
示せば十分である。$\mathcal{C}$ の各対象 $U$ に対し、
Remark \ref{sdga-remark-cone-identity} と同様に恒等写像
$\mathcal{A}_U \to \mathcal{A}_U$ の錐を $C_U$ と記す。
$$
\mathcal{G} = \bigoplus\nolimits_{k, U} j_{U!}C_U[k]
$$
は生成対象であると主張する。ここで和は $\mathcal{C}$ のすべての対象
$U$ と $k \in \mathbf{Z}$ にわたる。実際、微分次数付き
$\mathcal{A}$-加群 $\mathcal{M}$ が与えられ、
$\mathcal{G}$ から $\mathcal{M}$ への非零写像が一つもなければ、
すべての $k$ と $U$ に対して
\begin{align*}
\Hom_{\textit{Mod}(\mathcal{A})}(j_{U!}C_U[k], \mathcal{M}) \\
& =
\Hom_{\textit{Mod}(\mathcal{A}_U)}(C_U[k], \mathcal{M}|_U) \\
& =
\{(x, y) \in \mathcal{M}^{-k}(U) \times \mathcal{M}^{-k - 1}(U) \mid
x = \text{d}(y)\}
\end{align*}
は零である。したがって $\mathcal{M}$ は零である。
\end{proof}













\section{平坦分解}
\label{sdga-section-P-resolutions}

\noindent
本節は Differential Graded Algebra,
Section \ref{dga-section-P-resolutions} に対応する。

\medskip\noindent
$(\mathcal{C}, \mathcal{O})$ を環付きサイトとする。
$\mathcal{A}$ を $(\mathcal{C}, \mathcal{O})$ 上の
微分次数付き代数の層とする。右微分次数付き
$\mathcal{A}$-加群 $\mathcal{P}$ が次を満たすとき、
$\mathcal{P}$ は{\it 良い}と呼ぶ。
\begin{enumerate}
\item 関手
$\mathcal{N} \mapsto \mathcal{P} \otimes_\mathcal{A} \mathcal{N}$
は次数付き左 $\mathcal{A}$-加群の圏上で完全である。
\item $\mathcal{N}$ が非輪状な微分次数付き左
$\mathcal{A}$-加群ならば、
$\mathcal{P} \otimes_\mathcal{A} \mathcal{N}$ は非輪状である。
\item 環付きトポスの任意の射
$(f, f^\sharp) : (\Sh(\mathcal{C}'), \mathcal{O}')
\to (\Sh(\mathcal{C}), \mathcal{O})$、任意の微分次数付き
$\mathcal{O}'$-代数 $\mathcal{A}'$、および微分次数付き
$f^{-1}\mathcal{O}_\mathcal{D}$-代数の任意の写像
$\varphi : f^{-1}\mathcal{A} \to \mathcal{A}'$ に対して、
微分次数付き $\mathcal{A}'$-加群とみなした引き戻し
$f^*\mathcal{P}$ (Section \ref{sdga-section-functoriality-dg}) は
性質 (1) と (2) を満たす。
\end{enumerate}
第1の条件は、$\mathcal{P}$ が右次数付き
$\mathcal{A}$-加群として平坦であることを意味する。第2の条件は
$\mathcal{P}$ が Spaltenstein の意味で K-平坦であることを意味し
(Cohomology on Sites,
Section \ref{sites-cohomology-section-flat} を参照せよ)、
第3の条件はこれが任意の底変換後にも成り立つことを意味する。

\medskip\noindent
意外かもしれないが、良い加群は多数存在する。

\begin{lemma}
\label{sdga-lemma-supply-good}
$(\mathcal{C}, \mathcal{O})$ を環付きサイトとする。
$\mathcal{A}$ を $(\mathcal{C}, \mathcal{O})$ 上の
微分次数付き代数の層とする。$U \in \Ob(\mathcal{C})$ とする。
このとき $j_!\mathcal{A}_U$ は良い微分次数付き
$\mathcal{A}$-加群である。
\end{lemma}

\begin{proof}
$\mathcal{N}$ を次数付き左 $\mathcal{A}$-加群とする。
Lemma \ref{sdga-lemma-tensor-with-extension-by-zero} により、次数付き加群として
$$
j_!\mathcal{A}_U \otimes_\mathcal{A} \mathcal{N} =
j_!(\mathcal{A}_U \otimes_{\mathcal{A}_U} \mathcal{N}|_U) =
j_!(\mathcal{N}_U)
$$
である。$U$ への制限と $j_!$ はともに完全なので、
これは条件 (1) を証明する。
Lemma \ref{sdga-lemma-tensor-with-extension-by-zero-dg} を用いれば、
同じ議論が (2) にも適用できる。

\medskip\noindent
環付きトポスの射
$(f, f^\sharp) : (\Sh(\mathcal{C}'), \mathcal{O}')
\to (\Sh(\mathcal{C}), \mathcal{O})$、微分次数付き
$\mathcal{O}'$-代数 $\mathcal{A}'$、および微分次数付き
$f^{-1}\mathcal{O}$-代数の写像
$\varphi : f^{-1}\mathcal{A} \to \mathcal{A}'$ を考える。
環付きトポス $(\Sh(\mathcal{C}'), \mathcal{O}')$ に
微分次数付き $\mathcal{O}'$-代数の層 $\mathcal{A}'$ を備えたものに対し
$$
f^*j_!\mathcal{A}_U =
f^{-1}j_!\mathcal{A}_U \otimes_{f^{-1}\mathcal{A}} \mathcal{A}'
$$
が (1) と (2) を満たすことを示さなければならない。
これを証明する際、
$(\Sh(\mathcal{C}), \mathcal{O})$ と
$(\Sh(\mathcal{C}'), \mathcal{O}')$ を同値な環付きトポスで
置き換えてよい。したがって Modules on Sites, Lemma
\ref{sites-modules-lemma-morphism-ringed-topoi-comes-from-morphism-ringed-sites}
により、$f$ は連続関手 $u : \mathcal{C} \to \mathcal{C}'$
によって与えられるサイトの射
$f : \mathcal{C} \to \mathcal{C}'$ から生じるとしてよい。
この場合 $U' = u(U)$ とおき、対応する局所化射を
$j' : \Sh(\mathcal{C}'/U') \to \Sh(\mathcal{C}')$ と記す。
環付きトポスの射の可換正方形
$$
\xymatrix{
(\Sh(\mathcal{C}'/U'), \mathcal{O}'_{U'})
\ar[rr]_{(j', (j')^\sharp)} \ar[d]_{(f', (f')^\sharp)} & &
(\Sh(\mathcal{C}'), \mathcal{O}')
\ar[d]^{(f, f^\sharp)} \\
(\Sh(\mathcal{C}/U), \mathcal{O}_U)
\ar[rr]^{(j, j^\sharp)} & &
(\Sh(\mathcal{C}), \mathcal{O}).
}
$$
を得て、$f'_*(j')^{-1} = j^{-1}f_*$ である。Modules on Sites,
Lemma \ref{sites-modules-lemma-localize-morphism-ringed-sites}
を参照せよ。随伴の一意性により
$f^{-1}j_! = j'_!(f')^{-1}$ を得る。したがって
\begin{align*}
f^*j_!\mathcal{A}_U
& =
f^{-1}j_!\mathcal{A}_U \otimes_{f^{-1}\mathcal{A}} \mathcal{A}' \\
& =
j'_!(f')^{-1}\mathcal{A}_U \otimes_{f^{-1}\mathcal{A}} \mathcal{A}' \\
& =
j'_!\left(
(f')^{-1}\mathcal{A}_U \otimes_{f^{-1}\mathcal{A}|_{U'}}
\mathcal{A}'|_{U'}\right) \\
& =
j'_!\mathcal{A}'_{U'}
\end{align*}
となる。第1の等式は $j_!\mathcal{A}_U$ を
$\mathcal{A}'$ 上の微分次数付き加群へ引き戻す定義である。
第2の等式は $f^{-1}j_! = j'_!(f')^{-1}$ による。
第3の等式は、環付きサイト $(\mathcal{C}', f^{-1}\mathcal{O})$、
微分次数付き代数の層 $f^{-1}\mathcal{A}$、
$\mathcal{C}'/U'$ 上の微分次数付き加群
$(f')^{-1}\mathcal{A}_U$、および $\mathcal{C}'$ 上の
微分次数付き加群 $\mathcal{A}'$ に
Lemma \ref{sdga-lemma-tensor-with-extension-by-zero-dg} を適用したものである。
第4の等式は、もちろん
$(f')^{-1}\mathcal{A}_U = f^{-1}\mathcal{A}|_{U'}$ だから成立する。
したがって引き戻しも同じ種類の加群であり、それについて (1) と (2) は
上で既に証明した。
\end{proof}

\begin{lemma}
\label{sdga-lemma-good-admissible-ses}
$(\mathcal{C}, \mathcal{O})$ を環付きサイトとする。
$\mathcal{A}$ を $(\mathcal{C}, \mathcal{O})$ 上の
微分次数付き代数の層とする。
$0 \to \mathcal{P} \to \mathcal{P}' \to \mathcal{P}'' \to 0$
を微分次数付き $\mathcal{A}$-加群の許容短完全列とする。
これら三つの加群のうち二つが良ければ、残りの一つも良い。
\end{lemma}

\begin{proof}
条件 (1) については、この列が次数付きの水準で直和であるため直ちに従う。
条件 (2) については、任意の微分次数付き左 $\mathcal{A}$-加群に対し
$$
0 \to 
\mathcal{P} \otimes_\mathcal{A} \mathcal{N} \to
\mathcal{P}' \otimes_\mathcal{A} \mathcal{N} \to
\mathcal{P}'' \otimes_\mathcal{A} \mathcal{N} \to 0
$$
が微分次数付き $\mathcal{O}$-加群の許容短完全列であることに注意する
（微分を忘れれば、テンソル積は単に次数付き加群の圏でとられる）。
したがって、そのうち二つが $\mathcal{O}$-加群の複体として
完全ならば、三つ目も完全である。最後に、同じ議論により、
環付きトポスの射
$(f, f^\sharp) : (\Sh(\mathcal{C}'), \mathcal{O}')
\to (\Sh(\mathcal{C}), \mathcal{O})$、微分次数付き
$\mathcal{O}'$-代数 $\mathcal{A}'$、および微分次数付き
$f^{-1}\mathcal{O}$-代数の写像
$\varphi : f^{-1}\mathcal{A} \to \mathcal{A}'$ が与えられれば
$$
0 \to f^*\mathcal{P} \to f^*\mathcal{P}' \to f^*\mathcal{P}'' \to 0
$$
は微分次数付き $\mathcal{A}'$-加群の許容短完全列となり、
上と同じ議論が適用される。
\end{proof}


\begin{lemma}
\label{sdga-lemma-good-direct-sum}
$(\mathcal{C}, \mathcal{O})$ を環付きサイトとする。
$\mathcal{A}$ を $(\mathcal{C}, \mathcal{O})$ 上の
微分次数付き代数の層とする。良い微分次数付き
$\mathcal{A}$-加群の任意の直和は良い。良い微分次数付き
$\mathcal{A}$-加群のフィルター余極限は良い。
\end{lemma}

\begin{proof}
省略する。ヒント：直和とフィルター余極限はテンソル積および
引き戻しと可換する。
\end{proof}

\begin{lemma}
\label{sdga-lemma-good-quotient}
$(\mathcal{C}, \mathcal{O})$ を環付きサイトとする。
$\mathcal{A}$ を $(\mathcal{C}, \mathcal{O})$ 上の
微分次数付き代数の層とする。$\mathcal{M}$ を微分次数付き
$\mathcal{A}$-加群とする。次の性質をもつ微分次数付き
$\mathcal{A}$-加群の準同型
$\mathcal{P} \to \mathcal{M}$ が存在する。
\begin{enumerate}
\item $\mathcal{P} \to \mathcal{M}$ は全射である。
\item $\Ker(\text{d}_\mathcal{P}) \to \Ker(\text{d}_\mathcal{M})$
は全射である。
\item $\mathcal{P}$ は良い。
\end{enumerate}
\end{lemma}

\begin{proof}
三つ組 $(U, k, x)$ で、$U$ が $\mathcal{C}$ の対象、
$k \in \mathbf{Z}$、$x$ が $U$ 上の $\mathcal{M}^k$ の切断で
$\text{d}_\mathcal{M}(x) = 0$ を満たすものを考える。
このとき、$1$ を $x$ に送る微分次数付き
$\mathcal{A}_U$-加群の一意な射
$\varphi_x : \mathcal{A}_U[-k] \to \mathcal{M}|_U$
を得る。これは射
$\psi_x : j_{U!}\mathcal{A}_U[-k] \to \mathcal{M}$
に随伴する。$1 \in \mathcal{A}_U(U)$ は、次数 $k$ で微分が零であり、
$\psi_x$ により $x$ に送られる切断
$1 \in j_{U!}\mathcal{A}_U[-k](U)$ に対応することに注意せよ。
したがって写像
$$
\bigoplus\nolimits_{(U, k, x)} j_{U!}\mathcal{A}_U[-k]
\longrightarrow
\mathcal{M}
$$
を考えると、条件 (2) と (3) が満たされる。実際、
対象 $j_{U!}\mathcal{A}_U[-k]$ は良く
(Lemma \ref{sdga-lemma-supply-good})、良い対象の任意の直和は良い
(Lemma \ref{sdga-lemma-good-direct-sum})。

\medskip\noindent
次に三つ組 $(U, k, x)$ で、$U$ が $\mathcal{C}$ の対象、
$k \in \mathbf{Z}$、$x$ が $\mathcal{M}^k$ の切断
（微分で消えるとは限らない）であるものを考える。
Remark \ref{sdga-remark-cone-identity} と同様に、恒等写像
$\mathcal{A}_U \to \mathcal{A}_U$ の錐 $C_U$ を考えられる。
元 $x$ は写像
$\varphi_x : C_U[-k - 1] \to \mathcal{A}_U$
を定める。Remark \ref{sdga-remark-cone-identity} を参照せよ。
許容短完全列
$$
0 \to \mathcal{A}_U \to C_U \to \mathcal{A}_U[1] \to 0
$$
があるので、Lemma \ref{sdga-lemma-good-admissible-ses} と既に用いた
Lemma \ref{sdga-lemma-supply-good} により $j_{U!}C_U$ は良い加群である。
上と同様に、$\varphi_x$ に随伴する写像
$\psi_x : j_{U!}C_U \to \mathcal{M}$ の直和
$$
\bigoplus\nolimits_{(U, k, x)} j_{U!}C_U \longrightarrow \mathcal{M}
$$
は全射であると分かる。これと第1段落で構成した写像との直和をとればよい。
\end{proof}

\begin{remark}
\label{sdga-remark-sheaf-graded-sets}
$(\mathcal{C}, \mathcal{O})$ を環付きサイトとする。
$\mathcal{C}$ 上の{\it 次数付き集合の層}とは、集合の層
$\mathcal{S}$ と集合の層の写像
$\deg : \mathcal{S} \to \underline{\mathbf{Z}}$ の組である。
次数付き集合の層 $\mathcal{S}$ 上の自由
$\mathcal{O}$-加群である次数付き $\mathcal{O}$-加群を
$\mathcal{O}[\mathcal{S}]$ と記す。より正確には、
$\mathcal{O}[\mathcal{S}]$ の次数 $n$ の部分は、規則
$$
U \longmapsto
\bigoplus\nolimits_{s \in \mathcal{S}(U),\ \deg(s) = n} s \cdot \mathcal{O}(U)
$$
の層化である。零微分を備えるものとして、これを微分次数付き
$\mathcal{O}$-加群とみなすこともできる。
$\mathcal{A}$ を次数付き $\mathcal{O}$-代数の層とする。
同様に、$\mathcal{A}[\mathcal{S}]$ を、次数 $n$ の部分が規則
$$
U \longmapsto
\bigoplus\nolimits_{s \in \mathcal{S}(U)} s \cdot \mathcal{A}^{n - \deg(s)}(U)
$$
の層化である次数付き $\mathcal{A}$-加群と定義する。
$\mathcal{A}$ が微分次数付き $\mathcal{O}$-代数ならば、
すべての $s \in \mathcal{S}(U)$ に対して $\text{d}(s) = 0$
とおいて層化することにより、これを微分次数付き
$\mathcal{O}$-加群にする。
\end{remark}

\begin{lemma}
\label{sdga-lemma-free-graded-module-good}
$(\mathcal{C}, \mathcal{O})$ を環付きサイトとする。
$\mathcal{A}$ を微分次数付き $\mathcal{A}$-代数とする。
$\mathcal{S}$ を $\mathcal{C}$ 上の次数付き集合の層とする。
このとき、$\mathcal{S}$ 上の自由次数付き加群
$\mathcal{A}[\mathcal{S}]$ に
Remark \ref{sdga-remark-sheaf-graded-sets} の微分を備えたものは、
良い微分次数付き $\mathcal{A}$-加群である。
\end{lemma}

\begin{proof}
$\mathcal{N}$ を次数付き左 $\mathcal{A}$-加群とする。このとき
$$
\mathcal{A}[\mathcal{S}] \otimes_\mathcal{A} \mathcal{N} =
\mathcal{O}[\mathcal{S}] \otimes_\mathcal{O} \mathcal{N} =
\mathcal{N}[\mathcal{S}]
$$
である。ここで $\mathcal{N}[\mathcal{S}$ は、次数 $n$ の部分が
前層
$$
U \longmapsto
\bigoplus\nolimits_{s \in \mathcal{S}(U)} s \cdot \mathcal{N}^{n - \deg(s)}(U)
$$
に付随する層である次数付き $\mathcal{O}$-加群である。
$\mathcal{N} \to \mathcal{N}[\mathcal{S}]$ が完全関手であることは
明らかなので、$\mathcal{A}[\mathcal{S}$ は次数付き
$\mathcal{A}$-加群として平坦である。次に、$\mathcal{N}$ を
微分次数付き左 $\mathcal{A}$-加群とする。このとき
$$
H^*(\mathcal{A}[\mathcal{S}] \otimes_\mathcal{A} \mathcal{N}) =
H^*(\mathcal{O}[\mathcal{S}] \otimes_\mathcal{O} \mathcal{N})
$$
が次数付き $\mathcal{O}$-加群の層として成り立ち、これは
（$\mathcal{O})$ 上の平坦性により
$$
H^*(\mathcal{N})[\mathcal{S}]
$$
に等しい。したがって $\mathcal{N}$ が非輪状ならば
$\mathcal{A}[\mathcal{S}] \otimes_\mathcal{A} \mathcal{N}$
は非輪状である。

\medskip\noindent
最後に、環付きトポスの射
$(f, f^\sharp) : (\Sh(\mathcal{C}'), \mathcal{O}')
\to (\Sh(\mathcal{C}), \mathcal{O})$、微分次数付き
$\mathcal{O}'$-代数 $\mathcal{A}'$、および微分次数付き
$f^{-1}\mathcal{O}$-代数の写像
$\varphi : f^{-1}\mathcal{A} \to \mathcal{A}'$ を考える。このとき
$$
f^*\mathcal{A}[\mathcal{S}] = \mathcal{A}'[f^{-1}\mathcal{S}]
$$
であることは直ちに分かり、これでこの加群が良いことの証明が完了する。
\end{proof}


\begin{lemma}
\label{sdga-lemma-resolve}
$(\mathcal{C}, \mathcal{O})$ を環付きサイトとする。
$\mathcal{A}$ を $(\mathcal{C}, \mathcal{O})$ 上の
微分次数付き代数の層とする。$\mathcal{M}$ を微分次数付き
$\mathcal{A}$-加群とする。次の性質をもつ微分次数付き
$\mathcal{A}$-加群の準同型
$\mathcal{P} \to \mathcal{M}$ が存在する。
\begin{enumerate}
\item $\mathcal{P} \to \mathcal{M}$ は擬同型である。
\item $\mathcal{P}$ は良い。
\end{enumerate}
\end{lemma}

\begin{proof}[第1の証明]
次数 $n$ の部分が $\Ker(\text{d}_\mathcal{M}^n)$ である
次数付き集合の層
$\mathcal{S}_0$ (Remark \ref{sdga-remark-sheaf-graded-sets}) をとる。
微分次数付き加群の準同型
$$
\mathcal{P}_0 = \mathcal{A}[\mathcal{S}_0] \longrightarrow \mathcal{M}
$$
を考える。ここで左辺は
Remark \ref{sdga-remark-sheaf-graded-sets} のとおりであり、写像は
$\mathcal{S}_0$ の局所切断 $s$ を $\mathcal{M}^{\deg(s)}$ の
対応する局所切断に送る（これは微分の核に入るので、確かに
微分次数付き加群の写像である）。構成により、上の写像から誘導される
コホモロジーの層の写像
$H^n(\mathcal{P}_0) \to H^n(\mathcal{M})$ は全射である。
写像
$$
\mathcal{P}_0 \to \mathcal{P}_1 \to \mathcal{P}_2 \to \ldots \to \mathcal{M}
$$
を帰納的に構成する。もちろん、すべての $i$ に対して
$H^*(\mathcal{P}_i) \to H^*(\mathcal{M})$ は全射となる。
$\mathcal{P}_i \to \mathcal{M}$ が与えられたとき、
次数 $n$ の部分が
$$
\Ker(\text{d}_{\mathcal{P}_i}^{n + 1})
\times_{\mathcal{M}^{n + 1}, \text{d}}
\mathcal{M}^n
$$
である次数付き集合の層を $\mathcal{S}_{i + 1}$ と記す。
次数付き $\mathcal{A}$-加群として
$$
\mathcal{P}_{i + 1} = \mathcal{P}_i \oplus \mathcal{A}[\mathcal{S}_{i + 1}]
$$
とおき、微分および $\mathcal{M}$ への写像を次のように定義する。
\begin{enumerate}
\item $\mathcal{P}_i$ の局所切断に対しては、
$\mathcal{P}_i$ 上の微分と与えられた $\mathcal{M}$ への写像を用いる。
\item $\mathcal{S}_{i + 1}$ の局所切断 $s = (p, m)$ に対しては、
$\text{d}(s)$ を次数 $\deg(s) + 1$ の $\mathcal{P}_i$ の切断と
みなした $p$ とおき、$s$ を $\mathcal{M}$ の $m$ に送る。
\item Leibniz 則を満たすように微分を一意に延長する。
\end{enumerate}
$\text{d}(m)$ は $p$ の像であり $\text{d}(p) = 0$ なので、
この定義は意味をもつ。最後に
$\mathcal{P} = \colim \mathcal{P}_i$ とおき、
$\mathcal{M}$ への誘導写像を備える。

\medskip\noindent
写像 $\mathcal{P} \to \mathcal{M}$ は擬同型である。実際
$H^n(\mathcal{P}) = \colim H^n(\mathcal{P}_i)$ であり、
各 $i$ について写像
$H^n(\mathcal{P}_i) \to H^n(\mathcal{M})$ は全射で、その核は構成上
$H^n(\mathcal{P}_i) \to H^n(\mathcal{P}_{i + 1})$ により零に送られる。
各 $\mathcal{P}_i$ は良い。これは
Lemma \ref{sdga-lemma-free-graded-module-good} により $\mathcal{P}_0$ が良く、
各 $\mathcal{P}_{i + 1}$ は許容短完全列
$0 \to \mathcal{P}_i \to \mathcal{P}_{i + 1} \to
\mathcal{A}[\mathcal{S}_{i + 1}] \to 0$
の中項であり、帰納法によりその両端が良いからである。したがって
Lemma \ref{sdga-lemma-good-admissible-ses} により
$\mathcal{P}_{i + 1}$ は良い。最後に
Lemma \ref{sdga-lemma-good-direct-sum} により $\mathcal{P}$ は良い。
\end{proof}

\begin{proof}[第2の証明]
この証明を読む前に Differential Graded Algebra,
Lemma \ref{dga-lemma-resolve} の証明を読むことを強く勧める。
$\mathcal{M} = \mathcal{M}_0$ とおく。写像
$\mathcal{P}_i \to \mathcal{M}_i$ を
Lemma \ref{sdga-lemma-good-quotient} のように選び、短完全列
$$
0 \to \mathcal{M}_{i + 1} \to \mathcal{P}_i \to \mathcal{M}_i \to 0
$$
を帰納的に選ぶ。これにより「分解」
$$
\ldots \to \mathcal{P}_2 \xrightarrow{f_2}
\mathcal{P}_1 \xrightarrow{f_1} \mathcal{P}_0 \to \mathcal{M} \to 0
$$
を得る。次に、微分次数付き $\mathcal{A}$-加群 $\mathcal{P}$ を
次のように定義する。
\begin{enumerate}
\item 次数付き $\mathcal{A}$-加群として
$\mathcal{P} = \bigoplus_{a \leq 0} \mathcal{P}_{-a}[-a]$ とおく。
すなわち次数 $n$ の部分は
$\mathcal{P}^n = \bigoplus\nolimits_{a + b = n} \mathcal{P}_{-a}^b$
である。
\item $\mathcal{P}$ 上の微分は、二重複体に付随する全複体の構成と
同様に、$\mathcal{P}_{-a}^b$ の局所切断 $x$ に対して
$$
\text{d}_\mathcal{P}(x) = f_{-a}(x) + (-1)^a \text{d}_{\mathcal{P}_{-a}}(x)
$$
で与える。
\end{enumerate}
これらの規約のもとで $\mathcal{P}$ は確かに微分次数付き
$\mathcal{A}$-加群である。詳細は省略する。微分次数付き
$\mathcal{A}$-加群の写像 $\mathcal{P} \to \mathcal{M}$ があり、
$a < 0$ では直和因子 $\mathcal{P}_{-a}[-a]$ 上で零、
$a = 0$ では与えられた写像
$\mathcal{P}_0 \to \mathcal{M}$ である。
$$
\mathcal{P} = \colim_i F_i\mathcal{P}
$$
であることに注意せよ。ここで $F_i\mathcal{P} \subset \mathcal{P}$
は、台となる次数付き $\mathcal{A}$-加群が
$$
F_i\mathcal{P} =
\bigoplus\nolimits_{i \geq -a \geq 0} \mathcal{P}_{-a}[-a]
$$
である微分次数付き $\mathcal{A}$-部分加群である。写像
$$
0 \to F_1\mathcal{P} \to F_2\mathcal{P} \to F_3\mathcal{P} \to
\ldots \to \mathcal{P}
$$
はすべて許容単射であり、許容短完全列
$$
0 \to F_i\mathcal{P} \to
F_{i + 1}\mathcal{P} \to
\mathcal{P}_{i + 1}[i + 1] \to 0
$$
があることは直ちに分かる。帰納法と
Lemma \ref{sdga-lemma-good-admissible-ses} により、
$F_i\mathcal{P}$ は良い微分次数付き $\mathcal{A}$-加群である。
$\mathcal{P} = \colim F_i\mathcal{P}$ なので
Lemma \ref{sdga-lemma-good-direct-sum} により $\mathcal{P}$ は良い。

\medskip\noindent
最後に、$\mathcal{P} \to \mathcal{M}$ が擬同型であることを
示さなければならない。$\mathcal{C}$ が十分多くの点をもつならば、
これは茎で確認することにより、初等的な Homology,
Lemma \ref{homology-lemma-good-resolution-gives-qis} から従う。
一般には次のように論じられる（この証明はあまりにも長い。
初等的な証明と同じく局所切断を用いる別の議論もあるが、
それもかなり長い）。アーベル群の層の圏では
フィルター余極限が完全なので
$$
H^d(\mathcal{P}) = \colim H^d(F_i\mathcal{P})
$$
である。各 $i \geq 0$ と $d \in \mathbf{Z}$ に対して、
(a) 短完全列
$$
0 \to H^d(\mathcal{M}_{i + 1}[i]) \to
H^d(F_i\mathcal{P}) \to H^d(\mathcal{M}) \to 0
$$
があり、ここで第2の矢印は
$F_i\mathcal{P} \to \mathcal{P} \to \mathcal{M}$ から来ること、
および (b) 合成
$$
H^d(\mathcal{M}_{i + 1}[i]) \to
H^d(F_i\mathcal{P}) \to
H^d(F_{i + 1}\mathcal{P})
$$
が零であることを主張する。この主張が証明を完了するのに十分であることは
明らかである。

\medskip\noindent
主張を証明する。任意の $i \geq 0$ に対し、
$f_i$ の核としての $\mathcal{M}_{i + 1}$ の
$\mathcal{P}_i$ への包含から写像
$\mathcal{M}_{i + 1}[i] \to F_i\mathcal{P}$ が得られる。
$C_i$ を定義する $\mathcal{O}$-加群の複体の短完全列
$$
0 \to \mathcal{M}_{i + 1}[i] \to F_i\mathcal{P} \to C_i \to 0
$$
を考える。$C_0 = \mathcal{M}_0 = \mathcal{M}$ である。
また $C_i$ は、次数 $-i, -i + 1, \ldots, 0$ にそれぞれ列
$$
\mathcal{M}_i = \mathcal{P}_i/\mathcal{M}_{i + 1},
\mathcal{P}_{i - 1}, \ldots, \mathcal{P}_0
$$
をもつ二重複体 $C_i^{\bullet, \bullet}$ に付随する全複体である。
二重複体の写像
$C_i^{\bullet, \bullet} \to C_{i - 1}^{\bullet, \bullet}$ があり、
次数 $-i$ の列上では $0$、次数 $-i + 1$ では全射
$\mathcal{P}_{i - 1} \to \mathcal{M}_{i - 1}$、
他の列上では恒等写像である。したがって複体の写像
$$
C_i \longrightarrow C_{i - 1}
$$
を得る。これらは全射な擬同型である。実際、その核は次数
$-i, -i + 1$ の二つの列が $\mathcal{M}_i, \mathcal{M}_i$ であり、
その二列の間の写像が恒等写像である二重複体の全複体である。
得られる同一視
$H^d(C_i) = H^d(C_{i - 1} = \ldots = H^d(\mathcal{M})$
を用いると、上の複体の短完全列から既に長完全列
$$
H^d(\mathcal{M}_{i + 1}[i]) \to
H^d(F_i\mathcal{P}) \to H^d(\mathcal{M}) \to
H^{d + 1}(\mathcal{M}_{i + 1}[i])
$$
を得ることが分かる。しかし可換図式
$$
\xymatrix{
\mathcal{M}_{i + 2}[i + 1] \ar[r]_a & T_{i + 1} \ar[r] &
F_{i + 1}\mathcal{P} \ar[r] & C_{i + 1} \ar[d] \\
& \mathcal{M}_{i + 1}[i] \ar[r] \ar[u]^b &
F_i\mathcal{P} \ar[u] \ar[r] &
C_i
}
$$
もある。ここで $T_{i + 1}$ は次数 $-i - 1$ と $-i$ にそれぞれ
$\mathcal{P}_{i + 1}, \mathcal{M}_{i + 1}$ を列として置いた
二重複体の全複体である。言い換えれば、$T_{i + 1}$ は写像
$\mathcal{P}_{i + 1} \to \mathcal{M}_{i + 1}$ の錐のシフトであり、
$a$ は擬同型、写像 $a^{-1} \circ b$ は短完全列
$$
0 \to \mathcal{M}_{i + 2} \to \mathcal{P}_{i + 1} \to \mathcal{M}_{i + 1} \to 0
$$
に付随する $D(\mathcal{O})$ の識別三角形の第3写像のシフトである。
写像
$H^d(\mathcal{P}_{i + 1}) \to H^d(\mathcal{M}_{i + 1})$
は全射である。これは写像を
$\Ker(\text{d}_{\mathcal{P}_{i + 1}}) \to
\Ker(\text{d}_{\mathcal{M}_{i + 1}})$
が全射となるように選んだからである。したがって
$a^{-1} \circ b$ はコホモロジーの層上で零である。
これで主張の (b) が証明された。
$T_{i + 1}$ は複体の全射
$F_{i + 1}\mathcal{P} \to C_i$ の核なので、
コホモロジー長完全列の写像
$$
\xymatrix{
H^d(T_{i + 1}) \ar[r] &
H^d(F_{i + 1}\mathcal{P}) \ar[r] &
H^d(\mathcal{M}) \ar[r] &
H^{d + 1}(T_{i + 1}) \\
H^d(\mathcal{M}_{i + 1}[i]) \ar[r] \ar[u] &
H^d(F_i\mathcal{P}) \ar[r] \ar[u] &
H^d(\mathcal{M}) \ar[r] \ar[u] &
H^{d + 1}(\mathcal{M}_{i + 1}[i]) \ar[u]
}
$$
を得る。上の議論により、外側の縦写像は零であることが分かっている。
したがって写像
$H^d(F_{i + 1}\mathcal{P}) \to H^d(\mathcal{M})$
は全射であり、主張の (a) が従う。より正確には、$i > 0$ に対して
主張が従う。$i = 0$ の場合は読者に委ねる
（実際、この補題を得るにはすべての $i \gg 0$ に対して
主張を証明すれば十分である）。
\end{proof}


\begin{lemma}
\label{sdga-lemma-acyclic-good}
$(\mathcal{C}, \mathcal{O})$ を環付きサイトとする。
$\mathcal{A}$ を $(\mathcal{C}, \mathcal{O})$ 上の
微分次数付き代数の層とする。$\mathcal{P}$ を良い非輪状な
右微分次数付き $\mathcal{A}$-加群とする。
\begin{enumerate}
\item 任意の微分次数付き左 $\mathcal{A}$-加群 $\mathcal{N}$ に対し、
テンソル積 $\mathcal{P} \otimes_\mathcal{A} \mathcal{N}$ は
非輪状である。
\item 環付きトポスの任意の射
$(f, f^\sharp) : (\Sh(\mathcal{C}'), \mathcal{O}')
\to (\Sh(\mathcal{C}), \mathcal{O})$、任意の微分次数付き
$\mathcal{O}'$-代数 $\mathcal{A}'$、および微分次数付き
$f^{-1}\mathcal{O}$-代数の任意の写像
$\varphi : f^{-1}\mathcal{A} \to \mathcal{A}'$ に対して、
引き戻し $f^*\mathcal{P}$ は非輪状かつ良い。
\end{enumerate}
\end{lemma}

\begin{proof}
(1) の証明。Lemma \ref{sdga-lemma-resolve} により、
良い左微分次数付き $\mathcal{Q}$ と擬同型
$\mathcal{Q} \to \mathcal{N}$ を選べる。このとき
$\mathcal{Q}$ は良いので
$\mathcal{P} \otimes_\mathcal{A} \mathcal{Q}$ は非輪状である。
写像 $\mathcal{Q} \to \mathcal{N}$ の錐を
$\mathcal{N}'$ とする。$\mathcal{P}$ が良く、$\mathcal{N}'$ は
擬同型の錐として非輪状なので
$\mathcal{P} \otimes_\mathcal{A} \mathcal{N}'$ は非輪状である。
三角構造の構成により
$K(\textit{Mod}(\mathcal{A}, \text{d}))$ に識別三角形
$$
\mathcal{Q} \to \mathcal{N} \to \mathcal{N}' \to \mathcal{Q}[1]
$$
がある。$\mathcal{P} \otimes_\mathcal{A} -$ は識別三角形を
識別三角形に送るので、$K(\textit{Mod}(\mathcal{O}))$ に識別三角形
$$
\mathcal{P} \otimes_\mathcal{A} \mathcal{Q} \to
\mathcal{P} \otimes_\mathcal{A} \mathcal{N} \to
\mathcal{P} \otimes_\mathcal{A} \mathcal{N}' \to
\mathcal{P} \otimes_\mathcal{A} \mathcal{Q}[1]
$$
を得る。したがって結論が従う。

\medskip\noindent
(2) の証明。良い加群の定義により $f^*\mathcal{P}$ は良い。
ここで
$f^*\mathcal{P} = f^{-1}\mathcal{P} \otimes_{f^{-1}\mathcal{A}} \mathcal{A}'$
であった。$f^{-1}$ は完全なので、$f^{-1}\mathcal{P}$ は良い非輪状な
微分次数付き $f^{-1}\mathcal{A}$-加群である。したがって (1) により
$f^*\mathcal{P}$ は非輪状である。
\end{proof}








\section{次数付き加群の微分次数付き包}
\label{sdga-section-dg-hull}

\noindent
次数付き加群 $\mathcal{N}$ の微分次数付き包とは、
次の補題の関手 $G$ を適用して得られるものである。

\begin{lemma}
\label{sdga-lemma-dg-hull}
$(\mathcal{C}, \mathcal{O})$ を環付きサイトとする。
$\mathcal{A}$ を $(\mathcal{C}, \mathcal{O})$ 上の
微分次数付き代数の層とする。忘却関手
$F : \textit{Mod}(\mathcal{A}, \text{d}) \to \textit{Mod}(\mathcal{A})$
は左随伴
$G : \textit{Mod}(\mathcal{A}) \to
\textit{Mod}(\mathcal{A}, \text{d})$ をもつ。
\end{lemma}

\begin{proof}
$G$ の存在を証明するには随伴関手定理を用いればよい。
Categories, Theorem \ref{categories-theorem-adjoint-functor} を参照せよ
（$F$ と $G$ の役割を入れ替えていることに注意せよ）。
$F$ に対する完全性の条件は
Lemma \ref{sdga-lemma-dgm-abelian} により満たされる。
集合論的条件は次のように確認できる。次数付き
$\mathcal{A}$-加群 $\mathcal{N}$ が与えられたとする。任意の写像
$$
\varphi : \mathcal{N} \longrightarrow F(\mathcal{M})
$$
に対し、$\Im(\varphi) \subset F(\mathcal{M}')$ を満たす最小の
微分次数付き $\mathcal{A}$-部分加群
$\mathcal{M}' \subset \mathcal{M}$ を考えられる。
$\mathcal{M}'$ は、次数付き $\mathcal{A}$-加群の写像
$$
\mathcal{N} \oplus
\mathcal{N}[-1] \otimes_\mathcal{O} \mathcal{A}
\longrightarrow
\mathcal{M}
$$
で
$$
(n, \sum n_i \otimes a_i) \longmapsto
\varphi(n) + \sum \text{d}(\varphi(n_i)) a_i
$$
と定義されるものの像であることは明らかである。実際、この写像の像が
$\mathcal{M}$ の微分次数付き部分加群であることは容易に分かる。
したがって、このような $\mathcal{M}'$ の可能な同型類の数は
有界であり、結論を得る。
\end{proof}

\noindent
$(\mathcal{C}, \mathcal{O})$ を環付きサイトとする。
$\mathcal{A}$ を $(\mathcal{C}, \mathcal{O})$ 上の
微分次数付き代数の層とする。$\mathcal{M}$ を微分次数付き
$\mathcal{A}$-加群とし、$\textit{Mod}(\mathcal{A})$ に短完全列
$$
0 \to \mathcal{N} \to F(\mathcal{M}) \to \mathcal{N}' \to 0
$$
があるとする。このとき標準的な次数付き $\mathcal{A}$-加群準同型
$$
\overline{\text{d}} : \mathcal{N} \to \mathcal{N}'[1]
$$
を次のように得る。$\mathcal{N}$ の局所切断 $x$ が与えられたとき、
$x$ を $\mathcal{M}$ の局所切断とみなし、
$\text{d}_\mathcal{M}(x)$ の $\mathcal{N}'$ における像を
$\overline{\text{d}}(x)$ と記す。

\begin{lemma}
\label{sdga-lemma-dg-hull-acyclic}
Lemma \ref{sdga-lemma-dg-hull} の関手 $F, G$ は次の性質をもつ。
次数付き $\mathcal{A}$-加群 $\mathcal{N}$ が与えられたとき
\begin{enumerate}
\item 余単位 $\mathcal{N} \to F(G(\mathcal{N}))$ は単射である。
\item 写像 $\overline{\text{d}} : \mathcal{N} \to
\Coker(\mathcal{N} \to F(G(\mathcal{N})))[1]$ は同型である。
\item $G(\mathcal{N})$ は非輪状な微分次数付き
$\mathcal{A}$-加群である。
\end{enumerate}
\end{lemma}

\begin{proof}
性質 (3) は性質 (1) と (2) の帰結であることに注意する。
実際、$s$ が $F(G(\mathcal{N}))$ の非零局所切断で
$\text{d}(s) = 0$ ならば、$s$ は
$\mathcal{N} \to F(G(\mathcal{N}))$ の像に入らない。
したがって、$s$ の余核における像 $\overline{s}$ は、
$\mathcal{N}$ のある局所切断 $s'$ に対して
$\overline{\text{d}}(s')$ と書ける。すると
$s - \text{d}(s')$ は依然として $\text{d}$ の核に属し、
余単位の像に含まれるので、$s = \text{d}(s')$ である。

\medskip\noindent
一時的に、層 $\mathcal{A}$ をそれ自身上の（右）次数付き加群、
それぞれそれ自身上の（右）微分次数付き加群とみなしたものを
$\mathcal{A}_{gr}$、それぞれ $\mathcal{A}_{dg}$ と書く。
この補題で最も重要な場合は $G(\mathcal{A}_{gr})$ を理解することである。
もちろん $G(\mathcal{A}_{gr})$ は、関手
$$
\mathcal{M} \longmapsto
\Hom_{\textit{Mod}(\mathcal{A})}(\mathcal{A}_{gr}, F(\mathcal{M})) =
\Gamma(\mathcal{C}, \mathcal{M})
$$
を表現する $\textit{Mod}(\mathcal{A}, \text{d})$ の対象である。
Remark \ref{sdga-remark-cone-identity} により、この関手は
$C[-1]$ によって表現される。ここで $C$ は
$\mathcal{A}_{dg}$ の恒等写像の錐である。
$\textit{Mod}(\mathcal{A}, \text{d})$ に短完全列
$$
0 \to \mathcal{A}_{dg}[-1] \to C[-1] \to \mathcal{A}_{dg} \to 0
$$
があり、これは $\textit{Mod}(\mathcal{A})$ において余単位
$\mathcal{A}_{gr} \to F(C[-1])$ により分裂する。したがって
$G(\mathcal{A}_{gr})$ は性質 (1) と (2) を満たす。

\medskip\noindent
$U$ を $\mathcal{C}$ の対象とする。局所化射を
$j_U : \mathcal{C}/U \to \mathcal{C}$ と記し、
$\mathcal{A}$ の $U$ への制限を $\mathcal{A}_U$ と記す。
$\mathcal{A}_U$ を次数付き $\mathcal{A}_U$-加群とみなしたものを
$\mathcal{A}_{U, gr}$ と記す。忘却関手を
$F_U : \textit{Mod}(\mathcal{A}_U, \text{d}) \to
\textit{Mod}(\mathcal{A}_U)$ と記し、その随伴を $G_U$ と記す。
このとき可換図式
$$
\vcenter{
\xymatrix{
\textit{Mod}(\mathcal{A}, \text{d}) \ar[d]_{j_U^*} \ar[r]_F &
\textit{Mod}(\mathcal{A}) \ar[d]^{j_U^*} \\
\textit{Mod}(\mathcal{A}_U, \text{d}) \ar[r]^{F_U} &
\textit{Mod}(\mathcal{A}_U)
}
}
\quad\text{and}\quad
\vcenter{
\xymatrix{
\textit{Mod}(\mathcal{A}_U, \text{d}) \ar[r]_{F_U} \ar[d]_{j_{U!}} &
\textit{Mod}(\mathcal{A}_U) \ar[d]^{j_{U!}} \\
\textit{Mod}(\mathcal{A}, \text{d}) \ar[r]^F &
\textit{Mod}(\mathcal{A})
}
}
$$
を得る。これは Sections \ref{sdga-section-functoriality-graded},
\ref{sdga-section-functoriality-dg}, \ref{sdga-section-localize-graded},
\ref{sdga-section-localize-dg} における $j_U^*$ と $j_{U!}$ の構成による。
随伴の一意性から $j_{U!} \circ G_U = G \circ j_{U!}$ を得る。
$j_{U!}$ は完全関手なので、証明の前半で見た余単位
$\mathcal{A}_{U, gr} \to F_U(G_U(\mathcal{A}_{U, gr}))$
に対する性質 (1) と (2) から、余単位
$j_{U!}\mathcal{A}_{U, gr} \to F(G(j_{U!}\mathcal{A}_{U, gr})) =
j_{U!}F_U(G_U(\mathcal{A}_{U, gr}))$
に対する性質 (1) と (2) が従う。

\medskip\noindent
Lemma \ref{sdga-lemma-gm-grothendieck-abelian} の証明では、
$\textit{Mod}(\mathcal{A})$ の任意の対象が
$j_{U!}\mathcal{A}_{U, gr}$ のコピーの直和の商であることを見た。
$G$ は左随伴なので直和と可換する。したがって性質 (1) と (2) は、
それらを満たす対象の直和に対しても成立する。ゆえに
$\textit{Mod}(\mathcal{A})$ の任意の対象 $\mathcal{N}$ は、
$\mathcal{N}_1$ と $\mathcal{N}_0$ に対して (1) と (2) が成立する
完全列
$$
\mathcal{N}_1 \to \mathcal{N}_0 \to \mathcal{N} \to 0
$$
に入る。$G$ が右完全であることを用いて $\mathcal{N}$ に対する
(1) と (2) を導くことは読者に委ねる。
\end{proof}










\section{K-入射微分次数付き加群}
\label{sdga-section-K-injective}

\noindent
本節は、微分次数付き代数の層上の微分次数付き加群の層という設定における
Injectives, Section \ref{injectives-section-K-injective} に対応する。

\begin{lemma}
\label{sdga-lemma-characterize-injectives}
$(\mathcal{C}, \mathcal{O})$ を環付きサイトとする。
$\mathcal{A}$ を $(\mathcal{C}, \mathcal{O})$ 上の
次数付き代数の層とする。集合 $T$ と、各 $t \in T$ に対する
次数付き $\mathcal{A}$-加群の単射
$\mathcal{N}_t \to \mathcal{N}'_t$ が存在し、
$\textit{Mod}(\mathcal{A})$ の対象 $\mathcal{I}$ が入射的であることと、
すべての実線図式
$$
\xymatrix{
\mathcal{N}_t \ar[r] \ar[d] & \mathcal{I} \\
\mathcal{N}'_t \ar@{..>}[ru]
}
$$
に対して、図式を可換にする点線矢印が
$\textit{Mod}(\mathcal{A})$ に存在することとは同値である。
\end{lemma}

\begin{proof}
これは任意の Grothendieck アーベル圏で成立する。Injectives,
Lemma \ref{injectives-lemma-characterize-injective} を参照せよ。
Lemma \ref{sdga-lemma-gm-grothendieck-abelian} により、
圏 $\textit{Mod}(\mathcal{A})$ は Grothendieck アーベル圏である。
\end{proof}

\begin{definition}
\label{sdga-definition-graded-injective}
$(\mathcal{C}, \mathcal{O})$ を環付きサイトとする。
$(\mathcal{A}, \text{d})$ を $(\mathcal{C}, \mathcal{O})$ 上の
微分次数付き代数の層とする。微分次数付き
$\mathcal{A}$-加群 $\mathcal{I}$ が{\it 次数付き入射的}
\footnote{これは標準的な用語ではないかもしれない。}であるとは、
$\mathcal{M}$ を次数付き $\mathcal{A}$-加群とみなしたものが、
次数付き $\mathcal{A}$-加群の圏 $\textit{Mod}(\mathcal{A})$ の
入射対象であることをいう。
\end{definition}

\begin{remark}
\label{sdga-remark-why-graded-injective}
$(\mathcal{C}, \mathcal{O})$ を環付きサイトとする。
$(\mathcal{A}, \text{d})$ を $(\mathcal{C}, \mathcal{O})$ 上の
微分次数付き代数の層とする。$\mathcal{I}$ を次数付き入射的な
微分次数付き $\mathcal{A}$-加群とする。
$$
0 \to \mathcal{M}_1 \to \mathcal{M}_2 \to \mathcal{M}_3 \to 0
$$
を微分次数付き $\mathcal{A}$-加群の短完全列とする。
$\mathcal{I}$ は次数付き入射的なので、
$\Gamma(\mathcal{C}, \mathcal{O})$-加群の複体の短完全列
$$
0 \to
\Hom_{\textit{Mod}^{dg}(\mathcal{A}, \text{d})}(\mathcal{M}_3, \mathcal{I})
\to
\Hom_{\textit{Mod}^{dg}(\mathcal{A}, \text{d})}(\mathcal{M}_2, \mathcal{I})
\to
\Hom_{\textit{Mod}^{dg}(\mathcal{A}, \text{d})}(\mathcal{M}_1, \mathcal{I})
\to 0
$$
を得る。コホモロジーをとると、ホモトピー圏における準同型群の長完全列
$$
\xymatrix{
\Hom_{K(\textit{Mod}(\mathcal{A}, \text{d}))}(\mathcal{M}_3, \mathcal{I})
\ar[d] &
\Hom_{K(\textit{Mod}(\mathcal{A}, \text{d}))}(\mathcal{M}_3, \mathcal{I})[1]
\ar[d] \\
\Hom_{K(\textit{Mod}(\mathcal{A}, \text{d}))}(\mathcal{M}_2, \mathcal{I})
\ar[d] &
\Hom_{K(\textit{Mod}(\mathcal{A}, \text{d}))}(\mathcal{M}_2, \mathcal{I})[1]
\ar[d] \\
\Hom_{K(\textit{Mod}(\mathcal{A}, \text{d}))}(\mathcal{M}_1, \mathcal{I})
\ar[ruu]
&
\Hom_{K(\textit{Mod}(\mathcal{A}, \text{d}))}(\mathcal{M}_1, \mathcal{I})[1]
}
$$
を得る。重要なのは、この短完全列が許容であるとは仮定していないにも
かかわらず、これが得られることである
（したがって一般には、この短完全列はホモトピー圏の識別三角形を
定義しない）。
\end{remark}

\begin{lemma}
\label{sdga-lemma-product-graded-injective}
$(\mathcal{C}, \mathcal{O})$ を環付きサイトとする。
$(\mathcal{A}, \text{d})$ を $(\mathcal{C}, \mathcal{O})$ 上の
微分次数付き代数の層とする。$T$ を集合とし、各 $t \in T$ に対して
$\mathcal{I}_t$ を次数付き入射的な微分次数付き
$\mathcal{A}$-加群とする。このとき
$\prod \mathcal{I}_t$ は次数付き入射的な微分次数付き
$\mathcal{A}$-加群である。
\end{lemma}

\begin{proof}
入射対象の直積は入射的だからである。Homology,
Lemma \ref{homology-lemma-product-injectives} を参照せよ。
また、$\textit{Mod}(\mathcal{A}, \text{d})$ における直積は
忘却関手を介して $\textit{Mod}(\mathcal{A})$ における直積と両立する。
\end{proof}

\begin{lemma}
\label{sdga-lemma-characterize-graded-injectives-in-dg}
$(\mathcal{C}, \mathcal{O})$ を環付きサイトとする。
$(\mathcal{A}, \text{d})$ を $(\mathcal{C}, \mathcal{O})$ 上の
微分次数付き代数の層とする。集合 $T$ と、各 $t \in T$ に対する
非輪状な微分次数付き $\mathcal{A}$-加群の単射
$\mathcal{M}_t \to \mathcal{M}'_t$ が存在し、
$\textit{Mod}(\mathcal{A}, \text{d})$ の対象 $\mathcal{I}$ に対して
次は同値である。
\begin{enumerate}
\item $\mathcal{I}$ は次数付き入射的である。
\item すべての実線図式
$$
\xymatrix{
\mathcal{M}_t \ar[r] \ar[d] & \mathcal{I} \\
\mathcal{M}'_t \ar@{..>}[ru]
}
$$
に対して、図式を可換にする点線矢印が
$\textit{Mod}(\mathcal{A}, \text{d})$ に存在する。
\end{enumerate}
\end{lemma}

\begin{proof}
$T$ と $\mathcal{N}_t \to \mathcal{N}'_t$ を
Lemma \ref{sdga-lemma-characterize-injectives} のものとする。
忘却関手を
$F : \textit{Mod}(\mathcal{A}, \text{d}) \to \textit{Mod}(\mathcal{A})$
と記し、Lemma \ref{sdga-lemma-dg-hull} のように $F$ の左随伴関手を
$G$ とする。
$$
\mathcal{M}_t = G(\mathcal{N}_t) \to
G(\mathcal{N}'_t) = \mathcal{M}'_t
$$
とおく。Lemma \ref{sdga-lemma-dg-hull-acyclic} により、
これは非輪状な微分次数付き $\mathcal{A}$-加群の単射である。
$G$ は $F$ の左随伴なので、図式
$$
\xymatrix{
\mathcal{M}_t \ar[r] \ar[d] & \mathcal{I} \\
\mathcal{M}'_t \ar@{..>}[ru]
}
$$
に点線矢印が存在することと、図式
$$
\xymatrix{
\mathcal{N}_t \ar[r] \ar[d] & F(\mathcal{I}) \\
\mathcal{N}'_t \ar@{..>}[ru]
}
$$
に点線矢印が存在することとは同値である。したがって結論は、
矢印族 $\mathcal{N}_t \to \mathcal{N}_t'$ の選び方から従う。
\end{proof}



\begin{lemma}
\label{sdga-lemma-small-acyclics}
$(\mathcal{C}, \mathcal{O})$ を環付きサイトとする。
$(\mathcal{A}, \text{d})$ を $(\mathcal{C}, \mathcal{O})$ 上の
微分次数付き代数の層とする。集合 $S$ と、各 $s$ に対する
非輪状な微分次数付き $\mathcal{A}$-加群 $\mathcal{M}_s$ が存在し、
任意の非零で非輪状な微分次数付き $\mathcal{A}$-加群
$\mathcal{M}$ に対して、ある $s \in S$ と
$\textit{Mod}(\mathcal{A}, \text{d})$ における単射
$\mathcal{M}_s \to \mathcal{M}$ が存在する。
\end{lemma}

\begin{proof}
始める前に、規約によりサイト $\mathcal{C}$ の対象と射は集合をなし、
被覆の集合 $\text{Cov}(\mathcal{C})$ があることを思い出そう。
$\mathcal{F}$ が微分次数付き $\mathcal{A}$-加群ならば、
$U$ が $\mathcal{C}$ の対象を、$n \in \mathbf{Z}$ が整数を動くときの
$$
\coprod\nolimits_{(U, n)} \mathcal{F}^n(U)
$$
の濃度の和を $|\mathcal{F}|$ と定義する。
$|\Ob(\mathcal{C})|$、$|\text{Arrows}(\mathcal{C})|$、
$|\text{Cov}(\mathcal{C})|$、
$\{U_i \to U\}_{i \in I} \in \text{Cov}(\mathcal{C})$
に対する $\sup |I|$、および $|\mathcal{A}|$ より大きい
無限基数 $\kappa$ を選ぶ。

\medskip\noindent
$\mathcal{F} \subset \mathcal{M}$ を微分次数付き
$\mathcal{A}$-加群の包含とする。集合 $K$ と、各 $k \in K$ に対する
$\mathcal{C}$ の対象 $U_k$、整数 $n_k$、切断
$x_k \in \mathcal{M}^{n_k}(U_k)$ からなる三つ組
$(U_k, n_k, x_k)$ が与えられたとする。
$\mathcal{F}$ とすべての $k \in K$ に対する切断 $x_k$ を含む
最小の微分次数付き $\mathcal{A}$-部分加群
$\mathcal{F}' \subset \mathcal{M}$ を考えられる。
$$
(\mathcal{F}')^n(U) \subset \mathcal{M}^n(U)
$$
は、次の性質をもつ元 $x \in \mathcal{M}^n(U)$ の集合として
記述できる。すなわち、被覆
$\{f_i : U_i \to U\}_{i \in I} \in \text{Cov}(\mathcal{C})$
が存在して、各 $i \in I$ に対し有限集合 $T_i$ と射
$g_t : U_i \to U_{k_t}$ があり、
$$
f_i^*x = y_i +
\sum\nolimits_{t \in T_i} a_{it}g_t^*x_{k_t} + b_{it}g_t^*\text{d}(x_{k_t})
$$
となることである。ここで $y_i \in \mathcal{F}^n(U)$、
$a_{it} \in \mathcal{A}^{n - n_{k_t}}(U_i)$、
$b_{it} \in \mathcal{A}^{n - n_{k_t} - 1}(U_i)$ である。
（詳細は省略する。ヒント：これらの切断が確かに
$\mathcal{F}'$ に属することを確認し、逆にこの規則が
微分次数付き $\mathcal{A}$-部分加群を定義することを示す。）
この記述から
$|\mathcal{F}'| \leq \max(|\mathcal{F}|, |K|, \kappa)$ が従う。

\medskip\noindent
$\mathcal{M}$ を非零で非輪状な微分次数付き
$\mathcal{A}$-加群とする。このとき、ある整数 $n$、
$\mathcal{C}$ のある対象 $U$、および $U$ 上の
$\mathcal{M}^n$ の非零切断 $x$ が存在する。
$$
\mathcal{F}_0 \subset \mathcal{M}
$$
を $x$ を含む最小の微分次数付き $\mathcal{A}$-部分加群とする。
前段落により $|\mathcal{F}_0| \leq \kappa$ である。
$\mathcal{F}_0, \ldots, \mathcal{F}_n$ が与えられたとき、
$\mathcal{F}_{n + 1}$ を帰納的に次のように定義する。集合
$$
L = \{(U, n, x)\}
\{U_i \to U\}_{i \in I}, (x_i)_{i \in I})\}
$$
を考える。これは、$U$ が $\mathcal{C}$ の対象、
$n \in \mathbf{Z}$、$x \in \mathcal{F}_n(U)$ かつ
$\text{d}(x) = 0$ である三つ組の集合である。
$\mathcal{M}$ は非輪状なので、各三つ組
$l = (U_l, n_l, x_l) \in L$ に対して被覆
$\{(U_{l, i} \to U_l\}_{i \in I_l} \in \text{Cov}(\mathcal{C})$
と $x_{l, i} \in \mathcal{M}^{n_l - 1}(U_{l, i})$ を
$\text{d}(x_{l, i}) = x|_{U_{l, i}}$ となるように選べる。
そこで
$$
K = \{(U_{l, i}, n_l - 1, x_{l, i}) \mid l \in L, i \in I_l\}
$$
とおき、$\mathcal{F}_{n + 1}$ を $\mathcal{F}_n$ と切断
$x_{l, i}$ を含む $\mathcal{M}$ の最小の微分次数付き
$\mathcal{A}$-部分加群とする。
$|K| \leq \max(\kappa, |\mathcal{F}_n|)$ なので、帰納法により
$|\mathcal{F}_{n + 1}| \leq \kappa$ である。

\medskip\noindent
構成により、包含 $\mathcal{F}_n \to \mathcal{F}_{n + 1}$ が
コホモロジーの層上で誘導する写像は零である。したがって
$\mathcal{F} = \bigcup \mathcal{F}_n$ は
$|\mathcal{F}| \leq \kappa$ を満たす非零で非輪状な部分加群である。
$|\mathcal{F}|$ が有界な微分次数付き
$\mathcal{A}$-加群 $\mathcal{F}$ の同型類は集合をなすので、結論を得る。
\end{proof}

\begin{definition}
\label{sdga-definition-K-injective}
$(\mathcal{C}, \mathcal{O})$ を環付きサイトとする。
$(\mathcal{A}, \text{d})$ を $(\mathcal{C}, \mathcal{O})$ 上の
微分次数付き代数の層とする。微分次数付き
$\mathcal{A}$-加群 $\mathcal{I}$ が{\it K-入射的}であるとは、
任意の非輪状な微分次数付き $\mathcal{M}$ に対して
$$
\Hom_{K(\textit{Mod}(\mathcal{A}, \text{d}))}(\mathcal{M}, \mathcal{I}) = 0
$$
となることをいう。
\end{definition}

\noindent
Derived Categories, Definition \ref{derived-definition-K-injective}
との類似に注意せよ。

\begin{lemma}
\label{sdga-lemma-product-K-injective}
$(\mathcal{C}, \mathcal{O})$ を環付きサイトとする。
$(\mathcal{A}, \text{d})$ を $(\mathcal{C}, \mathcal{O})$ 上の
微分次数付き代数の層とする。$T$ を集合とし、各 $t \in T$ に対して
$\mathcal{I}_t$ を K-入射的な微分次数付き
$\mathcal{A}$-加群とする。このとき $\prod \mathcal{I}_t$ は
K-入射的な微分次数付き $\mathcal{A}$-加群である。
\end{lemma}

\begin{proof}
$\mathcal{K}$ を非輪状な微分次数付き $\mathcal{A}$-加群とする。
このとき
$$
\Hom_{\textit{Mod}^{dg}(\mathcal{A}, \text{d})}(\mathcal{K},
\prod\nolimits_{t \in T} \mathcal{I}_t)
=
\prod\nolimits_{t \in T}
\Hom_{\textit{Mod}^{dg}(\mathcal{A}, \text{d})}(\mathcal{K}, \mathcal{I}_t)
$$
である。これは $\textit{Mod}(\mathcal{A}, \text{d})$ で直積をとることが
次数付き $\mathcal{A}$-加群への忘却関手と可換するからである。
アーベル群の圏上で直積をとることは完全関手なので、結論を得る。
\end{proof}

\begin{lemma}
\label{sdga-lemma-first-property-dg-injective}
$(\mathcal{C}, \mathcal{O})$ を環付きサイトとする。
$(\mathcal{A}, \text{d})$ を $(\mathcal{C}, \mathcal{O})$ 上の
微分次数付き代数の層とする。$\mathcal{I}$ を
$\textit{Mod}(\mathcal{A}, \text{d})$ の K-入射的かつ
次数付き入射的な対象とする。
$\textit{Mod}(\mathcal{A}, \text{d})$ の任意の実線図式
$$
\xymatrix{
\mathcal{M} \ar[r]_a \ar[d]_b & \mathcal{I} \\
\mathcal{M}' \ar@{..>}[ru]
}
$$
で $b$ が単射、$\mathcal{M}$ が非輪状であるものに対し、
図式を可換にする点線矢印が存在する。
\end{lemma}

\begin{proof}
$\mathcal{M}$ は非輪状で $\mathcal{I}$ は K-入射的なので、
$a = \text{d}(h)$ となる次数 $-1$ の次数付き
$\mathcal{A}$-加群写像 $h : \mathcal{M} \to \mathcal{I}$ が存在する。
$\mathcal{I}$ は次数付き入射的で $b$ は単射なので、
$h = h' \circ b$ となる次数 $-1$ の次数付き
$\mathcal{A}$-加群写像 $h' : \mathcal{M}' \to \mathcal{I}$ が存在する。
このとき点線矢印として $a' = \text{d}(h')$ をとればよい。
\end{proof}

\begin{lemma}
\label{sdga-lemma-second-property-dg-injective}
$(\mathcal{C}, \mathcal{O})$ を環付きサイトとする。
$(\mathcal{A}, \text{d})$ を $(\mathcal{C}, \mathcal{O})$ 上の
微分次数付き代数の層とする。$\mathcal{I}$ を
$\textit{Mod}(\mathcal{A}, \text{d})$ の K-入射的かつ
次数付き入射的な対象とする。
$\textit{Mod}(\mathcal{A}, \text{d})$ の任意の実線図式
$$
\xymatrix{
\mathcal{M} \ar[r]_a \ar[d]_b & \mathcal{I} \\
\mathcal{M}' \ar@{..>}[ru]
}
$$
で $b$ が擬同型であるものに対し、図式をホモトピーの差を除いて
可換にする点線矢印が存在する。
\end{lemma}

\begin{proof}
$\mathcal{M}'$ を $\mathcal{M}'$ と $\mathcal{M}$ の恒等写像の錐
（これは非輪状である）との直和で置き換えることで、
$b$ はさらに単射であるとしてよい。このとき $b$ の余核
$\mathcal{Q}$ は非輪状である。ゆえに $\mathcal{I}$ は K-入射的なので
$$
\Hom_{K(\textit{Mod}(\mathcal{A}, \text{d}))}(\mathcal{Q}, \mathcal{I}) =
\Hom_{K(\textit{Mod}(\mathcal{A}, \text{d}))}(\mathcal{Q}, \mathcal{I})[1] = 0
$$
である。$\mathcal{I}$ は次数付き入射的なので、
Remark \ref{sdga-remark-why-graded-injective} により
$$
\Hom_{K(\textit{Mod}(\mathcal{A}, \text{d}))}(\mathcal{M}', \mathcal{I})
\longrightarrow
\Hom_{K(\textit{Mod}(\mathcal{A}, \text{d}))}(\mathcal{M}, \mathcal{I})
$$
は全単射であり、証明は完了する。
\end{proof}


\begin{lemma}
\label{sdga-lemma-better-set-of-monos}
$(\mathcal{C}, \mathcal{O})$ を環付きサイトとする。
$(\mathcal{A}, \text{d})$ を $(\mathcal{C}, \mathcal{O})$ 上の
微分次数付き代数の層とする。集合 $R$ と、各 $r \in R$ に対する
非輪状な微分次数付き $\mathcal{A}$-加群の単射
$\mathcal{M}_r \to \mathcal{M}'_r$ が存在し、
$\textit{Mod}(\mathcal{A}, \text{d})$ の対象 $\mathcal{I}$ に対して
次は同値である。
\begin{enumerate}
\item $\mathcal{I}$ は K-入射的かつ次数付き入射的である。
\item すべての実線図式
$$
\xymatrix{
\mathcal{M}_r \ar[r] \ar[d] & \mathcal{I} \\
\mathcal{M}'_r \ar@{..>}[ru]
}
$$
に対して、図式を可換にする点線矢印が
$\textit{Mod}(\mathcal{A}, \text{d})$ に存在する。
\end{enumerate}
\end{lemma}

\begin{proof}
$T$ と $\mathcal{M}_t \to \mathcal{M}'_t$ を
Lemma \ref{sdga-lemma-characterize-graded-injectives-in-dg} のものとし、
$S$ と $\mathcal{M}_s$ を Lemma \ref{sdga-lemma-small-acyclics} のものとする。
零写像にホモトピックな、非輪状な微分次数付き
$\mathcal{A}$-加群の単射
$\mathcal{M}_s \to \mathcal{M}'_s$ を選ぶ。
これは $\mathcal{M}'_s$ として恒等写像の錐をとれば可能である。
その場合、$\mathcal{M}'_s$ 上の恒等写像さえ零写像にホモトピックである。
Differential Graded Algebra, Lemma \ref{dga-lemma-id-cone-null} を参照せよ。
この補題は Section \ref{sdga-section-conclude-triangulated} の議論により
適用できる。$R = T \coprod S$ と与えられた写像が要件を満たすと主張する。

\medskip\noindent
(1) $\Rightarrow$ (2) は
Lemma \ref{sdga-lemma-first-property-dg-injective} から従う。

\medskip\noindent
(2) を仮定する。まず
Lemma \ref{sdga-lemma-characterize-graded-injectives-in-dg} により
$\mathcal{I}$ は次数付き入射的である。次に
$\mathcal{M}$ を非輪状な微分次数付き $\mathcal{A}$-加群とする。
$$
\Hom_{K(\textit{Mod}(\mathcal{A}, \text{d}))}(\mathcal{M}, \mathcal{I}) = 0
$$
を示さなければならない。証明は Injectives,
Lemma \ref{injectives-lemma-characterize-K-injective} の証明と
まったく同じである。

\medskip\noindent
序数 $\alpha$ に関する帰納法により、非輪状な微分次数付き部分加群
$\mathcal{K}_\alpha \subset \mathcal{M}$ を次のように構成する。
$\alpha = 0$ では $\mathcal{K}_0 = 0$ とおく。
$\alpha > 0$ では次のように進める。
\begin{enumerate}
\item $\alpha = \beta + 1$ かつ
$\mathcal{K}_\beta = \mathcal{M}$ ならば、
$\mathcal{K}_\alpha = \mathcal{K}_\beta$ と選ぶ。
\item $\alpha = \beta + 1$ かつ
$\mathcal{K}_\beta \not = \mathcal{M}$ ならば、
$\mathcal{M}/\mathcal{K}_\beta$ は非零で非輪状な微分次数付き
$\mathcal{A}$-加群である。
Lemma \ref{sdga-lemma-small-acyclics} により、ある $s \in S$ に対する
$\mathcal{M}_s$ と同型な微分次数付き
$\mathcal{A}$-部分加群
$\mathcal{N}_\alpha \subset \mathcal{M}/\mathcal{K}_\beta$ を選ぶ。
最後に $\mathcal{K}_\alpha \subset \mathcal{M}$ を
$\mathcal{N}_\alpha$ の逆像とする。
\item $\alpha$ が極限序数ならば
$\mathcal{K}_\beta = \colim \mathcal{K}_\alpha$ とおく。
\end{enumerate}
十分大きい序数 $\alpha$ に対し
$\mathcal{M} = \mathcal{K}_\alpha$ となることは明らかである。
$$
\Hom_{K(\textit{Mod}(\mathcal{A}, \text{d}))}(\mathcal{K}_\alpha, \mathcal{I})
$$
が零であることを $\alpha$ に関する超限帰納法で示す。
$\mathcal{K}_0$ は零なので $\alpha = 0$ では成り立つ。
$\beta$ に対して成り立ち、$\alpha = \beta + 1$ とする。
上のリストの (1) の場合には結論は明らかである。
(2) の場合には短完全列
$$
0 \to \mathcal{K}_\beta \to \mathcal{K}_\alpha \to \mathcal{N}_\alpha \to 0
$$
がある。Remark \ref{sdga-remark-why-graded-injective} と
$\mathcal{I}$ が次数付き入射的であることから完全列
$$
\Hom_{K(\textit{Mod}(\mathcal{A}, \text{d}))}(\mathcal{K}_\beta, \mathcal{I})
\to
\Hom_{K(\textit{Mod}(\mathcal{A}, \text{d}))}(\mathcal{K}_\alpha, \mathcal{I})
\to
\Hom_{K(\textit{Mod}(\mathcal{A}, \text{d}))}(\mathcal{N}_\alpha, \mathcal{I})
$$
を得る。帰納法により左の項は零である。仮定 (2) により右の項も零である。
実際、任意の写像 $\mathcal{M}_s \to \mathcal{I}$ は
$\mathcal{M}'_s$ を経由し、したがって零写像にホモトピックである。
ゆえに中央の群も零である。最後に $\alpha$ を極限序数とする。
次数付き $\mathcal{A}$-加群として
$\mathcal{K}_\alpha = \colim \mathcal{K}_\alpha$ でもあるので
$$
\Hom_{\textit{Mod}^{dg}(\mathcal{A}, \text{d})}
(\mathcal{K}_\alpha, \mathcal{I})
= \lim_{\beta < \alpha}
\Hom_{\textit{Mod}^{dg}(\mathcal{A}, \text{d})}
(\mathcal{K}_\beta, \mathcal{I})
$$
がアーベル群の複体として成り立つ。
これらの複体のコホモロジー群は
$K(\textit{Mod}(\mathcal{A}, \text{d}))$ におけるシフト間の射を計算する。
$\mathcal{I}$ は次数付き入射的なので、
Remark \ref{sdga-remark-why-graded-injective} により
この複体の系の遷移写像は全射である。さらに極限序数
$\beta \leq \alpha$ に対して、極限とその点での値は等しい。
したがって Homology, Lemma \ref{homology-lemma-ML-over-ordinals}
を適用して結論を得る。
\end{proof}

\begin{lemma}
\label{sdga-lemma-functor-set-of-monos}
$(\mathcal{C}, \mathcal{O})$ を環付きサイトとする。
$(\mathcal{A}, \text{d})$ を $(\mathcal{C}, \mathcal{O})$ 上の
微分次数付き代数の層とする。$R$ を集合とし、各 $r \in R$ に対して
非輪状な微分次数付き $\mathcal{A}$-加群の単射
$\mathcal{M}_r \to \mathcal{M}'_r$ が与えられているとする。
関手
$M : \textit{Mod}(\mathcal{A}, \text{d}) \to
\textit{Mod}(\mathcal{A}, \text{d})$
と自然変換 $j : \text{id} \to M$ で次を満たすものが存在する。
\begin{enumerate}
\item $j_\mathcal{M} : \mathcal{M} \to M(\mathcal{M})$ は
単射かつ擬同型である。
\item すべての実線図式
$$
\xymatrix{
\mathcal{M}_r \ar[r] \ar[d] & \mathcal{M} \ar[d]^{j_\mathcal{M}} \\
\mathcal{M}'_r \ar@{..>}[r] & M(\mathcal{M})
}
$$
に対して、図式を可換にする点線矢印が
$\textit{Mod}(\mathcal{A}, \text{d})$ に存在する。
\end{enumerate}
\end{lemma}

\begin{proof}
$M(\mathcal{M})$ を図式
$$
\xymatrix{
\bigoplus_{(r, \varphi)} \mathcal{M}_r \ar[r] \ar[d] &
\mathcal{M} \ar[d] \\
\bigoplus_{(r, \varphi)} \mathcal{M}'_r \ar[r] &
M(\mathcal{M})
}
$$
の押し出しと定義する。ここで直和は
$r \in R$ および
$\varphi \in
\Hom_{\textit{Mod}(\mathcal{A}, \text{d})}(\mathcal{M}_r, \mathcal{M})$
を満たすすべての対 $(r, \varphi)$ にわたる。
単射の押し出しは単射なので
$\mathcal{M} \to M(\mathcal{M})$ は単射である。
左の縦矢印の余核は非輪状なので、それと同型な
$\mathcal{M} \to M(\mathcal{M})$ の余核も非輪状であり、
したがって $\mathcal{M} \to M(\mathcal{M})$ は擬同型である。
性質 (2) は構成により成り立つ。この手続きを関手にできることの
確認は省略する。
\end{proof}

\begin{theorem}
\label{sdga-theorem-qis-into-dg-injective}
$(\mathcal{C}, \mathcal{O})$ を環付きサイトとする。
$(\mathcal{A}, \text{d})$ を $(\mathcal{C}, \mathcal{O})$ 上の
微分次数付き代数の層とする。任意の微分次数付き
$\mathcal{A}$-加群 $\mathcal{M}$ に対して、
$\mathcal{I}$ が次数付き入射的かつ K-入射的な
微分次数付き $\mathcal{A}$-加群であるような擬同型
$\mathcal{M} \to \mathcal{I}$ が存在する。
さらに、この構成は $\mathcal{M}$ に関して関手的である。
\end{theorem}

\begin{proof}
$R$ と $\mathcal{M}_r \to \mathcal{M}'_r$ を
Lemma \ref{sdga-lemma-better-set-of-monos} で得た
$\textit{Mod}(\mathcal{A}, \text{d})$ の射の集合とする。
$M$ と自然変換 $\text{id} \to M$ を、$R$ と
$\mathcal{M}_r \to \mathcal{M}'_r$ を用いて
Lemma \ref{sdga-lemma-functor-set-of-monos} で構成したものとする。
超限再帰を用い、関手の列 $M_\alpha$ と自然変換
$M_\beta \to M_\alpha$（$\alpha < \beta$）を次のように定義する。
\begin{enumerate}
\item $M_0 = \text{id}$ とする。
\item $M_{\alpha + 1} = M \circ M_\alpha$ とし、
$\beta < \alpha + 1$ に対する自然変換
$M_\beta \to M_{\alpha + 1}$ は、既に構成した
$M_\beta \to M_\alpha$ と $\text{id} \to M$ から来る写像
$M_\alpha \to M \circ M_\alpha$ によって与える。
\item $\alpha$ が極限序数ならば
$M_\alpha = \colim_{\beta < \alpha} M_\beta$ とし、余入射を
$\alpha < \beta$ に対する変換
$M_\beta \to M_\alpha$ とする。
\end{enumerate}
任意の微分次数付き $\mathcal{A}$-加群に対して、写像
$\mathcal{M} \to M_\beta(\mathcal{M}) \to M_\alpha(\mathcal{M})$
は単射な擬同型であることに注意せよ
（フィルター余極限が完全だからである）。

\medskip\noindent
$\textit{Mod}(\mathcal{A}, \text{d})$ は
Grothendieck アーベル圏であった。したがって Injectives,
Proposition \ref{injectives-proposition-objects-are-small}
（すべての $r \in R$ に対する $\mathcal{M}_r$ の直和に適用する）
により、各 $r \in R$ に対して $\mathcal{M}_r$ が単射に関して
$\alpha$-小となる極限序数 $\alpha$ が存在する。
$\mathcal{M} \to M_\alpha(\mathcal{M})$ は、
$\mathcal{M}$ を次数付き入射的かつ K-入射的な加群に埋め込む
所望の関手的な埋め込みであると主張する。

\medskip\noindent
実際、任意の写像
$\mathcal{M}_r \to M_\alpha(\mathcal{M})$ は、ある
$\beta < \alpha$ に対して $M_\beta(\mathcal{M})$ を経由する。
しかし $M$ の構成により、これは
$\mathcal{M}_r \to M_{\beta + 1}(\mathcal{M})
= M(M_\beta(\mathcal{M}))$
が $\mathcal{M}'_r$ を経由することを意味する。
$M_\beta(\mathcal{M}) \subset  M_{\beta + 1}(\mathcal{M})
\subset M_\alpha(\mathcal{M})$ なので、
$M_\alpha(\mathcal{M})$ への所望の分解を得る。
Lemma \ref{sdga-lemma-better-set-of-monos} における
$R$ と $\mathcal{M}_r \to \mathcal{M}'_r$ の選択により結論する。
\end{proof}












\section{導来圏}
\label{sdga-section-derived}

\noindent
本節は Differential Graded Algebra, Section
\ref{dga-section-derived} の類似である。

\medskip\noindent
$(\mathcal{C}, \mathcal{O})$ を環付きサイトとする。
$(\mathcal{A}, \text{d})$ を $(\mathcal{C}, \mathcal{O})$ 上の
微分次数付き代数の層とする。
$K(\textit{Mod}(\mathcal{A}, \text{d}))$ において擬同型を逆にすることにより、
導来圏 $D(\mathcal{A}, \text{d})$ を構成する。

\begin{lemma}
\label{sdga-lemma-cohomology-homological}
$(\mathcal{C}, \mathcal{O})$ を環付きサイトとする。
$(\mathcal{A}, \text{d})$ を $(\mathcal{C}, \mathcal{O})$ 上の
微分次数付き代数の層とする。Section \ref{sdga-section-modules} の関手
$H^0 : \textit{Mod}(\mathcal{A}, \text{d}) \to \textit{Mod}(\mathcal{O})$
は関手
$$
H^0 : K(\textit{Mod}(\mathcal{A}, \text{d})) \to \textit{Mod}(\mathcal{O})
$$
を経由する。この関手は Derived Categories, Definition
\ref{derived-definition-homological} の意味でホモロジー的である。
\end{lemma}

\begin{proof}
定義から直ちに交換図
$$
\xymatrix{
\textit{Mod}(\mathcal{A}, \text{d}) \ar[r] \ar[d] &
K(\textit{Mod}(\mathcal{A}, \text{d})) \ar[d] \\
\text{Comp}(\mathcal{O}) \ar[r] &
K(\textit{Mod}(\mathcal{O}))
}
$$
を得る。$H^0(\mathcal{M})$ は $\mathcal{M}$ の基礎にある
$\mathcal{O}$-加群の複体の第零コホモロジー層として定義されるから、
本補題は $\mathcal{O}$-加群の複体の場合から従う。これは
Derived Categories, Lemma
\ref{derived-lemma-cohomology-homological} の特別な場合である。
\end{proof}

\begin{lemma}
\label{sdga-lemma-acyclics}
$(\mathcal{C}, \mathcal{O})$ を環付きサイトとする。
$(\mathcal{A}, \text{d})$ を $(\mathcal{C}, \mathcal{O})$ 上の
微分次数付き代数の層とする。非輪状加群からなるホモトピー圏
$K(\textit{Mod}(\mathcal{A}, \text{d}))$ の充満部分圏 $\text{Ac}$ は、
$K(\textit{Mod}(\mathcal{A}, \text{d}))$ の厳密充満な
飽和三角部分圏である。
\end{lemma}

\begin{proof}
もちろん、$K(\textit{Mod}(\mathcal{A}, \text{d}))$ の対象
$\mathcal{M}$ が $\text{Ac}$ に属することと、すべての $i$ に対して
$H^i(\mathcal{M}) = H^0(\mathcal{M}[i])$ が零であることとは同値である。
したがって本補題は、Lemma \ref{sdga-lemma-cohomology-homological} および
Derived Categories, Lemma
\ref{derived-lemma-homological-functor-kernel} から従う。
Derived Categories, Definitions \ref{derived-definition-saturated},
\ref{derived-definition-triangulated-subcategory} および
Lemma \ref{derived-lemma-triangulated-subcategory} も参照せよ。
\end{proof}

\begin{lemma}
\label{sdga-lemma-qis}
$(\mathcal{C}, \mathcal{O})$ を環付きサイトとする。
$(\mathcal{A}, \text{d})$ を $(\mathcal{C}, \mathcal{O})$ 上の
微分次数付き代数の層とする。擬同型からなる部分類
$\text{Qis} \subset \text{Arrows}(K(\textit{Mod}(\mathcal{A}, \text{d})))$
を考える。これは飽和乗法的系であり、
$K(\textit{Mod}(\mathcal{A}, \text{d}))$ 上の三角構造と両立する。
\end{lemma}

\begin{proof}
$f , g : \mathcal{M} \to \mathcal{N}$ を
$\textit{Mod}(\mathcal{A}, \text{d})$ の互いにホモトピックな射とする。
このとき、$f$ が擬同型であることと $g$ が擬同型であることとは同値である。
実際、Lemma \ref{sdga-lemma-cohomology-homological} により、写像
$H^i(f) = H^0(f[i])$ と $H^i(g) = H^0(g[i])$ は同じである。
したがって、ホモトピー圏
$K(\textit{Mod}(\mathcal{A}, \text{d}))$ の射が擬同型であると言っても
曖昧さはない。「乗法的系」、「飽和」および
「三角構造と両立する」の定義については
Derived Categories, Definition \ref{derived-definition-localization} および
Categories, Definitions \ref{categories-definition-multiplicative-system},
\ref{categories-definition-saturated-multiplicative-system} を参照せよ。

\medskip\noindent
実際に本補題を証明するため、三角圏の完全関手の合成
$$
K(\textit{Mod}(\mathcal{A}, \text{d}))
\longrightarrow
K(\textit{Mod}(\mathcal{O}))
\longrightarrow
D(\mathcal{O})
$$
を考える。$K(\textit{Mod}(\mathcal{A}, \text{d}))$ の射
$f : \mathcal{M} \to \mathcal{N}$ が $\text{Qis}$ に属することと、
その像が $D(\mathcal{O})$ において同型であることとは同値である。
したがって本補題は Derived Categories, Lemma
\ref{derived-lemma-triangle-functor-localize} から従う。
\end{proof}

\noindent
Lemma \ref{sdga-lemma-qis} の状況では、Derived Categories, Proposition
\ref{derived-proposition-construct-localization} を適用して、
三角圏の完全関手
$$
Q :
K(\textit{Mod}(\mathcal{A}, \text{d}))
\longrightarrow
\text{Qis}^{-1}K(\textit{Mod}(\mathcal{A}, \text{d}))
$$
を得ることができる。しかし $\textit{Mod}(\mathcal{A}, \text{d})$ は
「大きい」圏、すなわちその対象が真の類をなす圏であるため、
$\mathcal{M}$ と $\mathcal{N}$ を与えたとき、
$\text{Qis}^{-1}K(\textit{Mod}(\mathcal{A}, \text{d}))$ の構成が
射の{\it 集合}
$$
\Mor_{\text{Qis}^{-1}K(\textit{Mod}(\mathcal{A}, \text{d}))}
(\mathcal{M}, \mathcal{N})
$$
を生じることは直ちには明らかでない。しかし、K-入射複体の構成により、
これは実際に成り立つ。すなわち Theorem
\ref{sdga-theorem-qis-into-dg-injective} により、擬同型
$s_0 : \mathcal{N} \to \mathcal{I}$ を選べる。ここで $\mathcal{I}$ は
次数付き入射的かつ K-入射的な微分次数付き $\mathcal{A}$-加群である。
上の集合の元は対
$(f : \mathcal{M} \to \mathcal{N}', s : \mathcal{N} \to \mathcal{N}')$
の同値類であることを思い出そう。ここで $f$ は
$K(\textit{Mod}(\mathcal{A}, \text{d}))$ の任意の射であり、
$s$ は擬同型である。Categories, Section
\ref{categories-section-localization} の左分数計算の記述を参照せよ。
Lemma \ref{sdga-lemma-second-property-dg-injective} により、点線の射
$$
\xymatrix{
\mathcal{M} \ar[rd]^f & &
\mathcal{N} \ar[ld]_s \ar[rd]^{s_0} \\
& \mathcal{N}' \ar@{..>}[rr]^{s'} & & \mathcal{I}
}
$$
を選んで、この図を（ホモトピー圏において）可換にできる。
したがって対 $(f, s)$ は対 $(s' \circ f, s_0)$ と同値であり、
同値類の集まりは集合をなすことが分かる。

\begin{definition}
\label{sdga-definition-derived-category}
$(\mathcal{C}, \mathcal{O})$ を環付きサイトとする。
$(\mathcal{A}, \text{d})$ を $(\mathcal{C}, \mathcal{O})$ 上の
微分次数付き代数の層とする。$\text{Qis}$ を
Lemma \ref{sdga-lemma-qis} のものとする。
{\it $(\mathcal{A}, \text{d})$ の導来圏}とは、上で詳しく論じた三角圏
$$
D(\mathcal{A}, \text{d}) =
\text{Qis}^{-1}K(\textit{Mod}(\mathcal{A}, \text{d}))
$$
のことである。
\end{definition}

\noindent
この構成に関するいくつかの事実を証明する。

\begin{lemma}
\label{sdga-lemma-kernel-localization}
Definition \ref{sdga-definition-derived-category} において、局所化関手
$Q : K(\textit{Mod}(\mathcal{A}, \text{d})) \to D(\mathcal{A}, \text{d})$
の核は Lemma \ref{sdga-lemma-acyclics} の圏 $\text{Ac}$ である。
\end{lemma}

\begin{proof}
これは Derived Categories, Lemma
\ref{derived-lemma-kernel-localization} と、
$0 \to \mathcal{M}$ が擬同型であることと $\mathcal{M}$ が
非輪状であることとの同値性から直ちに従う。
\end{proof}

\begin{lemma}
\label{sdga-lemma-H0-over-D}
Definition \ref{sdga-definition-derived-category} において、関手
$H^0 : K(\textit{Mod}(\mathcal{A}, \text{d})) \to
\textit{Mod}(\mathcal{O})$ はホモロジー的関手
$H^0 : D(\mathcal{A}, \text{d}) \to \textit{Mod}(\mathcal{O})$
を経由する。
\end{lemma}

\begin{proof}
Derived Categories, Lemma
\ref{derived-lemma-universal-property-localization} から直ちに従う。
\end{proof}

\noindent
導来圏の射集合を計算する、先に予告した補題を述べる。

\begin{lemma}
\label{sdga-lemma-hom-derived}
$(\mathcal{C}, \mathcal{O})$ を環付きサイトとする。
$(\mathcal{A}, \text{d})$ を $(\mathcal{C}, \mathcal{O})$ 上の
微分次数付き代数の層とする。$\mathcal{M}$ と $\mathcal{N}$ を
微分次数付き $\mathcal{A}$-加群とする。
$\mathcal{N} \to \mathcal{I}$ を擬同型とし、$\mathcal{I}$ は
次数付き入射的かつ K-入射的な微分次数付き $\mathcal{A}$-加群であるとする。
このとき
$$
\Hom_{D(\mathcal{A}, \text{d})}(\mathcal{M}, \mathcal{N}) =
\Hom_{K(\textit{Mod}(\mathcal{A}, \text{d}))}(\mathcal{M}, \mathcal{I})
$$
が成り立つ。
\end{lemma}

\begin{proof}
$\mathcal{N} \to \mathcal{I}$ は擬同型なので、
$$
\Hom_{D(\mathcal{A}, \text{d})}(\mathcal{M}, \mathcal{N}) =
\Hom_{D(\mathcal{A}, \text{d})}(\mathcal{M}, \mathcal{I})
$$
である。Definition \ref{sdga-definition-derived-category} に先立つ議論では、
Lemma \ref{sdga-lemma-second-property-dg-injective} を用いて、
$D(\mathcal{A}, \text{d})$ の任意の射
$\mathcal{M} \to \mathcal{I}$ が
$K(\textit{Mod}(\mathcal{A}, \text{d}))$ の射
$f : \mathcal{M} \to \mathcal{I}$ で表されることを示した。
いま $f, f' :  \mathcal{M} \to \mathcal{I}$ を
$K(\textit{Mod}(\mathcal{A}, \text{d}))$ の二つの射とする。
これらが $D(\mathcal{A}, \text{d})$ において同じ射を定めることと、
$K(\textit{Mod}(\mathcal{A}, \text{d}))$ の擬同型
$g : \mathcal{I} \to \mathcal{K}$ で
$g \circ f = g \circ f'$ を満たすものが存在することとは同値である。
Categories, Lemma
\ref{categories-lemma-equality-morphisms-left-localization} を参照せよ。
一方、Lemma \ref{sdga-lemma-second-property-dg-injective} により、
$K(\textit{Mod}(\mathcal{A}, \text{d}))$ において
$h \circ g = \text{id}_\mathcal{I}$ を満たす写像
$h : \mathcal{K} \to \mathcal{I}$ が存在する。
したがって $g \circ f = g \circ f'$ は $f = f'$ を含意し、
証明が完了する。
\end{proof}

\begin{lemma}
\label{sdga-lemma-derived-products}
$(\mathcal{C}, \mathcal{O})$ を環付きサイトとする。
$(\mathcal{A}, \text{d})$ を $(\mathcal{C}, \mathcal{O})$ 上の
微分次数付き代数の層とする。このとき
\begin{enumerate}
\item $D(\mathcal{A}, \text{d})$ は直和と直積の両方をもつ。
\item 直和は微分次数付き $\mathcal{A}$-加群の直和をとることで得られる。
\item 直積は K-入射的な微分次数付き加群の直積をとることで得られる。
\end{enumerate}
\end{lemma}

\begin{proof}
$\textit{Mod}(\mathcal{A}, \text{d})$ は任意の直和と直積をもつ
アーベル圏であり、これらが
$K(\textit{Mod}(\mathcal{A}, \text{d}))$ の直和と直積を与えることを用いる。
Lemmas \ref{sdga-lemma-dgm-abelian}, \ref{sdga-lemma-homotopy-direct-sums} を参照せよ。

\medskip\noindent
$\mathcal{M}_j$ を微分次数付き $\mathcal{A}$-加群の族とする。
直和 $\mathcal{M} = \bigoplus \mathcal{M}_j$ を
微分次数付き $\mathcal{A}$-加群として考える。
微分次数付き $\mathcal{A}$-加群 $\mathcal{N}$ に対し、擬同型
$\mathcal{N} \to \mathcal{I}$ を選ぶ。ここで $\mathcal{I}$ は
微分次数付き $\mathcal{A}$-加群として次数付き入射的かつ K-入射的である。
Theorem \ref{sdga-theorem-qis-into-dg-injective} を参照せよ。
Lemma \ref{sdga-lemma-hom-derived} を用いると
\begin{align*}
\Hom_{D(\mathcal{A}, \text{d})}(\mathcal{M}, \mathcal{N})
& =
\Hom_{K(\mathcal{A}, \text{d})}(\mathcal{M}, \mathcal{I}) \\
& =
\prod \Hom_{K(\mathcal{A}, \text{d})}(\mathcal{M}_j, \mathcal{I}) \\
& =
\prod \Hom_{D(\mathcal{A}, \text{d})}(\mathcal{M}_j, \mathcal{I})
\end{align*}
を得る。したがって補題の (2) に記したとおり、
$D(A, \text{d})$ に直和が存在する。

\medskip\noindent
$\mathcal{M}_j$ を微分次数付き $\mathcal{A}$-加群の族とする。
各 $j$ に対して擬同型
$\mathcal{M} \to \mathcal{I}_j$ を選ぶ。ここで $\mathcal{I}_j$ は
微分次数付き $\mathcal{A}$-加群として次数付き入射的かつ K-入射的である。
微分次数付き $\mathcal{A}$-加群の直積
$\mathcal{I} = \prod \mathcal{I}_j$ を考える。
Lemmas \ref{sdga-lemma-product-K-injective},
\ref{sdga-lemma-product-graded-injective} により、$\mathcal{I}$ は
微分次数付き $\mathcal{A}$-加群として次数付き入射的かつ K-入射的である。
微分次数付き $\mathcal{A}$-加群 $\mathcal{N}$ に対し、
Lemma \ref{sdga-lemma-hom-derived} を用いると
\begin{align*}
\Hom_{D(\mathcal{A}, \text{d})}(\mathcal{N}, \mathcal{I})
& =
\Hom_{K(\mathcal{A}, \text{d})}(\mathcal{N}, \mathcal{I}) \\
& =
\prod \Hom_{K(\mathcal{A}, \text{d})}(\mathcal{N}, \mathcal{I}_j) \\
& =
\prod \Hom_{D(\mathcal{A}, \text{d})}(\mathcal{N}, \mathcal{M}_j)
\end{align*}
を得る。したがって補題の (3) に記したとおり、
$D(\mathcal{A}, \text{d})$ に直積が存在する。
\end{proof}




\section{標準的な $\delta$-関手}
\label{sdga-section-canonical-delta-functor}

\noindent
$(\mathcal{C}, \mathcal{O})$ を環付きサイトとする。
$(\mathcal{A}, \text{d})$ を $(\mathcal{C}, \mathcal{O})$ 上の
微分次数付き代数の層とする。関手
$\textit{Mod}(\mathcal{A}, \text{d}) \to
K(\textit{Mod}(\mathcal{A}, \text{d}))$
を考える。この関手は一般には $\delta$-関手では{\bf ない}。
しかし、関手
$\textit{Mod}(\mathcal{A}, \text{d}) \to D(A, \text{d})$ は
$\delta$-関手となる。これを見るためには、アーベル圏
$\textit{Mod}(\mathcal{A}, \text{d})$ の短完全列
$$
0 \to \mathcal{K} \xrightarrow{a} \mathcal{L} \xrightarrow{b} \mathcal{M} \to 0
$$
に付随する射 $\delta$ を定義しなければならない。
射 $a$ の錐 $C(a)$ と、その標準射
$i : \mathcal{L} \to C(a)$ および
$p : C(a) \to \mathcal{K}[1]$ を考える。
Definition \ref{sdga-definition-cone} を参照せよ。
微分次数付き $\mathcal{A}$-加群の準同型
$$
q : C(a) \longrightarrow \mathcal{M}
$$
が存在する。これは Differential Graded Algebra, Lemma
\ref{dga-lemma-cone}
（Section \ref{sdga-section-conclude-triangulated} の議論により使用できる）を図
$$
\xymatrix{
\mathcal{K} \ar[r]_a \ar[d] &
\mathcal{L} \ar[d]^b \\
0 \ar[r] &
\mathcal{M}
}
$$
に適用して得られる。写像 $q$ は擬同型である。
例えば、これは $\mathcal{O}$-加群の複体の射の圏で成り立つからである。
Derived Categories, Section
\ref{derived-section-canonical-delta-functor} の議論を参照せよ。
Differential Graded Algebra, Lemma
\ref{dga-lemma-cone-homotopy}
（Section \ref{sdga-section-conclude-triangulated} の議論により使用できる）によれば、
三角形
$$
(\mathcal{K}, \mathcal{L}, C(a), a, i, -p)
$$
は $K(\textit{Mod}(\mathcal{A}, \text{d}))$ の識別三角形である。
局所化関手
$K(\textit{Mod}(\mathcal{A}, \text{d})) \to D(\mathcal{A}, \text{d})$
は完全だから、
$(\mathcal{K}, \mathcal{L}, C(a), a, i, -p)$ は
$D(\mathcal{A}, \text{d})$ の識別三角形である。
$q$ は擬同型なので、$D(\mathcal{A}, \text{d})$ において同型である。
したがって
$$
(\mathcal{K}, \mathcal{L}, \mathcal{M}, a, b, -p \circ q^{-1})
$$
は $D(\mathcal{A}, \text{d})$ の識別三角形である。
これは次の補題を示唆する。

\begin{lemma}
\label{sdga-lemma-derived-canonical-delta-functor}
$(\mathcal{C}, \mathcal{O})$ を環付きサイトとする。
$(\mathcal{A}, \text{d})$ を $(\mathcal{C}, \mathcal{O})$ 上の
微分次数付き代数の層とする。局所化関手
$\textit{Mod}(\mathcal{A}, \text{d}) \to D(\mathcal{A}, \text{d})$
は $\delta$-関手の自然な構造をもち、
$$
\delta_{\mathcal{K} \to \mathcal{L} \to \mathcal{M}} = - p \circ q^{-1}
$$
である。ここで $p$ と $q$ は上で説明したものとする。
\end{lemma}

\begin{proof}
複体の短完全列が与えられるたびに、この選択が識別三角形を与えることは
すでに見た。この構成の関手性を示さなければならない。
Derived Categories, Definition
\ref{derived-definition-delta-functor} を参照せよ。
Differential Graded Algebra, Lemma \ref{dga-lemma-cone}
（Section \ref{sdga-section-conclude-triangulated} の議論により使用できる）
から、少し作業すれば従う。Derived Categories, Lemma
\ref{derived-lemma-derived-canonical-delta-functor} と比較せよ。
\end{proof}

\begin{lemma}
\label{sdga-lemma-homotopy-colimit}
$(\mathcal{C}, \mathcal{O})$ を環付きサイトとする。
$(\mathcal{A}, \text{d})$ を $(\mathcal{C}, \mathcal{O})$ 上の
微分次数付き代数の層とする。$\mathcal{M}_n$ を
微分次数付き $\mathcal{A}$-加群の系とする。
このとき、$D(\mathcal{A}, \text{d})$ における導来余極限
$\text{hocolim} \mathcal{M}_n$ は、微分次数付き加群
$\colim \mathcal{M}_n$ によって表される。
\end{lemma}

\begin{proof}
$\mathcal{M} = \colim \mathcal{M}_n$ とおく。
Derived Categories, Lemma \ref{derived-lemma-compute-colimit}
を基礎にある $\mathcal{O}$-加群の複体に適用すると、
微分次数付き $\mathcal{A}$-加群の完全列
$$
0 \to \bigoplus \mathcal{M}_n \to \bigoplus \mathcal{M}_n \to \mathcal{M} \to 0
$$
を得る。Lemma \ref{sdga-lemma-derived-products} により、これらの直和は
$D(\mathcal{A}, \text{d})$ における直和でもある。
したがって、結果は Derived Categories, Definition
\ref{derived-definition-derived-colimit} における導来余極限の定義と、
複体の短完全列が識別三角形を与えるという事実
（Lemma \ref{sdga-lemma-derived-canonical-delta-functor}）から従う。
\end{proof}




\section{導来引き戻し}
\label{sdga-section-derived-pullback}

\noindent
$(f, f^\sharp) : (\Sh(\mathcal{C}), \mathcal{O}_\mathcal{C})
\to (\Sh(\mathcal{D}), \mathcal{O}_\mathcal{D})$
を環付きトポスの射とする。$\mathcal{A}$ を微分次数付き
$\mathcal{O}_\mathcal{C}$-代数とし、$\mathcal{B}$ を微分次数付き
$\mathcal{O}_\mathcal{D}$-代数とする。微分次数付き
$f^{-1}\mathcal{O}_\mathcal{D}$-代数の写像
$$
\varphi : f^{-1}\mathcal{B} \to \mathcal{A}
$$
が与えられているとする。スカラー制限とスカラー拡大の随伴により、
これは微分次数付き $\mathcal{O}_\mathcal{C}$-代数の写像
$\varphi : f^*\mathcal{B} \to \mathcal{A}$ と同じデータである。
同値な言い方をすれば、$\varphi$ は微分次数付き
$\mathcal{O}_\mathcal{D}$-代数の写像
$$
\varphi : \mathcal{B} \to f_*\mathcal{A}
$$
とみなせる。Remark \ref{sdga-remark-functoriality-dga} を参照せよ。

\medskip\noindent
上記に加えて、$\mathcal{A}'$ を第二の微分次数付き
$\mathcal{O}_\mathcal{C}$-代数とし、$\mathcal{N}$ を微分次数付き
$(\mathcal{A}, \mathcal{A}')$-双加群とする。この状況では関手
$$
\textit{Mod}(\mathcal{B}, \text{d})
\longrightarrow
\textit{Mod}(\mathcal{A}', \text{d}),\quad
\mathcal{M} \longmapsto f^*\mathcal{M} \otimes_{\mathcal{A}} \mathcal{N}
$$
を考えられる。Sections \ref{sdga-section-functoriality-dg},
\ref{sdga-section-dg-bimodules} の議論により、これは微分次数付き圏の関手
$$
\textit{Mod}^{dg}(\mathcal{B}, \text{d})
\longrightarrow
\textit{Mod}^{dg}(\mathcal{A}', \text{d}),\quad
\mathcal{M} \longmapsto f^*\mathcal{M} \otimes_{\mathcal{A}} \mathcal{N}
$$
に拡張される。したがって形式的に、三角圏の完全関手
\begin{equation}
\label{sdga-equation-pullback}
K(\textit{Mod}(\mathcal{B}, \text{d}))
\longrightarrow
K(\textit{Mod}(\mathcal{A}', \text{d})),\quad
\mathcal{M} \longmapsto f^*\mathcal{M} \otimes_{\mathcal{A}} \mathcal{N}
\end{equation}
も得られる。

\begin{lemma}
\label{sdga-lemma-derived-tensor-product}
上の状況で、関手 (\ref{sdga-equation-pullback}) と局所化関手
$K(\textit{Mod}(\mathcal{A}', \text{d})) \to D(\mathcal{A}', \text{d})$
との合成は左導来拡張
$D(\mathcal{B}, \text{d}) \to D(\mathcal{A}', \text{d})$ をもつ。
良い右微分次数付き $\mathcal{B}$-加群 $\mathcal{P}$ におけるその値は
$f^*\mathcal{P} \otimes_\mathcal{A} \mathcal{N}$ である。
\end{lemma}

\begin{proof}
任意の（右）微分次数付き $\mathcal{B}$-加群 $\mathcal{M}$ に対し、
$\mathcal{P}$ が良い微分次数付き $\mathcal{B}$-加群であるような擬同型
$\mathcal{P} \to \mathcal{M}$ が存在することを思い出そう。
Lemma \ref{sdga-lemma-resolve} を参照せよ。したがって
Derived Categories, Lemma
\ref{derived-lemma-find-existence-computes} により、
良い微分次数付き $\mathcal{B}$-加群の擬同型
$\mathcal{P} \to \mathcal{P}'$ が与えられたとき、誘導される写像
$$
f^*\mathcal{P} \otimes_\mathcal{A} \mathcal{N}
\longrightarrow
f^*\mathcal{P}' \otimes_\mathcal{A} \mathcal{N}
$$
が擬同型であることを示せば十分である。
$\mathcal{P} \to \mathcal{P}'$ の錐 $\mathcal{P}''$ は
Lemma \ref{sdga-lemma-good-admissible-ses} により良い微分次数付き
$\mathcal{A}$-加群である。
$K(\textit{Mod}(\mathcal{B}, \text{d}))$ に識別三角形
$$
\mathcal{P} \to \mathcal{P}' \to \mathcal{P}'' \to \mathcal{P}[1]
$$
があるため、$K(\textit{Mod}(\mathcal{A}', \text{d}))$ に識別三角形
$$
f^*\mathcal{P} \otimes_\mathcal{A} \mathcal{N} \to
f^*\mathcal{P}' \otimes_\mathcal{A} \mathcal{N} \to
f^*\mathcal{P}'' \otimes_\mathcal{A} \mathcal{N} \to
f^*\mathcal{P}[1] \otimes_\mathcal{A} \mathcal{N}
$$
を得る。Lemma \ref{sdga-lemma-acyclic-good} により、微分次数付き加群
$f^*\mathcal{P}'' \otimes_\mathcal{A} \mathcal{N}$ は非輪状であり、
証明は完了する。
\end{proof}

\begin{definition}
\label{sdga-definition-pullback}
導来テンソル積と導来引き戻し。
\begin{enumerate}
\item $(\mathcal{C}, \mathcal{O})$ を環付きサイトとする。
$\mathcal{A}$ と $\mathcal{B}$ を微分次数付き $\mathcal{O}$-代数とする。
$\mathcal{N}$ を微分次数付き
$(\mathcal{A}, \mathcal{B})$-双加群とする。
Lemma \ref{sdga-lemma-derived-tensor-product} で構成した関手
$D(\mathcal{A}, \text{d}) \to D(\mathcal{B}, \text{d})$ を
{\it 導来テンソル積}と呼び、
$- \otimes_\mathcal{A}^\mathbf{L} \mathcal{N}$ と書く。
\item $(f, f^\sharp) : (\Sh(\mathcal{C}), \mathcal{O}_\mathcal{C})
\to (\Sh(\mathcal{D}), \mathcal{O}_\mathcal{D})$
を環付きトポスの射とする。$\mathcal{A}$ を微分次数付き
$\mathcal{O}_\mathcal{C}$-代数とし、$\mathcal{B}$ を微分次数付き
$\mathcal{O}_\mathcal{D}$-代数とする。$\varphi : \mathcal{B} \to
f_*\mathcal{A}$ を微分次数付き
$\mathcal{O}_\mathcal{D}$-代数の準同型とする。
Lemma \ref{sdga-lemma-derived-tensor-product} で構成した関手
$D(\mathcal{B}, \text{d}) \to D(\mathcal{A}, \text{d})$ を
{\it 導来引き戻し}と呼び、$Lf^*$ と書く。
\end{enumerate}
\end{definition}

\noindent
この用語により、いくつかの明らかな両立性を表せる。

\begin{lemma}
\label{sdga-lemma-compose-pullback-tensor}
Lemma \ref{sdga-lemma-derived-tensor-product} の関手
$D(\mathcal{B}, \text{d}) \to D(\mathcal{A}', \text{d})$ は
$\mathcal{M} \mapsto
Lf^*\mathcal{M} \otimes_\mathcal{A}^\mathbf{L} \mathcal{N}$ に等しい。
\end{lemma}

\begin{proof}
これらの関手は良い微分次数付き加群で対象を表すことによって計算でき、
$\mathcal{P}$ が良い微分次数付き $\mathcal{B}$-加群ならば
$f^*\mathcal{P}$ は良い微分次数付き $\mathcal{A}$-加群であるから、
直ちに従う。
\end{proof}

\begin{lemma}
\label{sdga-lemma-compose-pullback}
$(f, f^\sharp) : (\Sh(\mathcal{C}), \mathcal{O})
\to (\Sh(\mathcal{C}'), \mathcal{O}')$ および
$(g, g^\sharp) : (\Sh(\mathcal{C}'), \mathcal{O}')
\to (\Sh(\mathcal{C}''), \mathcal{O}'')$
を環付きトポスの射とする。$\mathcal{A}$、$\mathcal{A}'$、
$\mathcal{A}''$ をそれぞれ微分次数付き $\mathcal{O}$-代数、
$\mathcal{O}'$-代数、$\mathcal{O}''$-代数とする。
$\varphi : \mathcal{A}' \to f_*\mathcal{A}$ および
$\varphi' : \mathcal{A}'' \to g_*\mathcal{A}'$ をそれぞれ
微分次数付き $\mathcal{O}'$-代数および
$\mathcal{O}''$-代数の準同型とする。このとき
$L(g \circ f)^* = Lf^* \circ Lg^* :
D(\mathcal{A}'', \text{d}) \to D(\mathcal{A}, \text{d})$
である。
\end{lemma}

\begin{proof}
これらの関手は良い微分次数付き加群で対象を表すことによって計算でき、
$\mathcal{P}$ が良い微分次数付き $\mathcal{A}$-加群である場合に
$f^*\mathcal{P}$ は良い微分次数付き $\mathcal{A}'$-加群であるから、
直ちに従う。
\end{proof}

\noindent
$(\mathcal{C}, \mathcal{O})$ を環付きサイトとする。
$\mathcal{A}$ と $\mathcal{B}$ を微分次数付き $\mathcal{O}$-代数とする。
$\mathcal{N} \to \mathcal{N}'$ を微分次数付き
$(\mathcal{A}, \mathcal{B})$-双加群の準同型とする。
このとき、$D(\mathcal{A}, \text{d})$ の $\mathcal{M}$ に関して関手的な
標準写像
$$
t :
\mathcal{M} \otimes_\mathcal{A}^\mathbf{L} \mathcal{N}
\longrightarrow
\mathcal{M} \otimes_\mathcal{A}^\mathbf{L} \mathcal{N}'
$$
を得る。これらは、三角圏の完全関手
$D(\mathcal{A}, \text{d}) \to D(\mathcal{B}, \text{d})$
の間の自然変換を定める。良い微分次数付き
$\mathcal{A}$-加群 $\mathcal{P}$ における $t$ の値は、明らかな写像
$$
\mathcal{P} \otimes_\mathcal{A}^\mathbf{L} \mathcal{N} =
\mathcal{P} \otimes_\mathcal{A} \mathcal{N}
\longrightarrow
\mathcal{P} \otimes_\mathcal{A} \mathcal{N}' =
\mathcal{P} \otimes_\mathcal{A}^\mathbf{L} \mathcal{N}'
$$
である。

\begin{lemma}
\label{sdga-lemma-tensor-symmetry}
上の状況で、$\mathcal{N} \to \mathcal{N}'$ がコホモロジー層上で
同型ならば、$t$ は関手
$(- \otimes_\mathcal{A}^\mathbf{L} \mathcal{N}) \to
(- \otimes_\mathcal{A}^\mathbf{L} \mathcal{N}')$
の同型である。
\end{lemma}

\begin{proof}
任意の良い微分次数付き $\mathcal{A}$-加群 $\mathcal{P}$ に対し、
$\mathcal{P} \otimes_\mathcal{A} \mathcal{N} \to
\mathcal{P} \otimes_\mathcal{A} \mathcal{N}'$
がコホモロジー層上の同型であることを示せば十分である。
そのため、$\mathcal{N} \to \mathcal{N}'$ を左微分次数付き
$\mathcal{A}$-加群の写像とみなしたときの錐を $\mathcal{N}''$ とする。
Definition \ref{sdga-definition-cone} を参照せよ。
（正確には $\mathcal{N}''$ も双加群であるが、これは必要ない。）
テンソル積構成の関手性（これは微分次数付き圏の関手である）により、
$\mathcal{P} \otimes_\mathcal{A} \mathcal{N}''$ は
$\mathcal{O}$-加群の複体として写像
$\mathcal{P} \otimes_\mathcal{A} \mathcal{N} \to
\mathcal{P} \otimes_\mathcal{A} \mathcal{N}'$
の錐である。したがって、
$\mathcal{P} \otimes_\mathcal{A} \mathcal{N}''$ が非輪状であることを
示せば十分である。これは $\mathcal{P}$ が良く、
$\mathcal{N}''$ が擬同型の錐として非輪状であることから従う。
\end{proof}

\begin{lemma}
\label{sdga-lemma-good-on-other-side}
$(\mathcal{C}, \mathcal{O})$ を環付きサイトとする。
$\mathcal{A}$ と $\mathcal{B}$ を微分次数付き $\mathcal{O}$-代数とする。
$\mathcal{N}$ を微分次数付き
$(\mathcal{A}, \mathcal{B})$-双加群とする。$\mathcal{N}$ が
左微分次数付き $\mathcal{A}$-加群として良いならば、
すべての微分次数付き $\mathcal{A}$-加群 $\mathcal{M}$ に対して
$\mathcal{M} \otimes_\mathcal{A}^\mathbf{L} \mathcal{N} =
\mathcal{M} \otimes_\mathcal{A} \mathcal{N}$ である。
\end{lemma}

\begin{proof}
$\mathcal{P}$ が良い（右）微分次数付き $\mathcal{A}$-加群であるような
擬同型 $\mathcal{P} \to \mathcal{M}$ をとる。
本補題を証明するには、
$\mathcal{P} \otimes_\mathcal{A} \mathcal{N} \to
\mathcal{M} \otimes_\mathcal{A} \mathcal{N}$
が擬同型であることを示さなければならない。
写像 $\mathcal{P} \to \mathcal{M}$ の錐 $C$ は
非輪状な右微分次数付き $\mathcal{A}$-加群である。
$\mathcal{N}$ は左微分次数付き $\mathcal{A}$-加群として良いと仮定したから、
$C \otimes_\mathcal{A} \mathcal{N}$ は非輪状である。
$C \otimes_\mathcal{A} \mathcal{N}$ は $\mathcal{O}$-加群の複体として
写像 $\mathcal{P} \otimes_\mathcal{A} \mathcal{N} \to
\mathcal{M} \otimes_\mathcal{A} \mathcal{N}$ の錐なので、結論を得る。
\end{proof}

\begin{lemma}
\label{sdga-lemma-compose-tensor}
$(\mathcal{C}, \mathcal{O})$ を環付きサイトとする。
$\mathcal{A}$、$\mathcal{A}'$、$\mathcal{A}''$ を
微分次数付き $\mathcal{O}$-代数とする。
$\mathcal{N}$ と $\mathcal{N}'$ をそれぞれ微分次数付き
$(\mathcal{A}, \mathcal{A}')$-双加群および
$(\mathcal{A}', \mathcal{A}'')$-双加群とする。標準写像
$$
\mathcal{N} \otimes_{\mathcal{A}'}^\mathbf{L} \mathcal{N}'
\longrightarrow
\mathcal{N} \otimes_{\mathcal{A}'} \mathcal{N}'
$$
が $D(\mathcal{A}'', \text{d})$ において擬同型であると仮定する。
このとき
$$
(\mathcal{M}
\otimes_\mathcal{A}^\mathbf{L} \mathcal{N})
\otimes_{\mathcal{A}'}^\mathbf{L} \mathcal{N}'
=
\mathcal{M}
\otimes_\mathcal{A}^\mathbf{L}
(\mathcal{N} \otimes_{\mathcal{A}'} \mathcal{N}')
$$
が関手 $D(\mathcal{A}, \text{d}) \to D(\mathcal{A}'', \text{d})$
として成り立つ。
\end{lemma}

\begin{proof}
良い微分次数付き $\mathcal{A}$-加群 $\mathcal{P}$ と擬同型
$\mathcal{P} \to \mathcal{M}$ を選ぶ。
Lemma \ref{sdga-lemma-resolve} を参照せよ。このとき
$$
\mathcal{M}
\otimes_\mathcal{A}^\mathbf{L}
(\mathcal{N} \otimes_{\mathcal{A}'} \mathcal{N}') =
\mathcal{P} \otimes_\mathcal{A} \mathcal{N}
\otimes_{\mathcal{A}'} \mathcal{N}'
$$
であり、
$$
(\mathcal{M}
\otimes_\mathcal{A}^\mathbf{L} \mathcal{N})
\otimes_{\mathcal{A}'}^\mathbf{L} \mathcal{N}' =
(\mathcal{P} \otimes_\mathcal{A} \mathcal{N})
\otimes_{\mathcal{A}'}^\mathbf{L} \mathcal{N}'
$$
である。したがって、標準写像
$$
(\mathcal{P} \otimes_\mathcal{A} \mathcal{N})
\otimes_{\mathcal{A}'}^\mathbf{L} \mathcal{N}'
\longrightarrow
\mathcal{P} \otimes_\mathcal{A} \mathcal{N}
\otimes_{\mathcal{A}'} \mathcal{N}'
$$
が擬同型であることを示さなければならない。
$\mathcal{Q}$ が良い左微分次数付き $\mathcal{A}'$-加群であるような擬同型
$\mathcal{Q} \to \mathcal{N}'$ を選ぶ
（Lemma \ref{sdga-lemma-resolve}）。
Lemma \ref{sdga-lemma-good-on-other-side} により、上の写像は
$\mathcal{O}$-加群の導来圏の写像として
$$
\mathcal{P} \otimes_\mathcal{A} \mathcal{N}
\otimes_{\mathcal{A}'} \mathcal{Q}
\longrightarrow
\mathcal{P} \otimes_\mathcal{A} \mathcal{N}
\otimes_{\mathcal{A}'} \mathcal{N}'
$$
である。仮定により
$\mathcal{N} \otimes_{\mathcal{A}'} \mathcal{Q} \to
\mathcal{N} \otimes_{\mathcal{A}'} \mathcal{N}'$ は擬同型であり、
$\mathcal{P}$ は良い微分次数付き $\mathcal{A}$-加群なので、
この写像は Lemma \ref{sdga-lemma-tensor-symmetry} により擬同型である
（左辺と右辺は
$\mathcal{P} \otimes_\mathcal{A}^\mathbf{L} (\mathcal{N}
\otimes_{\mathcal{A}'} \mathcal{Q})$ および
$\mathcal{P} \otimes_\mathcal{A}^\mathbf{L} (\mathcal{N}
\otimes_{\mathcal{A}'} \mathcal{N}')$ を計算している。
または、本補題の証明の議論を繰り返してもよい）。
\end{proof}




\section{導来順像}
\label{sdga-section-derived-pushforward}

\noindent
十分多くの K-入射対象が存在するため、ホモトピー圏上の任意の完全関手の
右導来関手をとることができる。

\begin{lemma}
\label{sdga-lemma-right-derived}
$(\mathcal{C}, \mathcal{O})$ を環付きサイトとする。
$(\mathcal{A}, \text{d})$ を $(\mathcal{C}, \mathcal{O})$ 上の
微分次数付き代数の層とする。このとき、三角圏の任意の完全関手
$$
T : K(\textit{Mod}(\mathcal{A}, \text{d})) \longrightarrow \mathcal{D}
$$
は右導来拡張
$RT : D(\mathcal{A}, \text{d}) \to \mathcal{D}$ をもつ。
次数付き入射的かつ K-入射的な微分次数付き
$\mathcal{A}$-加群 $\mathcal{I}$ におけるその値は
$T(\mathcal{I})$ である。
\end{lemma}

\begin{proof}
Theorem \ref{sdga-theorem-qis-into-dg-injective} により、
任意の（右）微分次数付き $\mathcal{A}$-加群 $\mathcal{M}$ に対し、
$\mathcal{I}$ が次数付き入射的かつ K-入射的な微分次数付き
$\mathcal{A}$-加群であるような擬同型
$\mathcal{M} \to \mathcal{I}$ が存在する。
したがって Derived Categories, Lemma
\ref{derived-lemma-find-existence-computes} により、
ともに次数付き入射的かつ K-入射的な微分次数付き
$\mathcal{A}$-加群の擬同型
$\mathcal{I} \to \mathcal{I}'$ が与えられたとき、
$T(\mathcal{I}) \to T(\mathcal{I}')$ が同型であることを示せば十分である。
これは、$\mathcal{I} \to \mathcal{I}'$ が
$K(\textit{Mod}(\mathcal{A}, \text{d}))$ において同型だから成り立つ。
例えば Lemma \ref{sdga-lemma-hom-derived} から従う
（または Lemma \ref{sdga-lemma-second-property-dg-injective} から導ける）。
\end{proof}

\noindent
これを適用できる関手をすでにいくつか見てきた。
二つの例を挙げる。

\begin{definition}
\label{sdga-definition-pushforward}
導来内部 Hom と導来順像。
\begin{enumerate}
\item $(\mathcal{C}, \mathcal{O})$ を環付きサイトとする。
$\mathcal{A}$ と $\mathcal{B}$ を微分次数付き $\mathcal{O}$-代数とする。
$\mathcal{N}$ を微分次数付き
$(\mathcal{A}, \mathcal{B})$-双加群とする。内部 Hom 関手
$\SheafHom_\mathcal{B}^{dg}(\mathcal{N}, -)$ の右導来拡張
$$
R\SheafHom_\mathcal{B}(\mathcal{N}, -) :
D(\mathcal{B}, \text{d})
\longrightarrow
D(\mathcal{A}, \text{d})
$$
を{\it 導来内部 Hom} と呼ぶ。
\item $(f, f^\sharp) : (\Sh(\mathcal{C}), \mathcal{O}_\mathcal{C})
\to (\Sh(\mathcal{D}), \mathcal{O}_\mathcal{D})$
を環付きトポスの射とする。$\mathcal{A}$ を微分次数付き
$\mathcal{O}_\mathcal{C}$-代数とし、$\mathcal{B}$ を微分次数付き
$\mathcal{O}_\mathcal{D}$-代数とする。
$\varphi : \mathcal{B} \to f_*\mathcal{A}$ を微分次数付き
$\mathcal{O}_\mathcal{D}$-代数の準同型とする。順像 $f_*$ の右導来拡張
$$
Rf_* :
D(\mathcal{A}, \text{d})
\longrightarrow
D(\mathcal{B}, \text{d})
$$
を{\it 導来順像}と呼ぶ。
\end{enumerate}
\end{definition}

\noindent
実際、
$Rf_* : D(\mathcal{A}, \text{d}) \to D(\mathcal{B}, \text{d})$
は基礎にある $\mathcal{O}$-加群の複体の導来順像と一致する。
Lemma \ref{sdga-lemma-pushforward-agrees} を参照せよ。

\medskip\noindent
これらの関手は導来引き戻しおよび導来テンソル積の随伴である。

\begin{lemma}
\label{sdga-lemma-derived-adjoint-tensor-hom}
$(\mathcal{C}, \mathcal{O})$ を環付きサイトとする。
$\mathcal{A}$ と $\mathcal{B}$ を微分次数付き $\mathcal{O}$-代数とする。
$\mathcal{N}$ を微分次数付き
$(\mathcal{A}, \mathcal{B})$-双加群とする。このとき
$$
R\SheafHom_\mathcal{B}(\mathcal{N}, -) :
D(\mathcal{B}, \text{d})
\longrightarrow
D(\mathcal{A}, \text{d})
$$
は
$$
- \otimes_\mathcal{A}^\mathbf{L} \mathcal{N} :
D(\mathcal{A}, \text{d})
\longrightarrow
D(\mathcal{B}, \text{d})
$$
の右随伴である。
\end{lemma}

\begin{proof}
Derived Categories, Lemma
\ref{derived-lemma-pre-derived-adjoint-functors-general} および
Lemma \ref{sdga-lemma-tensor-hom-adjunction-dg} から従う。
\end{proof}

\begin{lemma}
\label{sdga-lemma-derived-adjoint-push-pull}
$(f, f^\sharp) : (\Sh(\mathcal{C}), \mathcal{O}_\mathcal{C})
\to (\Sh(\mathcal{D}), \mathcal{O}_\mathcal{D})$
を環付きトポスの射とする。$\mathcal{A}$ を微分次数付き
$\mathcal{O}_\mathcal{C}$-代数とし、$\mathcal{B}$ を微分次数付き
$\mathcal{O}_\mathcal{D}$-代数とする。
$\varphi : \mathcal{B} \to f_*\mathcal{A}$ を微分次数付き
$\mathcal{O}_\mathcal{D}$-代数の準同型とする。このとき
$$
Rf_* : 
D(\mathcal{A}, \text{d})
\longrightarrow
D(\mathcal{B}, \text{d})
$$
は
$$
Lf^* :
D(\mathcal{B}, \text{d})
\longrightarrow
D(\mathcal{A}, \text{d})
$$
の右随伴である。
\end{lemma}

\begin{proof}
Derived Categories, Lemma
\ref{derived-lemma-pre-derived-adjoint-functors-general} および
Lemma \ref{sdga-lemma-adjunction-push-pull-dg} から従う。
\end{proof}

\noindent
次に、Section \ref{sdga-section-derived-pullback} で考えた状況を論じる。

\medskip\noindent
$(f, f^\sharp) : (\Sh(\mathcal{C}), \mathcal{O}_\mathcal{C})
\to (\Sh(\mathcal{D}), \mathcal{O}_\mathcal{D})$
を環付きトポスの射とする。$\mathcal{A}$ を微分次数付き
$\mathcal{O}_\mathcal{C}$-代数とし、$\mathcal{B}$ を微分次数付き
$\mathcal{O}_\mathcal{D}$-代数とする。微分次数付き
$f^{-1}\mathcal{O}_\mathcal{D}$-代数の写像
$$
\varphi : f^{-1}\mathcal{B} \to \mathcal{A}
$$
が与えられているとする。スカラー制限とスカラー拡大の随伴により、
これは微分次数付き $\mathcal{O}_\mathcal{C}$-代数の写像
$\varphi : f^*\mathcal{B} \to \mathcal{A}$ と同じデータである。
同値な言い方をすれば、$\varphi$ は微分次数付き
$\mathcal{O}_\mathcal{D}$-代数の写像
$$
\varphi : \mathcal{B} \to f_*\mathcal{A}
$$
とみなせる。Remark \ref{sdga-remark-functoriality-dga} を参照せよ。

\medskip\noindent
上記に加えて、$\mathcal{A}'$ を第二の微分次数付き
$\mathcal{O}_\mathcal{C}$-代数とし、$\mathcal{N}$ を微分次数付き
$(\mathcal{A}, \mathcal{A}')$-双加群とする。この状況では関手
$$
\textit{Mod}(\mathcal{A}', \text{d})
\longrightarrow
\textit{Mod}(\mathcal{B}, \text{d}),\quad
\mathcal{M} \longmapsto
f_*\SheafHom^{dg}_{\mathcal{A}'}(\mathcal{N}, \mathcal{M})
$$
を考えられる。Sections \ref{sdga-section-functoriality-dg},
\ref{sdga-section-dg-bimodules} の議論により、これは微分次数付き圏の関手
$$
\textit{Mod}^{dg}(\mathcal{A}', \text{d})
\longrightarrow
\textit{Mod}^{dg}(\mathcal{B}, \text{d}),\quad
\mathcal{M} \longmapsto
f_*\SheafHom^{dg}_{\mathcal{A}'}(\mathcal{N}, \mathcal{M})
$$
に拡張される。したがって形式的に、三角圏の完全関手
\begin{equation}
\label{sdga-equation-pushforward}
K(\textit{Mod}(\mathcal{A}', \text{d}))
\longrightarrow
K(\textit{Mod}(\mathcal{B}, \text{d})),\quad
\mathcal{M} \longmapsto
f_*\SheafHom^{dg}_{\mathcal{A}'}(\mathcal{N}, \mathcal{M})
\end{equation}
も得られる。

\begin{lemma}
\label{sdga-lemma-compose-pushforward-hom}
上の状況で、(\ref{sdga-equation-pushforward}) の右導来拡張を
$RT : D(\mathcal{A}', \text{d}) \to
D(\mathcal{B}, \text{d})$ と書く。このとき
$$
RT(\mathcal{M}) = Rf_* R\SheafHom(\mathcal{N}, \mathcal{M})
$$
が $\mathcal{M}$ に関して関手的に成り立つ。
\end{lemma}

\begin{proof}
Lemmas \ref{sdga-lemma-tensor-hom-adjunction-dg},
\ref{sdga-lemma-adjunction-push-pull-dg} により、関手
(\ref{sdga-equation-pushforward}) は関手
(\ref{sdga-equation-pullback}) の右随伴である。
Derived Categories, Lemma
\ref{derived-lemma-pre-derived-adjoint-functors-general} により、
関手 $RT$ は Lemma \ref{sdga-lemma-derived-tensor-product} の関手の
右随伴であり、Lemma \ref{sdga-lemma-compose-pullback-tensor} により
この関手は $Lf^*(-) \otimes_\mathcal{A}^\mathbf{L} \mathcal{N}$
に等しい。Lemmas \ref{sdga-lemma-derived-adjoint-tensor-hom},
\ref{sdga-lemma-derived-adjoint-push-pull} により、関手
$Lf^*(-) \otimes_\mathcal{A}^\mathbf{L} \mathcal{N}$ は
$Rf_* R\SheafHom(\mathcal{N}, -)$ の左随伴である。
したがって随伴の一意性から結論を得る。
\end{proof}

\begin{lemma}
\label{sdga-lemma-compose-pushforward}
$(f, f^\sharp) : (\Sh(\mathcal{C}), \mathcal{O})
\to (\Sh(\mathcal{C}'), \mathcal{O}')$ および
$(g, g^\sharp) : (\Sh(\mathcal{C}'), \mathcal{O}')
\to (\Sh(\mathcal{C}''), \mathcal{O}'')$
を環付きトポスの射とする。$\mathcal{A}$、$\mathcal{A}'$、
$\mathcal{A}''$ をそれぞれ微分次数付き $\mathcal{O}$-代数、
$\mathcal{O}'$-代数、$\mathcal{O}''$-代数とする。
$\varphi : \mathcal{A}' \to f_*\mathcal{A}$ および
$\varphi' : \mathcal{A}'' \to g_*\mathcal{A}'$ をそれぞれ
微分次数付き $\mathcal{O}'$-代数および
$\mathcal{O}''$-代数の準同型とする。このとき
$R(g \circ f)_* = Rg_* \circ Rf_* :
D(\mathcal{A}, \text{d}) \to D(\mathcal{A}'', \text{d})$
である。
\end{lemma}

\begin{proof}
Lemmas \ref{sdga-lemma-compose-pullback},
\ref{sdga-lemma-derived-adjoint-push-pull} と随伴の一意性から従う。
\end{proof}

\begin{lemma}
\label{sdga-lemma-compose-hom}
$(\mathcal{C}, \mathcal{O})$ を環付きサイトとする。
$\mathcal{A}$、$\mathcal{A}'$、$\mathcal{A}''$ を
微分次数付き $\mathcal{O}$-代数とする。
$\mathcal{N}$ と $\mathcal{N}'$ をそれぞれ微分次数付き
$(\mathcal{A}, \mathcal{A}')$-双加群および
$(\mathcal{A}', \mathcal{A}'')$-双加群とする。標準写像
$$
\mathcal{N} \otimes_{\mathcal{A}'}^\mathbf{L} \mathcal{N}'
\longrightarrow
\mathcal{N} \otimes_{\mathcal{A}'} \mathcal{N}'
$$
が $D(\mathcal{A}'', \text{d})$ において擬同型であると仮定する。
このとき
$$
R\SheafHom_{\mathcal{A}''}
(\mathcal{N} \otimes_{\mathcal{A}'} \mathcal{N}', -)
=
R\SheafHom_{\mathcal{A}'}(\mathcal{N},
R\SheafHom_{\mathcal{A}''}(\mathcal{N}', -))
$$
が関手 $D(\mathcal{A}'', \text{d}) \to D(\mathcal{A}, \text{d})$
として成り立つ。
\end{lemma}

\begin{proof}
Lemmas \ref{sdga-lemma-compose-tensor},
\ref{sdga-lemma-derived-adjoint-tensor-hom} と随伴の一意性から従う。
\end{proof}

\begin{lemma}
\label{sdga-lemma-pushforward-agrees}
$(f, f^\sharp) : (\Sh(\mathcal{C}), \mathcal{O}_\mathcal{C})
\to (\Sh(\mathcal{D}), \mathcal{O}_\mathcal{D})$
を環付きトポスの射とする。$\mathcal{A}$ を微分次数付き
$\mathcal{O}_\mathcal{C}$-代数とし、$\mathcal{B}$ を微分次数付き
$\mathcal{O}_\mathcal{D}$-代数とする。
$\varphi : \mathcal{B} \to f_*\mathcal{A}$ を微分次数付き
$\mathcal{O}_\mathcal{D}$-代数の準同型とする。図
$$
\xymatrix{
D(\mathcal{A}, \text{d}) \ar[d]_{Rf_*} \ar[rr]_{forget} & &
D(\mathcal{O}_\mathcal{C}) \ar[d]^{Rf_*} \\
D(\mathcal{B}, \text{d}) \ar[rr]^{forget} & &
D(\mathcal{O}_\mathcal{D})
}
$$
は可換である。
\end{lemma}

\begin{proof}
いくつかの圏を同一視すれば、本補題は
Lemma \ref{sdga-lemma-compose-pushforward} から直ちに従う。

\medskip\noindent
$\mathcal{O}_\mathcal{C}$ を次数 $0$ に置き、零微分を与えることで、
$(\mathcal{O}_\mathcal{C}, 0)$ を微分次数付き
$\mathcal{O}_\mathcal{C}$-代数とみなせる。明らかに
$$
\textit{Mod}(\mathcal{O}_\mathcal{C}, 0) =
\text{Comp}(\mathcal{O}_\mathcal{C})
\quad\text{and}\quad
D(\mathcal{O}_\mathcal{C}, 0) = D(\mathcal{O}_\mathcal{C})
$$
である。この同一視のもとで、忘却関手
$\textit{Mod}(\mathcal{A}, \text{d}) \to
\text{Comp}(\mathcal{O}_\mathcal{C})$
は Section \ref{sdga-section-functoriality-dg} で定義した「順像」
$\text{id}_{\mathcal{C}, *}$ である。これは環付きトポスの恒等射
$\text{id}_\mathcal{C} : (\mathcal{C}, \mathcal{O}_\mathcal{C}) \to
(\mathcal{C}, \mathcal{O}_\mathcal{C})$
および微分次数付き $\mathcal{O}_\mathcal{C}$-代数の写像
$(\mathcal{O}_\mathcal{C}, 0) \to (\mathcal{A}, \text{d})$
に対応する。$\text{id}_{\mathcal{C}, *}$ は完全なので、直ちに
$$
R\text{id}_{\mathcal{C}, *} = forget :
D(\mathcal{A}, \text{d}) \longrightarrow
D(\mathcal{O}_\mathcal{C}, 0) = D(\mathcal{O}_\mathcal{C})
$$
を得る。まったく同じ議論により
$$
R\text{id}_{\mathcal{D}, *} = forget :
D(\mathcal{B}, \text{d}) \longrightarrow
D(\mathcal{O}_\mathcal{D}, 0) = D(\mathcal{O}_\mathcal{D})
$$
を得る。さらに、Cohomology on Sites, Section
\ref{sites-cohomology-section-unbounded} における
$Rf_* : D(\mathcal{O}_\mathcal{C}) \to D(\mathcal{O}_\mathcal{D})$
の構成は、Definition \ref{sdga-definition-pushforward} における
$Rf_* : D(\mathcal{O}_\mathcal{C}, 0) \to D(\mathcal{O}_\mathcal{D}, 0)$
の構成と一致する。どちらの関手も、基礎にある加群の複体上の順像の
右導来拡張として定義されるからである。
Lemma \ref{sdga-lemma-compose-pushforward} により、
$Rf_* \circ R\text{id}_{\mathcal{C}, *}$ と
$R\text{id}_{\mathcal{D}, *} \circ Rf_*$ はともに
$f_* \circ forget = forget \circ f_*$ の導来関手であり、
したがって随伴の一意性により等しい。
\end{proof}

\begin{lemma}
\label{sdga-lemma-cohomology-ext}
$(\mathcal{C}, \mathcal{O})$ を環付きサイトとする。
$\mathcal{A}$ を微分次数付き $\mathcal{O}$-代数とする。
$\mathcal{M}$ を微分次数付き $\mathcal{A}$-加群とする。
$n \in \mathbf{Z}$ とする。このとき
$$
H^n(\mathcal{C}, \mathcal{M}) =
\Hom_{D(\mathcal{A}, \text{d})}(\mathcal{A}, \mathcal{M}[n])
$$
である。左辺では $\mathcal{M}$ を $\mathcal{O}$-加群の複体とみなした
コホモロジーをとる。
\end{lemma}

\begin{proof}
この公式を証明するため、
$$
R\Gamma(\mathcal{C}, \mathcal{M}) = \Gamma(\mathcal{C}, \mathcal{I})
$$
であることに注意する。ここで $\mathcal{M} \to \mathcal{I}$ は
次数付き入射的かつ K-入射的な微分次数付き
$\mathcal{A}$-加群 $\mathcal{I}$ への擬同型である
（Lemmas \ref{sdga-lemma-right-derived},
\ref{sdga-lemma-pushforward-agrees} を組み合わせよ）。
Lemma \ref{sdga-lemma-hom-derived} により
$$
\Hom_{D(\mathcal{A}, \text{d})}(\mathcal{A}, \mathcal{M}[n]) =
\Hom_{K(\textit{Mod}(\mathcal{A}, \text{d}))}(\mathcal{M}, \mathcal{I}[n]) =
H^0(\Gamma(\mathcal{C}, \mathcal{I}[n])) =
H^n(\Gamma(\mathcal{C}, \mathcal{I}))
$$
である。この二つの結果を合わせると等式を得る。
\end{proof}






\section{導来圏の同値}
\label{sdga-section-equivalence}

\noindent
本節は Differential Graded Algebra, Section
\ref{dga-section-equivalence} の類似である。

\begin{lemma}
\label{sdga-lemma-qis-equivalence}
$(\mathcal{C}, \mathcal{O})$ を環付きサイトとする。
$\varphi : \mathcal{A} \to \mathcal{B}$ を、コホモロジー層上の同型を
誘導する微分次数付き $\mathcal{O}$-代数の準同型とする。このとき
$$
D(\mathcal{A}, \text{d}) \longrightarrow D(\mathcal{B}, \text{d}), \quad
\mathcal{M}
\longmapsto
\mathcal{M} \otimes_\mathcal{A}^\mathbf{L} \mathcal{B}
$$
は圏の同値である。
\end{lemma}

\begin{proof}
制限関手
$$
\textit{Mod}^{dg}(\mathcal{B}, \text{d}) \to
\textit{Mod}^{dg}(\mathcal{A}, \text{d}),\quad
\mathcal{N} \mapsto res_\varphi \mathcal{N}
$$
が関手
$$
\textit{Mod}^{dg}(\mathcal{A}, \text{d}) \to
\textit{Mod}^{dg}(\mathcal{B}, \text{d}),\quad
\mathcal{M} \mapsto \mathcal{M} \otimes_\mathcal{A} \mathcal{B}
$$
の右随伴であることを思い出そう。
Section \ref{sdga-section-dg-bimodules} を参照せよ。
制限は擬同型を擬同型に写すので、（制限自身で与えられる）
左導来拡張を自明にもつ。Derived Categories, Lemma
\ref{derived-lemma-pre-derived-adjoint-functors-general} により、
この関手は $- \otimes_\mathcal{A}^\mathbf{L} \mathcal{B}$ の右随伴である。
随伴写像
$$
\mathcal{M} \to
res_\varphi(\mathcal{M} \otimes_\mathcal{A}^\mathbf{L} \mathcal{B})
$$
は、$\mathcal{A} \to \mathcal{B}$ が（左）微分次数付き
$\mathcal{A}$-加群の擬同型であるという仮定により、
$D(\mathcal{A}, \text{d})$ において同型である。
特に本補題の関手は充満忠実である。
Categories, Lemma \ref{categories-lemma-adjoint-fully-faithful} を参照せよ。
制限関手 $D(\mathcal{B}, \text{d}) \to D(\mathcal{A}, \text{d})$
の核が零であることは明らかである。したがって
Derived Categories, Lemma
\ref{derived-lemma-fully-faithful-adjoint-kernel-zero} から結論を得る。
\end{proof}











\section{微分次数付き代数の分解}
\label{sdga-section-resolution-dgas}

\noindent
本節は Differential Graded Algebra, Section
\ref{dga-section-resolution-dgas} の類似である。

\medskip\noindent
$(\mathcal{C}, \mathcal{O})$ を環付きサイトとする。Remark
\ref{sdga-remark-sheaf-graded-sets} と同様に、$\mathcal{C}$ 上の
次数付き集合の層 $\mathcal{S}$ を考える。
$r$ 重自己積
$\mathcal{S} \times \ldots \times \mathcal{S}$
を、規則
$\deg(s_1 \cdot \ldots \cdot s_r) = \sum \deg(s_i)$
をもつ次数付き集合の層とみなす。
ここで局所切断 $s_i \in \mathcal{S}(U)$、$i = 1, \ldots, r$
が与えられたとき、$s_1 \cdot \ldots \cdot s_r$ は
$\mathcal{S} \times \ldots \times \mathcal{S}$ の
$U$ 上の対応する切断を表すものとする。
$\mathcal{S}$ 上の自由次数付き $\mathcal{O}$-代数を
$\mathcal{O}\langle \mathcal{S} \rangle$ と書く。
より正確には、Remark \ref{sdga-remark-sheaf-graded-sets} の記法を用いて
$$
\mathcal{O}\langle \mathcal{S} \rangle =
\mathcal{O} \oplus
\bigoplus\nolimits_{r \geq 1}
\mathcal{O}[\mathcal{S} \times \ldots \times \mathcal{S}]
$$
とおく。これは連結
$$
(s_1 \cdot \ldots \cdot s_r)  (s'_1 \cdot \ldots \cdot s'_{r'}) =
s_1 \cdot \ldots s_r \cdot s'_1 \cdot \ldots \cdot s'_{r'}
$$
によって次数付き $\mathcal{O}$-代数の層となる。
$\mathcal{S}$ のすべての局所切断 $s$ に対して
$\text{d}(s) = 0$ とおき、Leibniz 則を用いて一意に拡張することにより、
$\mathcal{O}\langle \mathcal{S} \rangle$ に微分を与えられる。
ただし、他の微分も考えることが重要である。

\medskip\noindent
実際、次のような系が与えられたとする。
\begin{enumerate}
\item $i = 0, 1, 2, \ldots$ に対して次数付き集合の層
$\mathcal{S}_i$、
\item $i = 0, 1, 2, \ldots$ に対して次数 $1$ の
次数付き集合の層の写像
$$
\delta_{i + 1} : \mathcal{S}_{i + 1}
\longrightarrow
\mathcal{A}_i =
\mathcal{O}\langle \mathcal{S}_0 \amalg \ldots \amalg \mathcal{S}_i\rangle
$$
で、その像が標的上に帰納的に定義される微分の核に含まれるもの。
\end{enumerate}
より正確には、まず
$\mathcal{A}_0 = \mathcal{O}\langle \mathcal{S}_0 \rangle$
とおき、Leibniz 則を満たし、$\mathcal{S}$ の任意の局所切断 $s$ に対して
$\text{d}(s) = 0$ となる一意な微分を与える。
帰納的に、$\mathcal{A}_i$ 上に微分 $\text{d}$ が与えられたとする。
この微分を、Leibniz 則を満たし、
$$
\text{d}(s) = \delta(s)
$$
となる $\mathcal{A}_{i + 1}$ 上の一意な微分へ拡張する。ここで
$s$ が
$\mathcal{S}_0 \amalg \ldots \amalg \mathcal{S}_{i + 1}$
の直和因子 $\mathcal{S}_j$ に属するとき
$\delta(s) = \delta_j(s)$ とする。
$\delta(s)$ が帰納的に定義された微分の核に属するからこそ、
この定義は意味をもつ。

\begin{lemma}
\label{sdga-lemma-special-good}
上の状況における微分次数付き $\mathcal{O}$-代数
$$
\mathcal{A} = \colim \mathcal{A}_i
$$
は次の性質をもつ。環付きトポスの任意の射
$(f, f^\sharp) : (\Sh(\mathcal{C}'), \mathcal{O}')
\to (\Sh(\mathcal{C}), \mathcal{O})$
に対し、引き戻し $f^*\mathcal{A}$ は次数付き
$\mathcal{O}'$-加群として平坦であり、$\mathcal{O}'$-加群の
複体として K-平坦である。
\end{lemma}

\begin{proof}
$f^*\mathcal{A} = \colim f^*\mathcal{A}_i$ であり、
$$
f^*\mathcal{A}_i = \mathcal{O}'\langle
f^{-1}\mathcal{S}_0 \amalg \ldots \amalg f^{-1}\mathcal{S}_i\rangle
$$
であって、その微分は $f^{-1}\delta_{i + 1}$ を用いる上の帰納的手続きで
与えられることに注意する。したがって、$\mathcal{A}$ が次数付き
$\mathcal{O}$-加群として平坦であり、$\mathcal{O}$-加群の
複体として K-平坦であることを証明すれば十分である。
そのためには、各 $\mathcal{A}_i$ が次数付き
$\mathcal{O}$-加群として平坦であり、$\mathcal{O}$-加群の
複体として K-平坦であることを証明すれば十分である。
Lemma \ref{sdga-lemma-good-direct-sum} と比較せよ。


\medskip\noindent
$i \geq 1$ に対して
$\mathcal{S} = \mathcal{S}_0 \amalg \ldots \amalg \mathcal{S}_i$
と書く。このとき、次数付き $\mathcal{O}$-代数として
$\mathcal{A}_i = \mathcal{O}\langle \mathcal{S} \rangle$ である。
この代数の微分次数付き $\mathcal{O}$-部分加群による
フィルトレーションを構成する。

\medskip\noindent
$W = \mathbf{Z}_{\geq 0}^{i + 1}$ とおき、辞書式順序を入れる。
すなわち、$W$ の $w = (w_0, \ldots w_i)$ と
$w' = (w'_0, \ldots, w'_i)$ に対して
$$
w > w' \Leftrightarrow
\exists j,\ 0 \leq j \leq i :
w_i = w'_i,\ w_{i - 1} = w'_{i - 1},\ \ldots ,
\ w_{j + 1} = w'_{j + 1},\ w_j > w'_j
$$
などと定める。
$\mathcal{S} \times \ldots \times \mathcal{S}$ の
$U$ 上の切断 $s = s_1 \cdot \ldots \cdot s_r$ が与えられたとする。
各 $a$ に対して一意な $0 \leq j_a \leq i$ が存在して
$s_a \in \mathcal{S}_{j_a}(U)$ となるとき、
{\it $s$ の重みが定義される}という。この場合、重みを
$$
w(s) = (w_0(s), \ldots, w_i(s)) \in W,\quad
w_j(s) = |\{a \mid j_a = j\}|
$$
によって定める。
$\mathcal{S} \times \ldots \times \mathcal{S}$ の任意の切断の重みは
局所的に定義される。一定の重みをもつ切断からなる部分層への
互いに素な和の分解
$$
\mathcal{S} \times \ldots \times \mathcal{S} =
\coprod\nolimits_{w \in W} \left(
\mathcal{S} \times \ldots \times \mathcal{S}\right)_w
$$
が得られることは容易に確かめられる。
もちろん、与えられた $r$ に対して現れるのは
$\sum_{0 \leq j \leq i} w_j = r$ を満たす $w \in W$ だけである。
これに応じて分解
$$
\mathcal{A}_i = \mathcal{O} \oplus
\bigoplus\nolimits_{r \geq 1}
\bigoplus\nolimits_{w \in W}
\mathcal{O}[\left(\mathcal{S} \times \ldots \times \mathcal{S}\right)_w]
$$
を得る。証明の残りは次の自明な観察に基づく。
$r$、$w$ と
$\left(\mathcal{S} \times \ldots \times \mathcal{S}\right)_w$
の局所切断 $s = s_1 \cdot \ldots \cdot s_r$ が与えられると、
$$
\text{d}(s) \text{ は次の層の局所切断である： } \mathcal{O} \oplus
\bigoplus\nolimits_{r' \geq 1}
\bigoplus\nolimits_{w' \in W,\ w' < w}
\mathcal{O}[\left(\mathcal{S} \times \ldots \times \mathcal{S}\right)_{w'}]
$$
となる。その理由は、和をとると $\text{d}(s)$ となる各式
$$
(-1)^{\deg(s_1) + \ldots + \deg(s_{a - 1})}
s_1 \cdot \ldots s_{a - 1} \cdot \delta(s_a) \cdot
s_{a + 1} \cdot \ldots \cdot s_r
$$
において、元 $\delta(s_a)$ は局所的には
$s'_1 \cdot \ldots \cdot s'_{r'}$ という元の
$\mathcal{O}$-線形結合であり、$s'_{a'}$ は
ある $0 \leq j'_{a'} < j_a$ に対して
$\mathcal{S}_{j'_a}$ に属するからである。ここで $j_a$ は
$s_a$ が $\mathcal{S}_{j_a}$ の切断となる添字である。

\medskip\noindent
これは次のことを意味する。$w \in W$ に対して
$$
F_w \mathcal{A}_i = \mathcal{O} \oplus \bigoplus\nolimits_{r \geq 1}
\bigoplus\nolimits_{w' \in W,\ w' \leq w}
\mathcal{O}[\left(\mathcal{S} \times \ldots \times \mathcal{S}\right)_{w'}]
$$
とおく。上の観察により、これは微分次数付き
$\mathcal{O}$-部分加群である。右辺の微分を零とする
微分次数付き $\mathcal{A}$-加群の許容短完全列
$$
0 \to \colim_{w' < w} F_{w'}\mathcal{A}_i \to
F_w\mathcal{A}_i \to
\bigoplus\nolimits_{r \geq 1}
\mathcal{O}[\left(\mathcal{S} \times \ldots \times \mathcal{S}\right)_w]
\to 0
$$
を得る。

\medskip\noindent
最後に、順序集合 $W$ に関する超限帰納法で証明を終える。
微分次数付き複体 $F_0\mathcal{A}_0$ は直和因子
$\mathcal{O}$ であり、これは K-平坦かつ次数付き平坦である。
$w \in W$ とし、すべての $w' < w$ に対して
$F_{w'}\mathcal{A}_i$ について結果が成り立つならば、
Lemmas \ref{sdga-lemma-good-direct-sum},
\ref{sdga-lemma-good-admissible-ses}, and \ref{sdga-lemma-free-graded-module-good}
により $w$ に対する結果を得る。最後に
$\mathcal{A}_i = \colim_{w \in W} F_w\mathcal{A}_i$ であるから、
結論を得る。
\end{proof}

\begin{lemma}
\label{sdga-lemma-good-dga}
$(\mathcal{C}, \mathcal{O})$ を環付きサイトとする。
$(\mathcal{B}, \text{d})$ を微分次数付き $\mathcal{O}$-代数とする。
微分次数付き $\mathcal{O}$-代数の擬同型
$(\mathcal{A}, \text{d}) \to (\mathcal{B}, \text{d})$ で、
$\mathcal{A}$ が次数付き平坦であり、$\mathcal{O}$-加群の
複体として K-平坦であり、しかも環付きトポスの任意の射による
引き戻しの後にも同じことが成り立つものが存在する。
\end{lemma}

\begin{proof}
証明は Lemma \ref{sdga-lemma-resolve} の最初の証明とまったく同じであるが、
ここでは自由次数付き加群の代わりに自由次数付き代数を用いる。

\medskip\noindent
Lemma \ref{sdga-lemma-special-good} と同様に、
$$
\mathcal{A}_0 \to \mathcal{A}_1 \to \mathcal{A}_2 \to \ldots \to \mathcal{B}
$$
を構成することにより
$\mathcal{A} = \colim \mathcal{A}_i$ を構成する。
$\mathcal{S}_0$ を、次数 $n$ の部分が
$\Ker(\text{d}_\mathcal{B}^n)$ である次数付き集合の層
（Remark \ref{sdga-remark-sheaf-graded-sets}）とする。
微分次数付き加群の準同型
$$
\mathcal{A}_0 = \mathcal{O}\langle \mathcal{S}_0 \rangle
\longrightarrow
\mathcal{B}
$$
であって、$\mathcal{S}_0$ の局所切断 $s$ を
$\mathcal{A}^{\deg(s)}$ の対応する局所切断へ送るものを考える
（これは微分の核に属するため、実際にこの写像は
微分次数付き代数の写像である）。
構成により、コホモロジー層上の誘導写像
$H^n(\mathcal{A}_0) \to H^n(\mathcal{B})$ は全射であり、
したがってこのことはすべての $i$ に対して成り立ち続ける。

\medskip\noindent
構成の帰納段階を述べる。
$\mathcal{A}_i \to \mathcal{B}$ が与えられたとき、
次数 $n$ の部分が
$$
\Ker(\text{d}_{\mathcal{A}_i}^{n + 1})
\times_{\mathcal{B}^{n + 1}, \text{d}}
\mathcal{B}^n
$$
である次数付き集合の層を $\mathcal{S}_{i + 1}$ と書く。
これには標準写像
$$
\delta_{i + 1} : \mathcal{S}_{i + 1} \longrightarrow
\mathcal{A}_i
$$
が備わり、その像は構成により
$\text{d}_{\mathcal{A}_i}$ の核に含まれる。したがって
$\mathcal{A}_{i + 1} = \mathcal{O}\langle
\mathcal{S}_0 \amalg \ldots \mathcal{S}_{i + 1}\rangle$
は $\mathcal{A}_i$ 上の微分を拡張する微分をもつ。
本節冒頭の議論を参照せよ。
$\mathcal{A}_{i + 1}$ から $\mathcal{B}$ への写像は、
$\mathcal{A}_i$ 上では与えられた写像に制限され、
$\mathcal{S}_{i + 1}$ の局所切断 $s = (a, b)$ を
$\mathcal{B}$ の $b$ へ送る一意な次数付き代数の写像である。
$\text{d}(b)$ が $\mathcal{B}$ における $a$ の像であることによって
ちょうど、この写像は微分と両立する。

\medskip\noindent
写像 $\mathcal{A} \to \mathcal{B}$ は擬同型である。
実際、$H^n(\mathcal{A}) = \colim H^n(\mathcal{A}_i)$ であり、
各 $i$ に対して写像
$H^n(\mathcal{A}_i) \to H^n(\mathcal{B})$ は全射で、その核は
構成により写像
$H^n(\mathcal{A}_i) \to H^n(\mathcal{A}_{i + 1})$
によって零に写される。最後に、$\mathcal{A}$ の平坦性条件は
Lemma \ref{sdga-lemma-special-good} で示した。
\end{proof}




\section{その他}
\label{sdga-section-misc}

\noindent
$(f, f^\sharp) : (\Sh(\mathcal{C}), \mathcal{O})
\to (\Sh(\mathcal{C}'), \mathcal{O}')$ を環付きトポスの射とする。
$\mathcal{A}$ を微分次数付き $\mathcal{O}$-代数の層とする。
合成\footnote{より正確には、
$F(\mathcal{A}) \otimes_\mathcal{O}^\mathbf{L} F(\mathcal{A})
\to F(\mathcal{A} \otimes_\mathcal{O} \mathcal{A}) \to F(\mathcal{A})$
と書くべきである。ここで $F$ は $\mathcal{O}$-加群の複体への
忘却関手を表す。また、
$\mathcal{A} \otimes_\mathcal{O} \mathcal{A}$ は
Section \ref{sdga-section-tensor-product-dg} のテンソル積を表すので、
$F(\mathcal{A} \otimes_\mathcal{O} \mathcal{A}) =
\text{Tot}(F(\mathcal{A}) \otimes_\mathcal{O} F(\mathcal{A}))$
である。この列の最初の矢印は、二つの $\mathcal{O}$-加群の複体の
導来テンソル積から $\mathcal{O}$-加群の複体の通常のテンソル積への
標準写像である。}
$$
\mathcal{A} \otimes_\mathcal{O}^\mathbf{L} \mathcal{A}
\longrightarrow
\mathcal{A} \otimes_\mathcal{O} \mathcal{A}
\longrightarrow
\mathcal{A}
$$
と相対カップ積（Cohomology on Sites, Remark
\ref{sites-cohomology-remark-cup-product} および
Section \ref{sites-cohomology-section-cup-product} を参照）を用いると、
$D(\mathcal{O}')$ における乗法\footnote{ここおよび以下で
$Rf_* : D(\mathcal{O}) \to D(\mathcal{O}')$ は
Cohomology on Sites, Section
\ref{sites-cohomology-section-unbounded} 以下で調べられる導来関手である。}
$$
\mu :
Rf_*\mathcal{A} \otimes_{\mathcal{O}'}^\mathbf{L} Rf_*\mathcal{A}
\longrightarrow
Rf_*\mathcal{A}
$$
を得る。この乗法は結合的である。すなわち、図
$$
\xymatrix{
Rf_*\mathcal{A} \otimes_{\mathcal{O}'}^\mathbf{L} Rf_*\mathcal{A}
\otimes_{\mathcal{O}'}^\mathbf{L} Rf_*\mathcal{A}
\ar[rr]_-{\mu \otimes 1} \ar[d]_{1 \otimes \mu} & &
Rf_*\mathcal{A} \otimes_{\mathcal{O}'}^\mathbf{L} Rf_*\mathcal{A}
\ar[d]^\mu \\
Rf_*\mathcal{A} \otimes_{\mathcal{O}'}^\mathbf{L} Rf_*\mathcal{A}
\ar[rr]^-\mu & &
Rf_*\mathcal{A}
}
$$
は $D(\mathcal{O}')$ において可換である。これは
Cohomology on Sites, Lemma
\ref{sites-cohomology-lemma-cup-product-associative} から従う。
まったく同様に、右微分次数付き $\mathcal{A}$-加群
$\mathcal{M}$ が与えられると、$D(\mathcal{O}')$ における乗法
$$
\mu_\mathcal{M} :
Rf_*\mathcal{M} \otimes_{\mathcal{O}'}^\mathbf{L} Rf_*\mathcal{A}
\longrightarrow
Rf_*\mathcal{M}
$$
を得る。この乗法は上の $\mu$ と両立する。すなわち、図
$$
\xymatrix{
Rf_*\mathcal{M} \otimes_{\mathcal{O}'}^\mathbf{L} Rf_*\mathcal{A}
\otimes_{\mathcal{O}'}^\mathbf{L} Rf_*\mathcal{A}
\ar[rr]_-{\mu_\mathcal{M} \otimes 1} \ar[d]_{1 \otimes \mu} & &
Rf_*\mathcal{M} \otimes_{\mathcal{O}'}^\mathbf{L} Rf_*\mathcal{A}
\ar[d]^{\mu_\mathcal{M}} \\
Rf_*\mathcal{M} \otimes_{\mathcal{O}'}^\mathbf{L} Rf_*\mathcal{A}
\ar[rr]^-{\mu_\mathcal{M}} & &
Rf_*\mathcal{M}
}
$$
は $D(\mathcal{O}')$ において可換である。これも
Cohomology on Sites, Lemma
\ref{sites-cohomology-lemma-cup-product-associative} から従う。

\medskip\noindent
上の特殊な例として、$f$ を一点トポス $\Sh(pt)$ への射とする場合がある。
この場合、$\mu$ は単にカップ積写像
$$
R\Gamma(\mathcal{C}, \mathcal{A})
\otimes_{\Gamma(\mathcal{C}, \mathcal{O})}^\mathbf{L}
R\Gamma(\mathcal{C}, \mathcal{A})
\longrightarrow
R\Gamma(\mathcal{C}, \mathcal{A}),
\quad \eta \otimes \theta \mapsto \eta \cup \theta
$$
であり、同様に $\mu_\mathcal{M}$ はカップ積写像
$$
R\Gamma(\mathcal{C}, \mathcal{M})
\otimes_{\Gamma(\mathcal{C}, \mathcal{O})}^\mathbf{L}
R\Gamma(\mathcal{C}, \mathcal{A})
\longrightarrow
R\Gamma(\mathcal{C}, \mathcal{M}),
\quad \eta \otimes \theta \mapsto \eta \cup \theta
$$
である。一般には、Cohomology on Sites, Remark
\ref{sites-cohomology-remark-before-Leray} の同一視
$$
R\Gamma(\mathcal{C}, \mathcal{A}) =
R\Gamma(\mathcal{C}', Rf_*\mathcal{A})
\quad\text{and}\quad
R\Gamma(\mathcal{C}, \mathcal{M}) =
R\Gamma(\mathcal{C}', Rf_*\mathcal{M})
$$
を通じて、写像 $\mu_\mathcal{M}$ はコホモロジー上にカップ積を誘導する。
これを見るには Cohomology on Sites, Lemma
\ref{sites-cohomology-lemma-compose-cup-product} を用い、その第二の
トポスの射を上のような $\Sh(\mathcal{C}')$ から一点トポスへの射とすればよい。

\medskip\noindent
$\mathcal{M}_1 \to \mathcal{M}_2$ が右微分次数付き
$\mathcal{A}$-加群の準同型ならば、図
$$
\xymatrix{
Rf_*\mathcal{M}_1 \otimes_{\mathcal{O}'}^\mathbf{L} Rf_*\mathcal{A}
\ar[rr]_-{\mu_{\mathcal{M}_1}} \ar[d] & &
Rf_*\mathcal{M}_1 \ar[d] \\
Rf_*\mathcal{M}_2 \otimes_{\mathcal{O}'}^\mathbf{L} Rf_*\mathcal{A}
\ar[rr]^-{\mu_{\mathcal{M}_2}} & &
Rf_*\mathcal{M}_2
}
$$
は $D(\mathcal{O}')$ において可換である。これは相対カップ積が
関手的であることから従う。右微分次数付き $\mathcal{A}$-加群の短完全列
$$
0 \to \mathcal{M}_1 \xrightarrow{a} \mathcal{M}_2 \to \mathcal{M}_3 \to 0
$$
があるとする。このとき、図
$$
\xymatrix{
Rf_*\mathcal{M}_3 \otimes_{\mathcal{O}'}^\mathbf{L} Rf_*\mathcal{A}
\ar[rr]_-{\mu_{\mathcal{M}_3}} \ar[d]_{Rf_*\delta \otimes \text{id}} & & 
Rf_*\mathcal{M}_3 \ar[d]^{Rf_*\delta} \\
Rf_*\mathcal{M}_1[1] \otimes_{\mathcal{O}'}^\mathbf{L} Rf_*\mathcal{A}
\ar[rr]^-{\mu_{\mathcal{M}_1[1]}} & &
Rf_*\mathcal{M}_1[1]
}
$$
は $D(\mathcal{O}')$ において可換であると主張する。ここで
$\delta : \mathcal{M}_3 \to \mathcal{M}_1[1]$ は、
与えられた短完全列から生じる $D(\mathcal{O})$ の射である
（Derived Categories, Section
\ref{derived-section-canonical-delta-functor} を参照）。
この列が次数付き右 $\mathcal{A}$-加群の列として分裂する場合、
これは明らかである。この場合 $\delta$ は右 $\mathcal{A}$-加群の写像で
表せるので、上の議論を適用できるからである。
一般の場合、$a$ の錐と図
$$
\xymatrix{
\mathcal{M}_1 \ar[r]_a \ar[d] &
\mathcal{M}_2 \ar[r]_i \ar[d] &
C(a) \ar[r]_{-p} \ar[d]^q &
\mathcal{M}_1[1] \ar[d] \\
\mathcal{M}_1 \ar[r] &
\mathcal{M}_2 \ar[r] &
\mathcal{M}_3 \ar[r]^\delta &
\mathcal{M}_1[1]
}
$$
を用いて議論する。右側の正方形は、Derived Categories, Lemma
\ref{derived-lemma-derived-canonical-delta-functor} における
$\delta$ の定義により $D(\mathcal{O})$ で可換である。
さて、錐 $C(a)$ は右微分次数付き $\mathcal{A}$-加群の構造をもち、
$i$、$p$、$q$ は右微分次数付き $\mathcal{A}$-加群の準同型となる。
Definition \ref{sdga-definition-cone} を参照せよ。
したがって上の議論により、射 $q$ と $-p$ に対する対応する図が
可換であることが分かる。$q$ は $D(\mathcal{O})$ における同型なので、
望むとおり $\delta$ に対しても同じことが成り立つ。

\medskip\noindent
上の状況で、右微分次数付き $\mathcal{A}$-加群
$\mathcal{M}$ と
$$
\xi \in H^n(\mathcal{C}, \mathcal{M})
$$
を取る。言い換えると、$\xi$ は $\mathcal{M}$ を
$\mathcal{O}$-加群の複体とみたときのコホモロジーにおける
次数 $n$ のコホモロジー類である。
Lemma \ref{sdga-lemma-cohomology-ext} により、右微分次数付き
$\mathcal{A}$-加群の写像
$$
x : \mathcal{A} \rightarrow \mathcal{M}'[n]
\quad\text{and}\quad
s : \mathcal{M} \to \mathcal{M}'
$$
で、$s$ が擬同型であり、導来圏
$D(\mathcal{A}, \text{d})$ における射
$s[n]^{-1} \circ x$、したがって特に導来圏 $D(\mathcal{O})$ における
同じ射による $1 \in H^0(\mathcal{C}, \mathcal{A})$ の像が
$\xi$ となるものを構成できる。したがって、対応する
$D(\mathcal{O})$ の写像
$$
\xi' = (s[n])^{-1} \circ x : \mathcal{A} \longrightarrow \mathcal{M}[n]
$$
は次の二つの性質によって一意に特徴づけられる。
\begin{enumerate}
\item $\xi'$ は $D(\mathcal{A}, \text{d})$ における射へ持ち上げられる。
\item
$H^0(\mathcal{C}, \mathcal{M}[n]) = H^n(\mathcal{C}, \mathcal{M})$
において $\xi = \xi'(1)$ である。
\end{enumerate}
上で論じた $x$ および $s$ と相対カップ積との両立性を用いると、
環付きトポスの任意の射\footnote{例えば恒等射である。}
$(f, f^\sharp) : (\Sh(\mathcal{C}), \mathcal{O})
\to (\Sh(\mathcal{C}'), \mathcal{O}')$ に対し、$\xi'$ の導来順像
$$
Rf_*\xi' : Rf_*\mathcal{A}  \longrightarrow Rf_*\mathcal{M}[n]
$$
は上で構成した写像 $\mu$ および $\mu_{\mathcal{M}[n]}$ と両立する。
すなわち、図
$$
\xymatrix{
Rf_*\mathcal{A} \otimes_{\mathcal{O}'}^\mathbf{L} Rf_*\mathcal{A}
\ar[rr]_-\mu \ar[d]_{Rf_*\xi' \otimes \text{id}}  & &
Rf_*\mathcal{A} \ar[d]^{Rf_*\xi'} \\
Rf_*\mathcal{M}[n] \otimes_{\mathcal{O}'}^\mathbf{L} Rf_*\mathcal{A}
\ar[rr]^-{\mu_{\mathcal{M}[n]}} & &
Rf_*\mathcal{M}[n]
}
$$
は $D(\mathcal{O}')$ において可換である。
この両立性を一点トポスへの写像に適用すると、特に図
$$
\xymatrix{
R\Gamma(\mathcal{C}, \mathcal{A})
\otimes_{\Gamma(\mathcal{C}, \mathcal{O})}^\mathbf{L}
R\Gamma(\mathcal{C}, \mathcal{A})
\ar[d]_{\xi' \otimes \text{id}} \ar[r] &
R\Gamma(\mathcal{C}, \mathcal{A}) \ar[d]^{\xi'} \\
R\Gamma(\mathcal{C}, \mathcal{M}[n])
\otimes_{\Gamma(\mathcal{C}, \mathcal{O})}^\mathbf{L}
R\Gamma(\mathcal{C}, \mathcal{A})
\ar[r] &
R\Gamma(\mathcal{C}, \mathcal{M}[n])
}
$$
が可換であることが分かる。これと $\xi'(1) = \xi$ を合わせると、
コホモロジー上に誘導される写像
$$
\xi' : R\Gamma(\mathcal{C}, \mathcal{A}) \to
R\Gamma(\mathcal{C}, \mathcal{M}[n]), \quad \eta \mapsto \xi \cup \eta
$$
は、表示どおり $\xi$ による左カップ積で与えられる。



\section{圏上の微分次数付き加群}
\label{sdga-section-modules-cohomology}

\noindent
本節は Cohomology on Sites, Section
\ref{sites-cohomology-section-modules-cohomology} の続きである。

\medskip\noindent
$\mathcal{C}$ を圏とする。$\mathcal{C}$ を密着位相をもつサイトとみなす。
$\mathcal{O}$ を $\mathcal{C}$ 上の環の層とする。
$(\mathcal{A}, \text{d})$ を微分次数付き
$\mathcal{O}$-代数の層とする。言い換えると、$\mathcal{O}$ は圏
$\mathcal{C}$ 上の環の前層であり、$(\mathcal{A}, \text{d})$ は
$\mathcal{C}$ 上の微分次数付き $\mathcal{O}$-代数の前層である。
Categories, Definition \ref{categories-definition-presheaf} を参照せよ。

\begin{definition}
\label{sdga-definition-cartesian}
上の状況で、$D(\mathcal{A}, \text{d})$ の対象 $M$ のうち、
$\mathcal{C}$ のすべての $U \to V$ に対して標準写像
$$
R\Gamma(V, M) \otimes_{\mathcal{A}(V)}^\mathbf{L} \mathcal{A}(U)
\longrightarrow
R\Gamma(U, M)
$$
が $D(\mathcal{A}(U), \text{d})$ における同型となるものからなる
充満部分圏を
{\it $\mathit{QC}(\mathcal{A}, \text{d})$}
と書く。
\end{definition}

\begin{lemma}
\label{sdga-lemma-cartesian-good-sub}
上の状況で、部分圏 $\mathit{QC}(\mathcal{A}, \text{d})$ は
$D(\mathcal{A}, \text{d})$ の厳密充満な飽和三角部分圏であり、
任意の直和によって保たれる。
\end{lemma}

\begin{proof}
$U$ を $\mathcal{C}$ の対象とする。$\mathcal{C}$ 上の位相は密着位相なので、
関手 $\mathcal{F} \mapsto \mathcal{F}(U)$ は完全であり、
直和と可換である。したがって完全関手
$M \mapsto R\Gamma(U, M)$ は、$K$ を任意の微分次数付き
$\mathcal{A}$-加群 $\mathcal{M}$ で表し、$\mathcal{M}(U)$ を
取ることで計算される。ゆえに $R\Gamma(U, -)$ は直和と可換する。
Lemma \ref{sdga-lemma-derived-products} を参照せよ。
同様に、$\mathcal{C}$ の射 $U \to V$ が与えられると、
導来テンソル積関手
$- \otimes_{\mathcal{O}(A)}^\mathbf{L} \mathcal{A}(U) :
D(\mathcal{A}(V)) \to D(\mathcal{A}(U))$
は完全であり、直和と可換する。これらの観察から本補題は直ちに従う。
詳細は省略する。
\end{proof}

\begin{remark}
\label{sdga-remark-cartesian-brown}
上と同様に、$\mathcal{C}$ を密着位相をもつサイトとみなした圏、
$\mathcal{O}$ を $\mathcal{C}$ 上の環の層、
$(\mathcal{A}, \text{d})$ を微分次数付き
$\mathcal{O}$-代数の層とする。このとき Cohomology on Sites,
Proposition \ref{sites-cohomology-proposition-cartesian-brown}
の類似が、ほぼまったく同じ証明によって
$\mathit{QC}(\mathcal{A}, \text{d})$ に対して成り立つ。
\begin{enumerate}
\item 直和を直積へ移す任意の反変コホモロジー的関手
$H : \mathit{QC}(\mathcal{A}, \text{d}) \to \textit{Ab}$
は表現可能である。
\item 直和を直和へ移す三角圏の任意の完全関手
$F : \mathit{QC}(\mathcal{A}, \text{d}) \to \mathcal{D}$
は完全な右随伴をもつ。
\item 包含関手
$\mathit{QC}(\mathcal{A}, \text{d}) \to D(\mathcal{A}, \text{d})$
は完全な右随伴をもつ。
\end{enumerate}
これが必要になった場合には、ここで正確に定式化して証明することにする。
\end{remark}

\noindent
$u : \mathcal{C}' \to \mathcal{C}$ を圏の間の関手とする。
$\mathcal{C}$ と $\mathcal{C}'$ を密着位相をもつサイトとみなすと、
$u$ は連続かつ余連続な関手である。したがってトポスの射
$g : \Sh(\mathcal{C}') \to \Sh(\mathcal{C})$
を得る。Sites, Lemma
\ref{sites-lemma-cocontinuous-morphism-topoi} を参照せよ。
さらに、$\mathcal{C}$ 上の環の層 $\mathcal{O}$ と
$\mathcal{C}'$ 上の環の層 $\mathcal{O}'$、および写像
$g^\sharp : g^{-1}\mathcal{O} \to \mathcal{O}'$
が与えられているとする。対応する環付きトポスの射を単に
$g : (\Sh(\mathcal{C}'), \mathcal{O}') \to
(\Sh(\mathcal{C}), \mathcal{O})$
と書く。Modules on Sites, Section
\ref{sites-modules-section-ringed-topoi} を参照せよ。
最後に、$(\mathcal{A}, \text{d})$ を微分次数付き
$\mathcal{O}$-代数の層、$(\mathcal{A}', \text{d})$ を
微分次数付き $\mathcal{O}'$-代数の層とし、さらに
微分次数付き $\mathcal{O}'$-代数の写像
$\varphi : g^*\mathcal{A} \to \mathcal{A}'$
が与えられているとする
（Section \ref{sdga-section-functoriality-dg} を参照）。

\begin{lemma}
\label{sdga-lemma-cartesian-functorial}
$g : (\Sh(\mathcal{C}'), \mathcal{O}') \to
(\Sh(\mathcal{C}), \mathcal{O})$ と
$\varphi : g^*\mathcal{A} \to \mathcal{A}'$
を上のとおりとする。このとき関手
$Lg^* : D(\mathcal{A}, \text{d}) \to D(\mathcal{A}', \text{d})$
は $\mathit{QC}(\mathcal{A}, \text{d})$ を
$\mathit{QC}(\mathcal{A}', \text{d})$ に写す。
\end{lemma}

\begin{proof}
$U' \in \Ob(\mathcal{C}')$ とし、$\mathcal{C}$ におけるその像を
$U = u(U')$ とする。$pt$ をただ一つの対象とただ一つの射をもつ圏とする。
表示された環付きトポス
$(\Sh(pt), \mathcal{O}'(U'))$ と
$(\Sh(pt), \mathcal{O}(U))$ に、それぞれ微分次数付き代数
$\mathcal{A}'(U)$ と $\mathcal{A}(U)$ を与える。
もちろん、これらの上の微分次数付き加群の導来圏をそれぞれ
$D(\mathcal{A}'(U'), \text{d})$ および
$D(\mathcal{A}(U), \text{d})$ と同一視する。
すると、各頂点に対応する微分次数付き代数を備えた
環付きトポスの可換図
$$
\xymatrix{
(\Sh(pt), \mathcal{O}'(U')) \ar[rr]_{U'} \ar[d] & &
(\Sh(\mathcal{C}'), \mathcal{O}') \ar[d]^g \\
(\Sh(pt), \mathcal{O}(U)) \ar[rr]^U & &
(\Sh(\mathcal{C}), \mathcal{O})
}
$$
を得る。下の水平射に沿う引き戻しは
$D(\mathcal{A}, \text{d})$ の $M$ を、
$D(\mathcal{A}(U), \text{d})$ の対象とみなした
$R\Gamma(U, K)$ に写す。左の垂直矢印による引き戻しは $M$ を
$M \otimes_{\mathcal{A}(U)}^\mathbf{L} \mathcal{A}'(U')$
に写す。この図のいずれの向きに合成しても同じ結果を生じる
（Lemma \ref{sdga-lemma-compose-pullback}）。したがって
$$
R\Gamma(U', Lg^*K) =
R\Gamma(U, K) \otimes_{\mathcal{A}(U)}^\mathbf{L} \mathcal{A}'(U')
$$
を得る。最後に、$f' : U' \to V'$ を $\mathcal{C}'$ の射とし、
$\mathcal{C}$ におけるその像を
$f = u(f') : U = u(U') \to V = u(V')$ と書く。
$K$ が $\mathit{QC}(\mathcal{A}, \text{d})$ に属するならば
\begin{align*}
R\Gamma(V', Lg^*K) \otimes_{\mathcal{A}'(V')}^\mathbf{L} \mathcal{A}'(U')
& =
R\Gamma(V, K) \otimes_{\mathcal{A}(V)}^\mathbf{L} \mathcal{A}'(V')
\otimes_{\mathcal{A}'(V')}^\mathbf{L} \mathcal{A}'(U') \\
& =
R\Gamma(V, K) \otimes_{\mathcal{A}(V)}^\mathbf{L} \mathcal{A}'(U') \\
& =
R\Gamma(V, K) \otimes_{\mathcal{A}(V)}^\mathbf{L} \mathcal{A}(U)
\otimes_{\mathcal{A}(U)}^\mathbf{L} \mathcal{A}'(U') \\
& =
R\Gamma(U, K) \otimes_{\mathcal{A}(U)}^\mathbf{L} \mathcal{A}'(U') \\
& =
R\Gamma(U', Lg^*K)
\end{align*}
となり、望む結果を得る。ここでは上の観察を $U'$ と $V'$ の
両方に用いた。
\end{proof}



\section{圏上の微分次数付き加群（続）}
\label{sdga-section-modules-cohomology-bis}

\noindent
Section \ref{sdga-section-modules-cohomology} で導入した擬連接加群の概念に関する
結果をさらにいくつか展開する。

\begin{lemma}
\label{sdga-lemma-cartesian-cofinal-subcategory}
$\mathcal{C}, \mathcal{O}, \mathcal{A}$ を
Section \ref{sdga-section-modules-cohomology} のとおりとする。
$\mathcal{C}' \subset \mathcal{C}$ を次の性質をもつ充満部分圏とする。
すべての $U \in \Ob(\mathcal{C})$ に対し、矢印 $U \to U'$ の圏
$U/\mathcal{C}'$ は余フィルター圏である。
$\mathcal{O}', \mathcal{A}'$ をそれぞれ
$\mathcal{O}, \mathcal{A}$ の $\mathcal{C}'$ への制限とする。
このとき制限は同値
$\mathit{QC}(\mathcal{A}, \text{d}) \to
\mathit{QC}(\mathcal{A}', \text{d})$
を誘導する。
\end{lemma}

\begin{proof}
この関手の擬逆を構成する。すなわち、
$M'$ を $\mathit{QC}(\mathcal{A}', \text{d})$ の対象とする。
$M'$ を良い微分次数付き加群 $\mathcal{M}'$ で表せる。
Lemma \ref{sdga-lemma-resolve} を参照せよ。
このとき、すべての $U' \in \Ob(\mathcal{C}')$ に対して
微分次数付き $\mathcal{A}'(U')$-加群 $\mathcal{M}'(U)$ は
K-平坦かつ次数付き平坦であり、$\mathcal{C}'$ のすべての射
$U'_1 \to U'_2$ に対して写像
$$
\mathcal{M}'(U'_2) \otimes_{\mathcal{A}'(U'_2)} \mathcal{A}'(U'_1)
\longrightarrow
\mathcal{M}'(U'_1)
$$
は擬同型である（始域が導来テンソル積を表すからである）。
規則
$$
\mathcal{M}(U) =
\colim_{U \to U' \in U/\mathcal{C}'}
\mathcal{M}'(U') \otimes_{\mathcal{A}'(U')} \mathcal{A}(U)
$$
によって定義される微分次数付き
$\mathcal{A}$-加群 $\mathcal{M}$ を考える。
本補題の仮定により、これは複体のフィルター余極限である。
$M'$ は $\mathit{QC}(\mathcal{A}', \text{d})$ に属するので、
この系のすべての遷移写像は擬同型である。
フィルター余極限は完全なので、$U' \in \Ob(\mathcal{C}')$ をもつ
任意の射 $U \to U'$ に対し、$D(\mathcal{A}(U), \text{d})$ における
$\mathcal{M}(U)$ は
$\mathcal{M}'(U') \otimes_{\mathcal{A}'(U')} \mathcal{A}(U)$
と同型である。

\medskip\noindent
$\mathcal{M}$ は $\mathit{QC}(\mathcal{A}, \text{d})$ に属すると主張する。
実際、$\mathcal{C}$ における $U \to V$ が与えられたとき、
$V' \in \Ob(\mathcal{C}')$ をもつ写像 $V \to V'$ を選ぶ。
上の議論により、写像 $\mathcal{M}(V) \to \mathcal{M}(U)$ は写像
$$
\mathcal{M}'(V') \otimes_{\mathcal{A}'(V')} \mathcal{A}(V)
\longrightarrow
\mathcal{M}'(V') \otimes_{\mathcal{A}'(V')} \mathcal{A}(U)
$$
と同一視される。$\mathcal{M'}(V')$ は微分次数付き
$\mathcal{A}'(V')$-加群として K-平坦なので、主張が従う。

\medskip\noindent
自然写像 $\mathcal{M}|_{\mathcal{C}'} \to \mathcal{M}'$ は
$D(\mathcal{A}', d)$ における同型である。これは上の議論から直ちに従う。

\medskip\noindent
逆に、$\mathit{QC}(\mathcal{A}, \text{d})$ の対象 $E$ があるとする。
これを良い微分次数付き加群 $\mathcal{E}$ で表す。
$\mathcal{M}' = \mathcal{E}|_{\mathcal{C}'}$ とおく
（これも良い微分次数付き加群である）と、写像
$$
\mathcal{E} \to \mathcal{M}
$$
が存在することが分かる。これは $\mathcal{C}$ の $U$ 上では写像
$$
\mathcal{E}(U)
\longrightarrow
\colim_{U \to U' \in U/\mathcal{C}'}
\mathcal{E}(U') \otimes_{\mathcal{A}'(U')} \mathcal{A}(U)
$$
によって与えられ、同じ理由で擬同型である。
したがって、望むとおり、制限と上の構成は互いに擬逆な関手である。
\end{proof}

\begin{lemma}
\label{sdga-lemma-cartesian-qis-equivalence}
$\mathcal{C}, \mathcal{O}$ を Section
\ref{sdga-section-modules-cohomology} のとおりとする。
$\varphi : \mathcal{A} \to \mathcal{B}$ を、コホモロジー層上に
同型を誘導する微分次数付き $\mathcal{O}$-代数の準同型とする。
このとき Lemma \ref{sdga-lemma-qis-equivalence} の同値
$D(\mathcal{A}, \text{d}) \to D(\mathcal{B}, \text{d})$
は同値
$\mathit{QC}(\mathcal{A}, \text{d}) \to
\mathit{QC}(\mathcal{B}, \text{d})$
を誘導する。
\end{lemma}

\begin{proof}
次を示せば十分である。$\mathcal{C}$ の射 $U \to V$ と
$D(\mathcal{A}, \text{d})$ の $M$ が与えられたとき、
次の二条件は同値である。
\begin{enumerate}
\item $R\Gamma(V, M) \otimes_{\mathcal{A}(V)}^\mathbf{L} \mathcal{A}(U)
\to \Gamma(U, M)$ は $D(\mathcal{A}(U), \text{d})$ における同型である。
\item $R\Gamma(V, M \otimes_\mathcal{A}^\mathbf{L} \mathcal{B})
\otimes_{\mathcal{B}(V)}^\mathbf{L} \mathcal{B}(U)
\to \Gamma(U, M \otimes_\mathcal{A}^\mathbf{L} \mathcal{B})$
は $D(\mathcal{B}(U), \text{d})$ における同型である。
\end{enumerate}
$\mathcal{C}$ 上の位相は密着位相なので、これは単に
$\mathcal{A}(U) \to \mathcal{B}(U)$ と
$\mathcal{A}(V) \to \mathcal{B}(V)$ が擬同型であるという事実に帰着する。
詳細は省略する。
\end{proof}



\section{微分次数付き代数の逆系}
\label{sdga-section-qc-inverse-system}

\noindent
本節では、Section \ref{sdga-section-modules-cohomology} で論じた状況の
次の特殊な場合を考える。
\begin{enumerate}
\item $\mathcal{C}$ は、$i \leq j$ のとき、かつそのときに限り
一意な射 $i \to j$ をもつ圏 $\mathbf{N}$ である。
\item $\mathcal{O}$ は、与えられた環 $R$ を値にもつ
環の定数（前）層である。
\end{enumerate}
この設定では、微分次数付き $\mathcal{O}$-代数の層
$\mathcal{A}$ は、微分次数付き $R$-代数の逆系 $(A_n)$ と同じものである。
微分次数付き $\mathcal{A}$-加群の層 $\mathcal{M}$ は、
$M_n$ が微分次数付き $A_n$-加群であり、遷移写像
$M_{n + 1} \to M_n$ が $A_{n + 1}$-加群の写像となる逆系
$(M_n)$ と同じものである。本節では記法から微分を省略し、
Definition \ref{sdga-definition-cartesian} で
$\mathit{QC}(\mathcal{A}, \text{d})$ と書いた圏を
$\mathit{QC}(\mathcal{A})$ と書く。

\medskip\noindent
$\mathcal{B} = (B_n)$ を微分次数付き $R$-代数の第二の逆系とする。
pro-対象の射 $\varphi : (A_n) \to (B_n)$ が与えられたとき、
$\mathit{QC}(\mathcal{A})$ から $\mathit{QC}(\mathcal{B})$ への
完全関手を構成する。実際、Categories, Example
\ref{categories-example-pro-morphism-inverse-systems}
によれば、射 $\varphi$ は整数列
$\ldots \geq m(3) \geq m(2) \geq m(1)$ と
微分次数付き $R$-代数の可換図
$$
\xymatrix{
\ldots \ar[r] &
A_{m(3)} \ar[d]^{\varphi_3} \ar[r] &
A_{m(2)} \ar[d]^{\varphi_2} \ar[r] &
A_{m(1)} \ar[d]^{\varphi_1} \\
\ldots \ar[r] &
B_3 \ar[r] &
B_2 \ar[r] &
B_1
}
$$
によって与えられる。そこで、$\mathit{QC}(\mathcal{A})$ の対象を表す
良い微分次数付き $\mathcal{A}$-加群の層
$\mathcal{M} = (M_n)$ が与えられたとき、
$$
N_n = M_{m(n)} \otimes_{A_{m(n)}} B_n
$$
とおける。この逆系は $\mathit{QC}(\mathcal{B})$ の対象を定める。
なぜなら $A_{m(n)}$-加群 $M_{m(n)}$ は K-平坦だからである。
詳細は省略する。得られる関手がその構成で行った選択に依存しないことも
読者に委ねる。

\begin{lemma}
\label{sdga-lemma-pro-isomorphic}
上の状況で、$\mathcal{A} = (A_n)$ と $\mathcal{B} = (B_n)$ を
微分次数付き $R$-代数の逆系とする。
$\varphi : (A_n) \to (B_n)$ が pro-対象の同型ならば、
上で構成した関手
$\mathit{QC}(\mathcal{A}) \to \mathit{QC}(\mathcal{B})$
は同値である。
\end{lemma}

\begin{proof}
$\psi : (B_n) \to (A_n)$ を $\varphi$ の逆となる
pro-対象の射とする。Categories, Example
\ref{categories-example-pro-morphism-inverse-systems}
の議論によれば、$\varphi$ は上のような写像の系で与えられ、
$\psi$ は $n(1) < n(2) < \ldots$ と微分次数付き
$R$-代数の可換図
$$
\xymatrix{
\ldots \ar[r] &
B_{n(3)} \ar[d]^{\psi_3} \ar[r] &
B_{n(2)} \ar[d]^{\psi_2} \ar[r] &
B_{n(1)} \ar[d]^{\psi_1} \\
\ldots \ar[r] &
A_3 \ar[r] &
A_2 \ar[r] &
A_1
}
$$
によって与えられると仮定してよい。
$\varphi \circ \psi = \text{id}$ なので、必要なら関数
$n(\cdot)$ と $m(\cdot)$ の値を増やすことにより、
$B_{n(m(i))} \to A_{m(i)} \to B_i$ が遷移写像であると仮定できる。
したがって、関手の合成
$$
\mathit{QC}(\mathcal{B}) \to
\mathit{QC}(\mathcal{A}) \to
\mathit{QC}(\mathcal{B})
$$
は、$\mathit{QC}(\mathcal{B})$ の対象を表す良い微分次数付き
$\mathcal{B}$-加群の層 $\mathcal{N} = (N_n)$ を、値
$$
N'_i = N_{n(m(i))} \otimes_{B_{n(m(i))}} B_i
$$
をもつ逆系 $\mathcal{N}' = (N'_i)$ に写す。
これは $\mathcal{N}$ と標準的に擬同型である。
まさに $\mathcal{N}$ が $\mathit{QC}(\mathcal{B})$ の対象であり、
すべての $j$ に対して $N_j$ が K-平坦な微分次数付き加群だからである。
逆向きの合成についても同じことが成り立つので、結論を得る。
\end{proof}

\noindent
$\mathcal{C} = \mathbf{N}$ とし、$\mathcal{O}$ を環 $R$ を値にもつ
定数層とし、$\mathcal{A}$ が微分次数付き $R$-代数の逆系
$(A_n)$ で与えられるとする。
二つの左微分次数付き $\mathcal{A}$-加群
$\mathcal{N}$ と $\mathcal{N}'$ が、それぞれ逆系
$(N_n)$ と $(N'_n)$ によって与えられているとする。
したがって、各 $N_n$ と $N'_n$ は左微分次数付き $A_n$-加群である。
整数列
$$
1 = n_0 < n_1 < n_2 < n_3 < \ldots
$$
および写像
$$
N_{n_{2i}} \to N'_{n_{2i - 1}}
\quad\text{in}\quad
D(A_{n_{2i}}^{opp}, \text{d})
$$
と
$$
N'_{n_{2i + 1}} \to N'_{n_{2i}}
\quad\text{in}\quad
D(A_{n_{2i + 1}}^{opp}, \text{d})
$$
が存在し、合成 $N_{n_{2i}} \to N_{n_{2i - 2}}$ と
$N'_{n_{2i + 1}} \to N'_{2i - 1}$ がそれぞれの系の遷移写像で
与えられるとき、$(N_n)$ と $(N'_n)$ は
{\it 導来圏において pro-同型である}と暫定的にいう。

\begin{lemma}
\label{sdga-lemma-pro-isomorphic-systems-tensor}
$(N_n)$ と $(N'_n)$ が上の意味で導来圏において pro-同型ならば、
$D(\mathbf{N}, \mathcal{A})$ の任意の対象 $(M_n)$ に対して
$$
R\lim (M_n \otimes_{A_n}^\mathbf{L} N_n) =
R\lim (M_n \otimes_{A_n}^\mathbf{L} N'_n)
$$
が $D(R)$ において成り立つ。
\end{lemma}

\begin{proof}
仮定により、$D(R)$ の逆系
$(M_n \otimes_{A_n}^\mathbf{L} N_n)$ は、通常の意味で
$D(R)$ の逆系
$(M_n \otimes_{A_n}^\mathbf{L} N'_n)$ と pro-同型である。
したがって結果は More on Algebra, Lemma
\ref{more-algebra-lemma-pro-isomorphism-Rlim} による。
\end{proof}

\begin{lemma}
\label{sdga-lemma-different-tensors}
\begin{reference}
これは \cite[Lemma 3.5.4]{BS} の変種である。
\end{reference}
$R$ を環とし、$f_1, \ldots, f_r \in R$ とする。
$K_n$ を $f_1^n, \ldots, f_r^n$ に関する Koszul 複体とし、
微分次数付き $R$-代数とみなす。
$(M_n)$ を $D(\mathbf{N}, (K_n))$ の対象とする。
このとき任意の $t \geq 1$ に対して
$$
R\lim (M_n \otimes_R^\mathbf{L} K_t) =
R\lim (M_n \otimes_{K_n}^\mathbf{L} K_t)
$$
が $D(R)$ において成り立つ。
\end{lemma}

\begin{proof}
$t \geq 1$ を固定する。$n \geq t$ に対して、微分次数付き
$R$-代数 $K_t$ を左微分次数付き $K_n$-加群とみなしたものを
${}_nK_t$ と書く。次に注意する。
$$
M_n \otimes_R^\mathbf{L} K_t =
M_n \otimes_{K_n}^\mathbf{L} (K_n \otimes_R^\mathbf{L} K_t) =
M_n \otimes_{K_n}^\mathbf{L} (K_n \otimes_R K_t)
$$
したがって Lemma \ref{sdga-lemma-pro-isomorphic-systems-tensor} により、
$({}_nK_t)$ と $(K_n \otimes_R K_t)$ が導来圏において
pro-同型であることを示せば十分である。乗法写像
$$
K_n \otimes_R K_t \longrightarrow {}_nK_t
$$
は左微分次数付き $K_n$-加群の写像である。
ゆえに証明を終えるには、すべての $n \geq 1$ に対して
$N > n$ と $D(K_N^{opp}, \text{d})$ における写像
$$
{}_NK_t \longrightarrow {}_NK_n \otimes_R K_t
$$
で、乗法写像との合成が（どちらの向きでも）遷移写像となるものが
存在することを示せば十分である。これは Divided Power Algebra,
Lemma \ref{dpa-lemma-construct-some-maps} において明示的な構成により
行われている。
\end{proof}

\begin{proposition}
\label{sdga-proposition-derived-complete}
$R$ をネーター環とし、$I \subset R$ をイデアルとする。
次の三つの圏は標準的に同値である。
\begin{enumerate}
\item $\mathcal{A}$ を、$R$-代数の逆系 $A_n = R/I^n$ に対応する
$\mathbf{N}$ 上の $R$-代数の層とする。圏
$\mathit{QC}(\mathcal{A})$。
\item $I$ の生成元 $f_1, \ldots, f_r$ を選ぶ。
$\mathcal{B}$ を、$f_1^n, \ldots, f_r^n$ に関する
Koszul 代数の逆系に対応する $\mathbf{N}$ 上の微分次数付き
$R$-代数の層とする。圏 $\mathit{QC}(\mathcal{B})$。
\item 導来完備対象からなる充満部分圏
$D_{comp}(R, I) \subset D(R)$。
More on Algebra, Definition
\ref{more-algebra-definition-derived-complete} およびその後の文章を参照せよ。
\end{enumerate}
\end{proposition}

\begin{proof}
環付きトポスの明らかな射
$f : (\Sh(\mathbf{N}), \mathcal{A}) \to (\Sh(pt), R)$
を考え、随伴関手 $Lf^*$ と $Rf_*$ を考える。
前者は関手
$$
F : D_{comp}(R, I) \longrightarrow \mathit{QC}(\mathcal{A})
$$
に制限され、K-平坦複体 $K^\bullet$ で表される
$D_{comp}(R, I)$ の対象 $K$ を
$\mathit{QC}(\mathcal{A})$ の対象
$(K^\bullet \otimes_R R/I^n)$ に写す。
後者は関手
$$
G : \mathit{QC}(\mathcal{A}) \longrightarrow D_{comp}(R, I)
$$
に制限され、$\mathit{QC}(\mathcal{A})$ の対象
$(M_n^\bullet)$ を $R\lim M_n^\bullet$ に写す。
出力が導来完備であることは、例えば More on Algebra, Lemma
\ref{more-algebra-lemma-naive-derived-completion} による。
また More on Algebra, Proposition
\ref{more-algebra-proposition-noetherian-naive-completion-is-completion}
から $G \circ F = \text{id}$ が従う。
したがって $F$ と $G$ が互いに擬逆な同値であることを見るには、
$G$ の核が零であることを見れば十分である
（Derived Categories, Lemma
\ref{derived-lemma-fully-faithful-adjoint-kernel-zero} を参照）。
しかし、これを直接示すのは容易ではないようである！

\medskip\noindent
本段落では $\mathit{QC}(\mathcal{A})$ と
$\mathit{QC}(\mathcal{B})$ が同値であることを示す。
$\mathcal{B} = (B_n)$ と書く。ここで $B_n$ は次数
$-r, -r + 1, \ldots, 0$ の余鎖複体とみなした Koszul 複体である。
Divided Power Algebra, Remark
\ref{dpa-remark-pro-system-koszul}
（ただし鎖複体を余鎖複体に変える）により、
$1 < n_1 < n_2 < \ldots$ と、微分次数付き $R$-代数の写像
$B_{n_i} \to E_i \to R/(f_1^{n_i}, \ldots, f_r^{n_i})$
および $E_i \to B_{n_{i - 1}}$ で、図
$$
\xymatrix{
B_{n_1} \ar[d] &
B_{n_2} \ar[d] \ar[l] &
B_{n_3} \ar[d] \ar[l] &
\ldots \ar[l] \\
E_1 \ar[d] &
E_2 \ar[l] \ar[d] &
E_3 \ar[l] \ar[d] &
\ldots \ar[l] \\
B_1 &
B_{n_1} \ar[l] &
B_{n_2} \ar[l] &
\ldots \ar[l]
}
$$
が微分次数付き $R$-代数の可換図となり、
$E_i \to R/(f_1^{n_i}, \ldots, f_r^{n_i})$ が擬同型となるものを
見いだせる。したがって次を得る。
\begin{enumerate}
\item $\mathit{QC}(\mathcal{B})$ と $\mathit{QC}((E_i))$ の間に同値がある。
\item $\mathit{QC}((E_i))$ と
$\mathit{QC}((R/(f_1^{n_i}, \ldots, f_r^{n_i})))$ の間に同値がある。
\item $\mathit{QC}((R/(f_1^{n_i}, \ldots, f_r^{n_i})))$ と
$\mathit{QC}(\mathcal{A})$ の間に同値がある。
\end{enumerate}
実際、(1) には Lemma \ref{sdga-lemma-pro-isomorphic} を上の図に適用できる。
この図は $(E_i)$ と $(B_n)$ が pro-同型であることを示す。
(2) には Lemma \ref{sdga-lemma-cartesian-qis-equivalence} を、
擬同型の逆系
$E_i \to R/(f_1^{n_i}, \ldots, f_r^{n_i})$
に適用できる。(3) には Lemma \ref{sdga-lemma-pro-isomorphic} と、
逆系 $(R/I^n)$ および
$(R/(f_1^{n_i}, \ldots, f_r^{n_i})$ が pro-同型であるという
初等的事実を適用できる。

\medskip\noindent
証明の最初の段落とまったく同様に、随伴関手\footnote{これらの関手が
上の同値を通じて、先に定義した $F$ と $G$ と両立することを示せる。}
$$
F' : D_{comp}(R, I) \longrightarrow \mathit{QC}(\mathcal{B})
\quad\text{and}\quad
G' : \mathit{QC}(\mathcal{B}) \longrightarrow D_{comp}(R, I).
$$
を定義できる。前者は、K-平坦複体 $K^\bullet$ で表される
$D_{comp}(R, I)$ の対象 $K$ を
$\mathit{QC}(\mathcal{B})$ の対象
$(K^\bullet \otimes_R B_n)$ に写す。後者は
$\mathit{QC}(\mathcal{B})$ の対象 $(M_n)$ を $R\lim M_n$ に写す。
上と同様に議論すれば、$G'$ の核が零であることを示せば十分である。
そこで、$\mathcal{M} = (M_n)$ を $\mathcal{B}$ 上の良い
微分次数付き加群の層で、$G'$ の核に属する
$\mathit{QC}(\mathcal{B})$ の対象を表すものとする。このとき
$$
0 = R\lim M_n \Rightarrow
0 = (R\lim M_n) \otimes_R^\mathbf{L} B_t =
R\lim (M_n \otimes_R^\mathbf{L} B_t)
$$
である。Lemma \ref{sdga-lemma-different-tensors} により
$R\lim (M_n \otimes_R^\mathbf{L} B_t) =
R\lim (M_n \otimes_{B_n}^\mathbf{L} B_t)$
である。$(M_n)$ は $\mathit{QC}(\mathcal{B})$ の対象なので、
逆系 $M_n \otimes_{B_n}^\mathbf{L} B_t$ は最終的に
値 $M_t$ の定数系となる。したがって望むとおり $M_t = 0$ である。
\end{proof}

\begin{remark}
\label{sdga-remark-proregular}
$R$ を環とし、$f_1, \ldots, f_r \in R$ をイデアル $I$ を
生成する元の列とする。$K_n$ を $f_1^n, \ldots, f_r^n$ に関する
Koszul 複体とし、微分次数付き $R$-代数とみなす。
すべての $n$ に対して $m > n$ が存在し、
$K_m \to K_n$ が次数 $0$ 以外のコホモロジー上に零写像を誘導するとき、
$f_1, \ldots, f_r$ は{\it 弱 pro-正則列}であるという。
この場合、Proposition \ref{sdga-proposition-derived-complete} の証明の議論は
$R$ がネーター環でない場合にも引き続き成り立つ。
特に、$I$ が弱 pro-正則列によって生成されるならば、
$\mathit{QC}(\{R/I^n\})$ は $R$-線形三角圏として、
導来完備対象の圏 $D_{comp}(R, I)$ と同値であることが分かる。
必要が生じた場合には、ここでこれを正確に述べて証明することにする。
\end{remark}
